\documentclass[a4paper,english,12pt]{amsart}
\synctex=1

\usepackage[margin=3cm]{geometry} % Added by Takehiko. To determine the size of margins

\usepackage[latin1]{inputenc}
\usepackage[bottom]{footmisc}
\usepackage{mathrsfs}
\usepackage{amssymb}
\usepackage{amsmath} 
\usepackage{amsthm}
\usepackage{commath}
\usepackage{mathtools}
\usepackage{amsfonts}
\usepackage{tikz}
\usepackage{bm}
\usepackage{accents}
\usepackage{tikz-cd}
\usetikzlibrary{matrix}
\usepackage{stmaryrd}
\usepackage{hyperref}
\usepackage[shortlabels]{enumitem}
\newcommand{\PP}{\mathbb{P}}
\newcommand{\abss}{\lvert\cdot\rvert}
\newcommand{\QQ}{\mathbb{Q}}
\newcommand{\ddd}{\mathbf d}
\newcommand{\prodin}{\prod _{i=1} ^n}
\newcommand{\prodjn}{\prod _{j=1} ^n}
\newcommand{\gk}{g_K}
\newcommand {\Gm}{\mathbb{G}_m}
\newcommand{\Ga}{\mathbb {G} _a}
\newcommand{\GG}{ G}
\newcommand{\RR}{\mathbb{R}}
\newcommand{\Rgz}{\mathbb{R} _{>0}}
\newcommand{\CC}{\mathbb{C}}
\newcommand{\CCC}{\mathcal C}
\newcommand{\ccc}{\mathbf c}
\newcommand{\oF}{\overline F}
\newcommand{\ZZ}{\mathbb{Z}}
\newcommand{\NN}{\mathbb N}
\newcommand{\DD}{\mathcal{D}}
\newcommand{\HH}{\mathcal H}
\newcommand{\II}{\mathcal I}
\newcommand{\XX}{\mathcal{X}}
\newcommand{\doots}{,\dots ,}
\newcommand{\AAK}{\mathbb{A} _{K}}
\newcommand{\AAF}{\mathbb{A} _F}
\newcommand{\AAFt}{\mathbb{A} _F^\times}
\newcommand{\AAFj}{\mathbb {A}_F^1}
\newcommand{\AAFja}{(\mathbb A_F^1)_\aaa}
\newcommand{\AAFtn}{(\mathbb A_F^\times)^n}
\newcommand{\AAFta}{(\mathbb A_F^\times)_\aaa}
\newcommand{\AAS}{\mathbb {A}_S}
\newcommand{\acZ}{\accentset{\circ}Z}
\newcommand{\AK}{\mathfrak A_K}
\newcommand{\Akk}{\mathfrak A_4}
\newcommand{\RRa}{\mathbb R ^{a_0,\dots ,a_n}}
\newcommand{\OOK}{\mathcal O _K}
\newcommand{\OO}{\mathcal O}
\newcommand{\OFS}{\mathcal O_{F,S}}
\newcommand{\EEE}{\mathcal E}
\newcommand{\FFF}{\mathcal F}
\newcommand{\FF}{\mathbb F}
\newcommand{\Ft}{F ^\times}
\newcommand{\Ftn}{F^{\times n}}
\newcommand{\ff}{\mathfrak f}
\newcommand{\Fv}{F_v}
\newcommand{\Fvt}{F_v^\times}
\newcommand{\Fvn}{F^n_v}
\newcommand{\Fvtn}{(F_v^{\times})^n}
\newcommand{\fvt}{f_v^\#}
\newcommand{\sG}{\prescript{}{\sigma}G}
\newcommand{\sN}{\prescript{}{\sigma}N}
\newcommand{\sA}{\prescript{}{\sigma}A}
\newcommand{\sak}{\prescript{}{\sigma}(\mathfrak A_4)}
\newcommand{\muAF}{\mu_{\AAF}}
\newcommand{\Hvt}{H_v^\#}
\newcommand{\pr}{pr}
\newcommand{\OOF}{\mathcal O _F}
\newcommand{\Ov}{\mathcal O _v}
\newcommand{\OOFt}{\OOF ^\times}
\newcommand{\OOT}{\mathcal O_{T}}
\newcommand{\OOS}{\mathcal O _{S}}
\newcommand{\Ovn}{\mathcal O _v ^{\times n}}
\newcommand{\ovk}{\omega _v ^{(4,6)}}
\newcommand{\vp}{v_{\mathfrak p}}
\newcommand{\ppp}{v}
\newcommand{\vMF}{v\in M_F}
\newcommand{\vMFz}{v\in M_F ^0}
\newcommand{\vMFi}{v\in M_F ^\infty}
\newcommand{\Mfi}{M_F^\infty}
\newcommand{\omI}{\overline{\mathfrak I}}
\newcommand{\vMFC}{v\in M_F^{\CC}}
\newcommand{\vMFR}{v\in M_F^{\RR}}
\newcommand{\qqq}{\mathfrak q}
\newcommand{\XXX}{\mathscr X}
\newcommand{\YYY}{\mathscr Y}
\newcommand{\ZZZ}{\mathscr Z}
\newcommand{\LLL}{\mathscr L}
\newcommand{\AAQ}{\mathbb A _{\mathbb Q}}
\newcommand{\ZZp}{\mathbb Z_p}
\newcommand{\QQp}{\mathbb Q_p}
\newcommand{\sss}{\mathbf{s}}
\newcommand{\vvv}{\mathbf v}
\newcommand{\sps}{\mathbf{s}\in\RR^n_{>0}+i\RR^n}
\newcommand{\xxx}{\mathbf{x}}
\newcommand{\rrr}{\mathbf{r}}
\newcommand{\llll}{\boldsymbol{\ell}}
\newcommand{\txxx}{\widetilde \xxx}
\newcommand{\wx}{\widetilde\xxx}
\newcommand{\wy}{\widetilde\yyy}
\newcommand{\ttt}{\mathbf{t}}
\newcommand{\aaa}{\mathbf{a}}
\newcommand{\all}{\mathbf{\alpha}}
\newcommand{\AAA}{\mathbb A}
\newcommand{\AAj}{\mathbb {A} ^1}
\newcommand{\bbb}{\mathbf{b}}
\newcommand{\yyy}{\mathbf{y}}
\newcommand{\zzz}{\mathbf{z}}
\newcommand{\Zo}{\accentset{\circ}{Z}}
\newcommand{\www}{\mathbf{w}}
\newcommand{\kkk}{\mathbf{k}}
\newcommand{\chiv}{\chi_v}
\newcommand{\chivk}{\chi_v^{(k)}}
\newcommand{\uuu}{\mathbf{u}}
\newcommand{\ujn}{(u_j)_{j=1}^{n-1}}
\newcommand{\ujkn}{(u_j)_{\substack{j=1\\j\neq k}}^{n-1}}
\newcommand{\jed}{\mathbf {1}}
\newcommand{\mmm}{\mathbf{m}}
\newcommand{\iR}{i\mathbb R}
\newcommand{\sJ}{\prescript{}{\sigma}J}
\newcommand{\Hv}{H _v}
\newcommand{\Fvun}{F_v^{\un}}
\newcommand{\Fvtame}{F_v^{\tame}}
\newcommand{\Fvs}{\overline{F_v}}
\newcommand{\Hav}{H _v ^\aaa}
\newcommand{\wHav}{\widehat{H_v^\aaa}}
\newcommand{\wH}{\widehat H}
\newcommand{\wchi}{\widetilde{\chi}}
\newcommand{\PPP}{\mathscr P}
\newcommand{\Rt}{\RR^\times}
\newcommand{\TT}{\mathscr T}
\newcommand{\TTT}{\mathscr T}
\newcommand{\cqp}{a}
\newcommand{\muup}{\widehat{\muu}_v ^{(p)}}
\newcommand{\piv}{\pi_v}
\newcommand{\GLn}{GL_n}
\newcommand{\pivv}{|\pi _v|_v}
\newcommand{\BGA}{BG(\AAA)}
\newcommand{\Gfs}{G(\overline{F})}
\newcommand{\muu}{\boldsymbol{\mu}}
\newcommand{\piomin}{\pi_0(\mathcal{J}BG)_{\min}}
\def\card #1{\mathopen| #1 \mathclose|}
\DeclareMathOperator{\an}{an}
\DeclareMathOperator{\Ab}{Ab}
\DeclareMathOperator{\Aut}{Aut}
\DeclareMathOperator{\Autt}{\underline{Aut}}
\DeclareMathOperator{\cond}{cond}
\DeclareMathOperator{\coun}{count}
\DeclareMathOperator{\Conj}{Conj}
\DeclareMathOperator{\disc}{disc}
\DeclareMathOperator{\ind}{ind}
\DeclareMathOperator{\End}{End}
\DeclareMathOperator{\Et}{\acute Et}
\DeclareMathOperator{\vol}{vol}
\DeclareMathOperator{\ord}{ord}
\DeclareMathOperator{\Res}{Res}
\DeclareMathOperator{\Pic}{Pic}
\DeclareMathOperator{\wPic}{\widehat{\Pic}}
\DeclareMathOperator{\Proj}{Proj}
\DeclareMathOperator{\Peyre}{Peyre}
\DeclareMathOperator{\sett}{set}
\DeclareMathOperator{\Tam}{Tam}
\DeclareMathOperator{\coker}{coker}
\DeclareMathOperator{\Hom}{Hom}
\DeclareMathOperator{\Gal}{Gal}
\DeclareMathOperator{\GL}{GL}
\DeclareMathOperator{\SL}{SL}
\DeclareMathOperator{\cycl}{cycl}
\DeclareMathOperator{\Aff}{Aff}
\DeclareMathOperator{\Sing}{Sing}
\DeclareMathOperator{\val}{val}
\DeclareMathOperator{\Div}{Div}
\DeclareMathOperator{\modd}{mod}
\DeclareMathOperator{\Log}{Log}
\DeclareMathOperator{\Ta}{T}
\DeclareMathOperator{\Imm}{Im}
\DeclareMathOperator{\Ideal}{Ideal}
\DeclareMathOperator{\discrete}{discrete}
\DeclareMathOperator{\Jac}{Jac}
\DeclareMathOperator{\lcm}{lcm}
\DeclareMathOperator{\Id}{Id}
\DeclareMathOperator{\prim}{prim}
\DeclareMathOperator{\Elli}{Ell}
\DeclareMathOperator{\supp}{supp}
\DeclareMathOperator{\stable}{stable}
\DeclareMathOperator{\Spec}{Spec}
\DeclareMathOperator{\Zar}{Zar}
\DeclareMathOperator{\fppf}{fppf}
\DeclareMathOperator{\fin}{fin}
\DeclareMathOperator{\aaaa}{(a)}
\DeclareMathOperator{\bbbb}{(b)}
\DeclareMathOperator{\cont}{cont}
\DeclareMathOperator{\rk}{rk}
\DeclareMathOperator{\smooth}{smooth}
\DeclareMathOperator{\Cl}{Cl}
\DeclareMathOperator{\Ann}{Ann}
\DeclareMathOperator{\charr}{char}
\DeclareMathOperator{\Tr}{tr}
\DeclareMathOperator{\tr}{tr}
\DeclareMathOperator{\Reg}{Reg}
\DeclareMathOperator{\Falt}{Falt}
\DeclareMathOperator{\Frob}{Frob}
\DeclareMathOperator{\un}{un}
\DeclareMathOperator{\age}{age}
\DeclareMathOperator{\inv}{inv}
\DeclareMathOperator{\tame}{{tame}}
\DeclareMathOperator{\pijbg}{\pi_0(\mathcal J_0BG)_{V-\min}}
\DeclareMathOperator\et{\acute et}
\def\no{n\textsuperscript{0}\,}
\def\sPic{\mathop{\mathscr P\mathit{ic}}}
\DeclareMathOperator{\Br}{Br}
\def\thesection{\arabic{section}}

\newtheorem{mydef}[equation]{Definition}
\newtheorem{lem}[equation]{Lemma}
\newtheorem{thm}[equation]{Theorem}

\newtheorem{conj}[equation]{Conjecture}

\newtheorem{prop}[equation]{Proposition}

\newtheorem{rem}[equation]{Remark}

\newtheorem{cor}[equation]{Corollary}

\newtheorem{exam}[equation]{Example}

\newtheorem*{theorem*}{Theorem}
\newtheorem{quest}[equation]{Question}
\numberwithin{equation}{subsection}
\usepackage{blindtext}

  \DeclareFontFamily{U}{wncy}{}
    \DeclareFontShape{U}{wncy}{m}{n}{<->wncyr10}{}
    \DeclareSymbolFont{mcy}{U}{wncy}{m}{n}
    \DeclareMathSymbol{\Sh}{\mathord}{mcy}{"58} 
\keywords{Malle conjecture, $G$-torsor, Poisson formula, Generalized discriminant}
\title{Torsors for finite group schemes of bounded height}
\author[Darda]{Ratko Darda}
\address{Department of Mathematics and Computer Science\\ University of Basel}
\email{ratko.darda@gmail.com}
\author[Yasuda]{Takehiko Yasuda}
\address{Department of Mathematics\\ Graduate School of Sciences\\Osaka University}
\email{yasuda.takehiko.sci@osaka-u.ac.jp}
\begin{document}
\begin{abstract}
Let~$F$ be a global field. Let~$G$ be a non trivial finite \' etale tame $F$-group scheme. We define height functions on the set of $G$-torsors over~$F,$ which generalize the usual heights such as discriminant. As an analogue of the Malle conjecture for group schemes, we formulate a conjecture on the asymptotic behaviour of the number of $G$-torsors over~$F$ of bounded height. This is a special case of our more general ``Stacky Batyrev-Manin" conjecture \cite{dardayasudabm}. The conjectured asymptotic is proven for the case when~$G$ is commutative. When~$F$ is a number field, the leading constant is expressed as a product of certain arithmetic invariants of~$G$ and a volume of a space attached to~$G$. Moreover, an equidistribution property of $G$-torsors in the space is established. %Furthermore, an equidistribution of $G$-torsors is established. % in the product space $\prod_{\vMF}BG(\Fv)$. 
%As an application of our theory, we prove that the inverse Galois problem has an affirmative answer for what we call semicommutative group schemes. % (the finite group schemes which can be obtained from finite commutative group schemes by performing iteratively semidirect products and taking quotients). %in the weak Malle conjecture for the alternating group~$\Akk$ embedded in its degree~$4$ and degree~$12$ permutation representations.
\end{abstract}
\maketitle
\section{Introduction}
\subsection{Malle conjecture} A classical result, the Hermite-Minkowski theorem, states that the number of extensions of a number field~$F$ having bounded discriminant is finite. It is a very natural question to ask what is this number. Moreover, one can fix  an invariant, such as the Galois group of the Galois closure, and count such extensions. In this context, we have the famous Malle conjecture:
\begin{conj}[{Malle \cite{Malle}}] \label{classmalle}
Let~$G$ be a non-trivial group which is embedded as a transitive subgroup of the group of permutations~$\mathfrak S_n$ for some $n\geq 1$. For an extension $K/F$ we write $N(\Delta(K/F))$ for the norm of its discriminant and $\Gal(K/F)=G$ if the Galois group of the Galois closure of~$K/F$ is isomorphic to~$G$ as a permutation group. There exists~$C>0$ such that
\begin{equation*}\#\{K/F|[K:F]=n, \Gal(K/F)=G, N(\Delta(K/F))<B\} \sim_{B\to\infty} C B^a\log(B)^{b-1},\end{equation*}
for some explicit invariants~$a=a(G\leq \mathfrak S_n)$ and~$b=b(G\leq \mathfrak S_n, F)$.
\end{conj}
\subsubsection{}The conjecture can be stated for global fields. We list a few of the cases when the Malle conjecture is known to be true:
\begin{itemize}
\item~$G$ is abelian, of order coprime to the characteristic of~$F$ if the characteristic is positive, embedded in its regular representation~\cite{Wright};
\item~$F$ is not of characteristic~$2$, and~$G=\mathfrak S_i$ for $i\in\{3,4,5\}$ in its natural representation \cite{https://doi.org/10.48550/arxiv.1512.03035}, building upon \cite{dave}, \cite{densityquartic}, \cite{densityquintic};
\item~$F$ is a number field, $G$ is~$\mathfrak S_3$ embedded in its regular representation~ \cite{sexticsthree};
\item~$F$ is a number field and $G=\mathfrak S_i\times A$ for $i\in\{3, 4, 5\}$, where~$\mathfrak S_i$ is in its natural representation and~$A$ is abelian in its regular representation \cite{https://doi.org/10.48550/arxiv.2004.04651}, \cite{wangmalle}.
\end{itemize}
A counterexample to Conjecture \ref{classmalle} has been given in \cite{Kluners} by Kl\" uners for $G=(\ZZ/3\ZZ)\wr (\ZZ/2\ZZ)$ embedded in its degree~$6$ representation. %The conjecture admits counterexamples, as we will see shortly. 
It rests unproven for some usual groups such as the alternating group~$\Akk$ embedded in its standard or regular representation. %A weaker version has been given by Malle in \cite{mallefirst}: it predicts that the number of points is bounded below by~$A_1B^a$ and above by~$A_2B^{a+\epsilon}$ for some $A_1, A_2>0$. No counterexamples are known to this version. 
\subsubsection{}\label{malleortorsors} Let~$G_0$ be a point stabilizer for a permutation group~$G\hookrightarrow~\mathfrak S_n.$ For a {\it connected} $G$-torsor~$\Spec(L),$ one has that~$L^{G_0}$ is a degree~$n$ extension of~$F$ having~$G$ for the Galois group of its Galois closure. The fibers of the map $\Spec(L)\mapsto L^{G_0}$ have all equal finite cardinality (by e.g. \cite[Lemma 5.8]{https://doi.org/10.48550/arxiv.1408.3912}). %the s to the Taking the Galois closure of a finite extension $K/F$, defines a {\it connected}~$G$-torsor~$\widetilde K$ over~$F$. Note that the map $K\mapsto\widetilde K$ is $\# G_0$ to~$1$. It admits a section, which is given by taking the subfield fixed by~$G_0$.
We deduce that the Malle conjecture can be studied as a question of counting connected $G$-torsors~$\Spec(L)$, with respect to the discriminant of ~$L^{G_0}$. %(the only difference in the corresponding numbers being the factor $\# G_0$).
\subsection{Non-constant $G$} In this article, we will study the question of counting $G$-torsors for non-constant finite group schemes~$G$. This question appears natural on its own, but let us provide some more motivation to study it. %SAY AN EXAMPLE OF COUNTING POINTS OF STACKS.
\subsubsection{} Our first motivation comes from recent attempts \cite{darda:tel-03682761}, \cite{Ellenberg} to extend a conjecture of Manin on number of rational points of bounded height on {\it Fano} varieties \cite{frankemanintschinkel} to {\it algebraic stacks}. %This was a topic of some recent research \cite{darda2021rational}, \cite{Ellenberg}.
The articles try to explain {\it \` a la Manin} some counting results of points of stacks (such as~\cite{Hortsch}), as well as similarities between the Manin and the Malle conjectures, observed by the second named author in \cite{Yasuda}. For that purpose, definitions of heights for stacks are provided. The generalization of the Manin conjecture for stacks has not been completed: in \cite{darda:tel-03682761} one is restricted to {\it weighted projective stacks}, while in \cite{Ellenberg}, the authors predict the exponent of~$B$ up to an~$\epsilon$ in an asymptotic formula on the number of points, but no leading constant or the exponent of the~$\log(B)$. In the article \cite{dardayasudabm} we will work on this question.

For a finite group scheme~$G$, we have the {\it classifying stack}~$BG$. It is of dimension zero and its set of rational points~$BG(F)$ is precisely the set of $G$-torsors. It looks like a very natural candidate to try to formulate more precise conjecture for the numbers of rational points of bounded height (a ``Manin conjecture"). If~$G$ is constant, the Malle conjecture predicts the number of rational points of~$BG$ for the height given by the discriminant (or the discriminant of a fixed field as above for some embedding $G\hookrightarrow \mathfrak S_n$). But it does not tell anything for non-constant~$G$ and hence does not explain counting results such as \cite[Corollary 9.2.5.10]{darda:tel-03682761} when~$G=\muu_m$ is the finite group scheme of $m$-th roots of unity. %REMOVE ELLENBERG PREDICTSEven for~$G$ constant, there are counting results for some heights different than discriminants, such as conductor (e.g. \cite[Theorem 3.1]{onprobabilitieslocal}). A wide range of heights arising from representations has been provided by the work of Wood and the second named author \cite{woodyasuda} (the heights of \cite{Ellenberg} for the case~$BG$ are precisely those of \cite{woodyasuda}). It is an interesting question to predict the number of $G$-torsors for such heights.
%
%
%In Conjecture \ref{malletale}, we predict the number of rational points of~$BG$ for~$G$ finite tame and \' etale, and heights which are more general than the ones of \cite{woodyasuda} (and hence, than those of \cite{Ellenberg}, because for stacks~$BG$, they are precisely the ones of \cite{woodyasuda}). %This is achieved in t. We define general heights for this stacks and conjecture a  the number of points, see \ref{}. We do not predict the leading constant, it may be a very hard to  an interesting question, that we do not investigate in this article. Moreover, it is possible to extend our conjecture to  case for which Malle conjecture for, while in the second paper 
 %second named author has observed in \cite{Yasuda} similarities between Malle conjecture and a conjecture of Manin . Namely, the predicted asymptotic in the latter conjecture is also of the form $CB^{a}\log(B)^{b-1}$ for certain $C,a, b>0$. This was recently investigated further by the first named author in \cite{darda2021rational} and by Ellenberg, Satriano and Zureick-Brown in \cite{Ellenberg}. In both articles, the idea is to formulate  Manin conjecture for certain algebraic stacks. %The Malle conjecture should correspond to the {\it classifying stack}~$BG$. 
%Arguably, one of the first candidates to do so are the classifying stacks~$BG$, for~$G$ finite group schemes, the points of which are~$G$-torsors, The constant case corresponds to Malle conjecture. In \cite{darda2021rational}, Manin conjecture is formulated for so-called {\it weighted projective stacks}, which includes the case when~$G=\mu_n$ is the group scheme of~$n$-th roots of unity. In \cite{Ellenberg}, based on some heuristics, the authors interpret~$a$-invariants of both conjecture in the same way. In our next article \cite{dardayasudabm}, we propose a version of Manin conjecture, which unconditionally interprets both~$a$ and~$b$ invariants in the same way of both conjectures.
\subsubsection{} Our second motivation is that torsors of non-constant group schemes may influence the counting of the torsors of constant groups.  Let~$\sG$ be an inner twist (see Paragraph \ref{paragtwist}) of a finite group~$G$. It is a finite group scheme, but usually non-constant. There exists a canonical bijection from the set of~$\sG$-torsors to the set of~$G$-torsors. For a closed subgroup~$R$ of~$\sG,$ we have a map from the set of~$R$-torsors to the set of~$\sG$-torsors which has fibers of bounded cardinality. We obtain a canonical map, with fibers of bounded cardinality, from the set of~$R$-torsors to the set of~$G$-torsors. Thus the torsors of the non-constant finite group~$R$ may influence the number of~$G$-torsors. %Such a phenomenon occurs e.g. for~$G=\Akk$ embedded in its regular representation, when a majority of torsors arise in this way (see Subsection \ref{atorsors}), or for $G=(\ZZ/3\ZZ)\wr (\ZZ/2\ZZ)$ embedded in its degree~$6$ representation % in the Kl\" uners' counterexample to Conjecture~\ref{classmalle} from \cite{Kluners}, 
%when even more torsors than predicted by Conjecture \ref{classmalle} arise in this way (a discussion of this phenomenon will be provided in our next article \cite{dardayasudabm}. %
%
%Let us present counterexample to Conjecture \ref{classmalle}, which is due to Kl\" uners \cite{Kluners}. We set~$F=\QQ$ and~$G$ to be the subgroup of~$S_6$ given by~$<(1,2,3), (4,5,6), (1,4), (2,5), (3,6)>.$ There is a map from the set see below. constructs the following counterexample to Conjecture~\ref{classmalle}. Let us set~$F=\QQ$ and $G=(\ZZ/3\ZZ)\wr (\ZZ/2\ZZ).$ We embed~$G$ in its degree~$6$ representation. One can verify that Conjecture~\ref{classmalle} predicts that~$b=1$. It turns out that there are too many extension which contain $\QQ(\xi_3).$ EXPLAIN IN OUR LANGUAGE.
\subsection{Principal results}
Let us state our prediction for the number of $G$-torsors of bounded height and our principal results. Let~$F$ be a global field and let~$G$ be a non-trivial finite \' etale {\it tame} $F$-group scheme (we say that a finite group scheme defined over a field is tame, if the characteristic of the field is zero or if the order of the group is coprime to the characteristic). %We write~$BG(F)$ for the set of~$G$-torsors defined over~$F$.
\subsubsection{}In Definition \ref{definitionofheight}, we define height functions $H:BG(F)\to\RR_{>0}$ which satisfy that the number of $G$-torsors of bounded height is finite (Proposition \ref{estimatenumber}). The heights capture some usual notions of ``size" of $G$-torsors (e.g. the notion of a discriminant). For heights, we have notions of its~$a$ and~$b$ invariants. 

We fix a height $H:BG(F)\to\RR_{>0}$. As we have seen above, for every inner twist~$\sG$ and every closed subgroup~$R$ of~$\sG$, we have a canonical map $BR(F)\to BG(F)$. We call such a morphism {\it breaking thin}\footnote{The terminology is borrowed from \cite{geometricconsitencymanin}.}, if in the lexicographic order, the pair consisting of~$a$ and~$b$ invariants of the pullback height (that is the composite function $BR(F)\to BG(F)\xrightarrow{H}~\RR_{>0}$) is bigger than the pair consisting of~$a$ and~$b$ invariants of~$H$. Let us say that an element of~$BG(F)$ is secure, if it is not in the image of a breaking thin map $BR(F)\to BG(F)$. As an analogue of the Malle conjecture, we propose the following conjecture.
\begin{conj}\label{introversionconj}
\begin{enumerate}
\item Suppose that~$F$ is a number field. There exists~$C>0$ such that $$\#\{x\in BG(F)|\text{$x$ is secure and }H(x)\leq B\}\sim_{B\to \infty} CB^{a(H)}\log (B)^{b(H)-1}.$$
\item Suppose that~$F$ is a function field. One has that $$\#\{x\in BG(F)|\text{$x$ is secure and }H(x)\leq B\}\asymp_{B\to \infty} B^{a(H)}\log (B)^{b(H)-1}. $$
\end{enumerate}
\end{conj}
A modification of Conjecture~\ref{classmalle} has been proposed by Turk\"eli~\cite{turkeli}. The author changes the exponent of $\log(B)$ by taking into account twists of normal subgroups~$N$ of~$G$ with~$G/N$ abelian. As part of our motivation is to extend the Manin conjecture for stacks, removal of torsors seems more natural than changing the logarithmic exponent: this what is usually done for the Manin conjecture e.g. in~\cite{peyrethin} and~\cite{geometricconsitencymanin}. In our next article \cite{dardayasudabm}, we will formulate a Manin type conjecture for certain stacks, the definition of a secure element carries naturally to this context.
\subsubsection{} We suppose that~$G$ is commutative. Then one has that every element of~$BG(F)$ is secure. The following theorem is one of our main results.
\begin{thm}\label{comcasmalle}
Conjecture \ref{introversionconj} is valid for commutative~$G$.
\end{thm}
%Note that for an inner twist $G^{\sigma}$ of~$G$, we have that $G^{\sigma}=G$. Moreover, it is clear that
When~$F$ is a number field and~$G$ is constant, this generalizes the above mentioned result of \cite{Wright} for heights given by discriminants and the result of Wood \cite[Theorem 3.1.]{onprobabilitieslocal} for so-called {\it fair} heights. It also generalizes the result of the first named author \cite[Theorem 8.2.2.12]{darda:tel-03682761} and \cite[Corollary 9.2.5.10]{darda:tel-03682761} for the case~$F$ is a number field and $G=~\muu_m,$ and {\it quasi-toric} and {\it quasi-discriminant} heights, respectively. We will discuss how it is compared with a result of Alberts and O'Dorney \cite{odorney} in Paragraph~\ref{detailscommutative}. 

We will define in Subsection \ref{secequid} a Radon measure~$\omega_H$ on the product space $\prod_{\vMF} BG(F_v)$ (the product goes over the set of places~$M_F$ of~$F$ and for~$v\in M_F$, the set $BG(F_v)$ is the set of $G$-torsors over~$F_v$). In Theorem \ref{satfunf} and Theorem \ref{thmofequidistribution}, we will establish an equidistribution property of $G$-torsors. Our motivation comes from the Manin conjecture, for which the analogous study of equidistribution was proposed by Peyre in \cite{Peyre}. The results are stronger than Theorem \ref{comcasmalle} and give us the following formula for the number of $G$-torsors satisfying certain local conditions.
\begin{thm}\label{equintro} Let $i:BG(F)\to\prod_{\vMF}BG(F_v)$ be the diagonal map. Let~$U\subset \prod_{\vMF}BG(F_v)$ be an open.
\begin{enumerate}  
\item Suppose that~$F$ is a number field and that the boundary of~$U$ is $\omega_H$-negligible. One has that \begin{equation*}\frac{\#\{x\in BG(F)|i(x)\in U, H(x)\leq B\}}{\# G(F)}\sim_{B\to\infty}\frac{\#\Sh^2(G)\cdot\omega_H(U)}{\# G^*(F)\cdot (b(c)-1)!}B^{a(H)}\log(B)^{b(H)-1},
\end{equation*}
where~$G^*$ denotes the Cartier dual of~$G$ and~$\Sh^2(G)$ is the second Tate-Shafarevich group of~$G$. 
\item Suppose that~$F$ is a function field. If $\omega_H(U)>0$, one has that $$\#\{x\in BG(F)|i(x)\in U, H(x)\leq B\}\asymp_{B\to\infty} B^{a(H)}\log(B)^{b(H)-1}.$$
\end{enumerate}
\end{thm}
The support of the measure~$\omega_H$ is given by the closure $\overline{i(BG(F))}\subset\prod_{\vMF}BG(F_v)$ which may be only a proper subset (when $G=\muu_m$, the conditions for the full support are given by the Grunwald-Wang theorem). In the situation of Part (1), we obtain that:
$$\lim_{B\to\infty}\frac{\#\{x\in BG(F)|i(x)\in U, H(x)\leq B\}}{\#\{x\in BG(F)|H(x)\leq B\}}=\frac{\omega_H(U)}{\omega_H(\prod_{\vMF}BG(F_v))}.$$
Hence, in particular, when~$G$ is constant and~$F$ is a number field, we have obtained an analogous formula to \cite[Theorem 2.1]{onprobabilitieslocal} without fairness restriction.  We will apply this result in our future paper \cite{dardayasudasemicomm} to solve the {\it inverse Galois problem} for so-called {\it semicommutative} finite group schemes.
%\subsubsection{}It is possible to choose local conditions so that any global torsor satisfying them is connected. Hence, we obtain (in Paragraph \ref{conntorsor}) that a positive proportion of torsors are connected (i.e. the corresponding algebras are fields):
%\begin{thm}\label{connsssintro} Suppose that~$G$ is commutative.
%One has that $$\#\{x\in BG(F)|x\text{ is connected}, H(x)\leq B\}\asymp_{B\to\infty} B^{a(H)}\log(B)^{b(H)-1}.$$
%\end{thm}
%When~$F$ is a number field, this generalizes the result of \cite[Proposition 9.2.7.1]{darda:tel-03682761} for the case $G=\muu_m$. In this situation, with the additional assumption that $i=\sqrt{-1}\in F$ or that $4\nmid m$, it is possible to derive a precise asymptotic formula for the number of connected $\mu_m$-torsors \cite[Theorem 9.2.7.3]{darda:tel-03682761}.
%\subsubsection{} As an application of our theory, we solve the inverse Galois problem for {\it semicommutative} finite group schemes, that is, we prove existence of connected $G$-torsors for the finite group schemes which can be obtained from finite commutative group schemes by iterated performing semidirect products with commutative kernels and taking quotients by normal subgroups. The constant semicommutative groups are precisely those which can be realized as Galois groups by successive solution to {\it split embedding problems with abelian kernels} and taking intermediate Galois extensions \cite[Chapter IV, Section 2.2]{inversegalois}. The methods of realization, however, do not work for the non-constant case.
%
%A semicommutative finite \' etale group scheme~$G$ can be written as $G=\langle A, K\rangle$, where~$\iota:A\hookrightarrow G$ is normal and commutative and~$K\lneq G$ is semicommutative. We prove:
%\begin{thm}\label{semicintro}
%Suppose that~$G$ is a nontrivial semicommutative finite \' etale and tame $F$-group scheme. Write~$G=\langle A, K\rangle$ as above. There exists $C>0$ such that $$\#\{x\in BG(F)| x\text{ is connected, } H(x)\leq B\}\geq  CB^{a(\iota^*H)},$$where $\iota^*H$ is the pullback height. In particular, there exist infinitely many connected $G$-torsors.
%\end{thm}
%The constant alternating group $G=\mathfrak A_4$ is semicommutative (non-constant examples with $G(\oF)=\mathfrak A_4$ do exist, as discussed in Example~\ref{altasexam}). Our result implies that if the characteristic of~$F$ is not~$2$ or~$3$, the number of $\Akk$-fields of bounded discriminant is growing at least as $C B^{\frac{1}{2}}.$  For the case~$F$ is a number field this was established in \cite[Corollary 1.8]{statfirstg}. The result gives the lower bound predicted by a weaker version of the Malle conjecture \cite{mallefirst} which predicts that the number of extensions of bounded discriminant is at least $CB^{a(G\hookrightarrow S)}$. For constant groups, let us mention other cases not listed before when the lower bound is known to be true:
%\begin{itemize}
%\item $F$ is a number field and~$G$ is nilpotent embedded in its regular representation~\cite{nilpotentmalle}
%\item $F=\QQ$ and~$G=D_{2p}$ is the dihedral group of order~$2p$ where~$p$ is a prime, embedded in its degree~$p$ and degree~$2p$ permutation representations \cite{dihedkluners}.
%\end{itemize}
%For any finite solvable (constant) group~$G$, based on Shafarevich's result on the inverse Galois problem for solvable groups, Alberts establishes in~\cite{statfirstg} a lower bound on number of $G$-extensions. 
\subsection{Overview}We present basic concepts and results of each of sections. We fix a non-trivial finite \' etale tame $F$-group scheme~$G$. 
\subsubsection{} In Section \ref{sectionheights}, a theory of heights is proposed. For that purpose a finite \' etale tame scheme~$G_*$ is defined. When~$G$ is constant its $F$-points are ``$F$-conjugacy classes" from \cite{Malle}. For every place~$v$ such that~$G$ is unramified and tame at~$v$, we have a {\it residue} map $\psi_v:BG(F_v)\to G_*(\Fv)\subset G_*(\oF)$. The scheme $G_*$ and the residue map are very convenient tools to define heights. We define a {\it counting function} to be $\Gamma_F$-invariant function $c: G_*(\oF)\to\RR_{\geq 0}$ with $c(x)=0$ if and only if $x=1$. One easily defines its~$a$ and~$b$ invariants (the earlier mentioned~$a$ and~$b$ invariants of heights are in fact invariants of counting functions). Heights are now defined to be the functions $H:BG(F)\to\RR_{>0}$ for which $v$-adic components~$H_v:BG(F_v)\to\RR_{>0}$ for almost all~$v$ are defined by $H_v(x)=q_v^{c(\psi_v(x))}$. At the end of the section, we state our conjecture. 
%\subsubsection{} In Section we derive Theorem~\ref{semicintro} from Theorem~\ref{equintro}.
\subsubsection{} \label{detailscommutative} Section \ref{sectioncommutativecase} is devoted to the proof of Theorem \ref{equintro}. Let us sketch it. It uses harmonic analysis. The idea originates from a method used to count rational points on equivariant compactifications of algebraic groups, developed initially by Batyrev and Tschinkel for the case of anisotropic tori \cite{aniso}. In the context of counting field extensions, harmonic analysis was first used by Frei, Loughran and Newton in \cite{hassenormprin} (they also used it in \cite{prescnorms}). The method has been used in \cite{odorney} to count elements of $BG(F)$ when~$F$ is a number field and for so-called ``Frobenian invariants". It seems to us that there may be a gap in their analysis: in the proof of Theorem 1.1, the claimed meromorphic extension for non-trivial characters to the left of~$1/a_{\inv}(\mathcal L)$ in their notation does not appear to be automatic.  Without the meromorphic extension for non-trivial characters, one would only obtain an upper bound, such as~$o(B^{a(c)+\epsilon})$. We were told by Alberts that the gap should not be hard to fill by using a result of \cite{statfirstg}, which is based on a result of \cite{prescnorms}. For the sake of completeness, we give a proof using different arguments.

By using Tauberian results, the task becomes to analyse the Dirichlet series $Z(s):=\frac{1}{\# G(F)}\sum_{x\in BG(F)} H(x)^{-s}$ for $s\in\CC$. It follows from polynomial bounds on the number of extensions of fixed degree and of bounded discriminant (c.f. \cite{SchmidtNF} and \cite[Lemma 2.4]{Gmethodant}) that the series converges for $\Re(s)\gg 0$ and defines a holomorphic function in a complex halfplane.  One has a map $$i:BG(F)\to BG(\AAA):=\sideset{}{'}\prod_vBG(F_v),$$ where the restricted product is taken with respect to the unramified cohomology. It follows from the Poitou-Tate exact sequence that the map~$i$ is of finite kernel and of discrete, closed and cocompact image.  By applying the Poisson formula for the inclusion $i(BG(F))\subset BG(\AAA),$ the function~$Z$ is written as a finite sum of Fourier transforms of the extended height $H:BG(\AAA)\to\RR_{>0}$ at automorphic characters (i.e. the ones vanishing at~$i(BG(F))$).   %The global Tate duality identifies the group of the automorphic characters with the group~$BG^*(F)$, where~$G^*$ is the Cartier dual of~$G$. 
We ``approach" these Fourier transforms by $L$-functions of certain Galois representations. Namely, the representations are ``twists" by characters, of  restricted Galois representations induced from Galois submodule of~$G_*(\oF)$ where $c:G_*(\oF)\to\RR_{\geq 0}$ attains its minimal strictly positive value. The analytic properties of~$Z$ are now deduced from analytic properties of $L$-functions. 
%
%At the end of the section we prove Theorem \ref{equintro}.
\subsection{Acknowledgements} This work was supported by JSPS KAKENHI Grant Number JP18H01112. This work has been done during a post-doctoral stay of the first named author at Osaka University. During the stay, he  was supported by JSPS Postdoctoral Fellowship for Research in Japan. The authors would like to thank Antoine Chambert-Loir, Daniel Loughran and Julian Demeio for useful comments. 
\subsection{Conventions} We will use notation~$F$ for a global field. We denote by~$M_F$ (respectively,~$M_F^0$ and~$M_F^\infty$) the set of its places (respectively, of finite places and of infinite places). %We denote by~$\OO_F$ the ring of integers of~$F$.

We fix algebraic closures of~$F$ and of~$F_v$ for $\vMF$ and embeddings of the algebraic closure of~$F$ in each of the algebraic closures of~$F_v$. We denote by~$\overline F$ and for $\vMF$ by $\overline{F_v}$ the separable closure of~$F$ in and~$F_v$ in the chosen algebraic closures.

%We write $\AAA$ for the ring of the adeles of $F$.
If~$v$ is a finite place of~$F$, we denote by~$\Ov$ the ring of integer of~$F_v$, by~$\kappa_v$ the residue field at~$v$ and by~$q_v$ its cardinality.  In the function field case, we write~$q$ for the cardinality of the field of constants of~$F$. The notation~$\Gamma_F$ and~$\Gamma_v$ will be used to denote the absolute Galois group of~$F$ and~$F_v$, respectively.

Let $f, g:\RR_{\geq 0}\to\RR_{\geq 0}$ be two functions, such that for $B\gg 0$ one has that $g(B)\neq 0$. We write $f\sim_{B\to\infty}g$ if $\lim_{B\to\infty}\frac{f(B)}{g(B)}=1$. We write $f\asymp_{B\to\infty}g$ if there are constants $C_1, C_2>0$ such that for every~$B$ big enough one has that $ C_1 g(B)\leq f(B)\leq C_2 g(B)$.
\section{Heights} \label{sectionheights} In this section we will define heights that we will be using to count elements of~$BG(F)$.
\subsection{Cohomologies}\label{equivalenceofcat} We start by defining cohomological groups that we work with. 
\subsubsection{}Let~$M$ be a profinite group and let~$J$ be a finite $M$-group (that is, a finite group endowed with a continuous $M$-action). We denote by $H^1(M, J)$ the pointed set given by the corresponding non-abelian degree~$1$ cohomology \cite[Chapter I, Section 5.2]{Cohomologiegalois}. It is the set of finite sets endowed with compatible~$M$ and~$J$ actions modulo $M$-equivariant bijections such that the map $$J\times A\xrightarrow{(g,x)\mapsto (g\cdot x,x)} A\times A$$ is a bijection. The distinguished element of~$H^1(M,J)$ is given by the class of~$J$, with the action being given by the left multiplication. % which are also $M$-groups and the two actions for which modulo isomorphisms. 
The pointed set $H^1(M, J)$ also allows the following description: it is the pointed set \begin{align*}Z^1(M,J):=\{f:M\to J\text{ continuous}|\forall\gamma,\delta\in M:f(\gamma\delta)=(f(\gamma))(\gamma\cdot f(\delta))\}
\end{align*} modulo the equivalence relation $$f\sim f'\iff \exists g\in J: \forall \gamma\in M: f'(\gamma)=g^{-1}f(\gamma)(\gamma\cdot g).$$
In the case~$J$ is commutative, one has that $Z^1(M,J)$ and $H^1(M, J)$ have structures of abelian groups.
\subsubsection{}Let~$S$ be a locally Noetherian connected scheme and let~$s$ be a geometric point of~$S$. Let $\et(S)$ be the category having finite \' etale $S$-schemes for its objects and \' etale morphisms of $S$-schemes for its morphisms. We let $\pi_1(S, s)$ be the fundamental group of~$S$ at~$s$ and let $\pi_1(S,s)-set$ be the category of {\it finite} sets endowed with a continuous $\pi_1(S,s)$-action. By \cite[\' Expos\' e V, Section 7]{SGAj}, one has an equivalence $$F_S:\et (S)\to \pi_1(S,s)-set.$$ %of the full subcategory of the category of $S$-schemes given by the $S$-schemes which are \' etale over~$S$ and the category of the finite sets endowed with a continuous action of $\pi_1(S, s)$. 
Let~$S'$ be another locally Noetherian connected scheme endowed with a morphism $f:S'\to S$, and~$s'$ its geometric point such that $f(s')=~s$. We have a continuous homomorphism $\pi_1(S', s')\to \pi_1(S,s)$. Let us denote by~$F_S^{-1}$ a quasi-inverse of~$F_S$ and by $f^*:\et(S)\to\et(S')$ the functor given by the base change to~$S'$. The functor $$F_S\circ f^*\circ F_S^{-1}:\pi_1(S,s)-set\to\pi_1(S', s')-set$$ is isomorphic to the functor given by the the homomorphism $\pi_1(S', s')\to \pi_1(S,s)$. 

Let~$G$-be a finite \' etale group scheme over~$S$. We note a canonical identification $$G(S)=F_S(G)^{\pi^1(S,s)}\hspace{2cm} (=:H^0(S, G)).$$ A (left) $G$-torsor over~$S$ is a scheme~$P$ endowed with an {\it fppf} morphism to~$S$ and  with a (left) $G$-action such that the morphism $G\times _SP\to P\times _SP$ given by $(g, p)\mapsto (g\cdot p, p)$ is an isomorphism. A morphism of $G$-torsors $P\to P'$ is a $G$-equivariant morphism $P\to P'$. The trivial~$G$-torsor is $P=G$ endowed with the action given by the left multiplication.
\begin{mydef}
We denote by $BG(S)$ the pointed set of isomorphism classes of $G$-torsors. If $S=\Spec(A)$ for a ring~$A$, we use also notation $BG(A)$.
\end{mydef} 
 As~$G$ is assumed \' etale, by \cite[Chapter III, Proposition 4.2]{Milne}, one has for a $G$-torsor~$P$ that $P\in\et (S)$. Obviously, any \' etale morphism is automatically {\it fppf}, thus~$P\in\et(S)$ is a $G$-torsor if and only if $G\times_S P\to P\times_S P$ is an isomorphism, and using the equivalence~$F_S$, if and only if $F_S(G)\times F_S(P)\to F_S(P)\times F_S(P)$ is an isomorphism in $\pi_1(S,s)-set$. We have hence established a canonical identification of~$BG(S)$ with the set $H^1(\pi_1(S,s), F_S(G))$. % in the sense of . In particular, we have a cocycle definition of $BG(S)$: it is the set \begin{align*}Z^1(\pi_1(S,s), F_S(G))&:=\{w:\pi_1(S,s)\to F_S(G)\text{ continuous and }\\&\quad\quad\quad\forall\gamma,\delta\in\pi_1(S,s):\theta(\gamma\delta)=(w(\gamma))(\gamma\cdot w(\delta))\}
%\end{align*} modulo the equivalence relation $$w\sim w'\iff \exists g\in F_S(G): \forall \gamma\in\pi_1(S,s): w'(\gamma)=g^{-1}w(\gamma)(\gamma\cdot g).$$
%{\color{red} VERIFY AGAIN, maybe use MILNE Chapter V}The functor~$F_S$ induces an equivalence of the abelian subcategories $\Ab(\et(S))$ and $\Ab(\pi_1(S,s)-set)$ given by the abelian group objects and the additive morphisms. If the group scheme~$G$ is commutative (i.e.~$G$ and hence~$F_S(G)$ are abelian group objects in the corresponding categories), then the functor $\Hom_{\Ab(\et(S))}(-,G)$ and the functor $\Hom_{\Ab(\pi_1(S,s)-set)}(-, F_S(G))$ are additive. We obtain the derived cohomological functors. By using identifications from \cite[Chapter II, Section 2.3]{Cohomologiegalois}, in the degree~$0$ and~$1$, this functors define structures of abelian groups on $H^0(S, G)$ and $H^0(\pi_1(S,s), G)$ and $H^1(S, G) $ and $H^1(\pi_1(S,s), G) $, which hence, coincide for the above identifications.
\subsubsection{} We apply results of the previous paragraph. Let~$F$ be a global field. Let $S=\Spec(A)$, where $$A\in\{F\}\cup\cup_{\vMFz}\{\Ov\}\cup\cup_{\vMF}\{F_v\}.$$ %We may write $H^{i}(A, -)$ for $H^i(\Spec(A),-)$. 
By~$s$, we will denote each time a geometric generic point of a scheme~$S$. Note that for every \' etale morphism $S'\to S,$ every point in the geometric fiber over~$s$ is defined over the separable closure of the function field of~$S$, so that $F_{S}(S')$ is given by the points in the separable closure of the function field.

By \cite[\' Expos\' e V, Proposition 8.1]{SGAj}, one has an identification $\pi^1(\Spec(F), s)=\Gamma_F$ and for finite places~$v$, an identification $\pi^1(\Spec(F_v), s)=~\Gamma_v$. For the last identification, by \cite[\' Expos\' e V, Proposition 8.2]{SGAj}, the group $\pi^1(\Spec(\Ov), s)$ is the quotient group of $\pi^1(\Spec(\Fv), s)$ corresponding to the quotient group~$\Gamma_v^{\un}$ of~$\Gamma_v$. 

%Let us suppose that $$S\in\{\Spec(F),\Spec(F_v),\Spec(\Ov)\}$$ and 
Let~$G$ be a finite \' etale $S$-group scheme. %We deduce identifications of pointed sets (of abelian groups if~$G$ is commutative) $$H^{i}(S, G)=H^{i}(\pi^1(S,s), F_S(G))$$ for $i\in\{0,1\}.$ 
When $S=\Spec(\Ov)$, we obtain that the map of pointed sets %(or the abelian group homomorphism, if~$G$ is commutative ) 
$$BG(\Ov)\to BG(F_v)$$ identifies with the unramified cohomology $$H^1(\Gamma_v^{\un}, G(\overline {F_v}))\hookrightarrow H^1(\Gamma_{F_v}, G(\overline {F_v})).$$ (The last map is an inclusion by \cite[Chapter I, \S 5.8 (a) and Chapter I \S 2.6 (b)]{Cohomologiegalois}.) When $S=\Spec(F)$, for every place~$v$, the inclusion $\overline F\hookrightarrow\overline F_v$, yields a map $$BG(F)\to BG(F_v)$$ which corresponds to the map on cohomologies induced by the inclusion $\Gamma_v\subset\Gamma_F$. 
\subsection{Definition of $G_*$}For an \' etale $F$-group scheme~$G$, we are going to define a pointed \' etale scheme~$G_*$. If~$G$ is commutative, then~$G_*$ has a structure of a commutative group scheme. When~$G$ is a constant group scheme, the $\overline F$-points of~$G_*$ correspond to $F$-conjugacy classes from \cite{Malle}. 
\subsubsection{}
%\subsection{Map $H(F_v, G)\to G_*(\oF)$}
%\begin{lem}[NOT THE ABELIAN ONE.]
%Let $G$ be an \' etale group scheme over~$\Spec(F)$. Let $M_F^{\infty}\subset\Sigma\subset M_F$ be finite, such that~$G$ admits an \' etale model~$\mathcal G$ over $\OO_{F, \Sigma}$. Let $x\in H^1(F, G)$. For almost all places~$v\in M_F-\Sigma$, one has that the image of~$x$ in $H^1(\Fv, G)$ lies in $H^1(\Ov, \mathcal G_{\Ov}).$
%\end{lem}
%\subsection{Adelic situation}
%\subsection{Cohomologies}
%Let $G$ be a commutative finite \' etale group scheme over $F$. There exists a finite set $\Sigma$ of finite places of $F$ such that outside of $\Sigma$, the group scheme $G$ admits a commutative finite \' etale model $\mathcal G$ over the open subscheme $S-\Sigma\subset S$ (where $\Sigma$ is identified with the set of corresponding prime ideals in $\OO_{F}$).
%
%If $R$ is a ring, we will identify $BG(R)$ with the fppf-cohomology group $$BG(R)=H^1_{\fppf}(R,G).$$Furthermore, the group scheme $G$ is \' etale, hence smooth, thus by \cite[Theorem B.17]{Cesnavicius} one has an equality \begin{equation}
%\label{equalcohom}
%BG(R)=H^1_{\fppf}(R,G)=H^1_{\et}(R,G).
%\end{equation}
%When $R$ is a field, we furthermore (REFERENCE?) can identify the \' etale cohomology with the Galois cohomology, and we have that
%\begin{equation}
%\label{equalcohomf}
%BG(R)=H^1_{\fppf}(R,G)=H^1_{\et}(R,G)=H^1(\Gal(\overline{R}/R),G(\overline{R})),
%\end{equation} where $\overline{R}$ stands for the separable closure of the field $R$.
%For every $v\in M^0_F-\Sigma$, one can (REFERENCE?) identify $B\mathcal G(\Ov)$ with the unramified cohomology group 
%\begin{equation}\label{unram}
%B\mathcal G(\Ov)=H^1_{\et}(\Ov,G)=H^1_{\un}(\Gal(\overline{F_v}/\Fv),G(\overline{F_v})),
%\end{equation}
%where 
%\begin{multline*}
%%H^1_{\un}(\Gal(\overline{F_v}/\Fv),G(\overline{F_v}))\\:=\ker(H^1(\Gal(\overline{F_v}/\Fv),G(\overline{F_v}))\to H^1(I_v, G(\overline{F_v}))),
%\end{multline*}
%where $I_v$ is the inertia group at the place $v$.
%\subsection{Construction of homomorphism $BG(F_v)\to G_*(\overline F)$}

Let $M$ be a profinite group, let $N$ be a profinite $M$-group and let $J$ be a finite $M$-group. %(not necessarily commutative) 
%(the actions are understood to be continuous). 
We denote by $\Hom(N,J)$ the group of (automatically continuous) homomorphisms $N\to J$. The group $\Hom(N,J)$ is endowed with a canonical structure of an $M$-group:
%s on $\Hom(N,J)$ by 
$$\gamma\cdot f:=\gamma\circ f\circ\gamma^{-1}.$$
%The action respects multiplication in $\Hom(N,J)$ in the sense that $$\gamma\cdot (\phi_1\phi_2)=(\gamma\circ\phi_1\circ\gamma^{-1})\cdot(\gamma\circ \phi_2\circ\gamma^{-1})=(x\mapsto \gamma(\phi_1(\gamma^{-1}(x))\phi_2(\gamma^{-1}(x))).$$  Moreover, the action is continuous  
The group $\Hom(N,J)$ is also endowed with a continuous action of~$J$ given by conjugation: $$j\cdot f:= (x\mapsto (j)(f(x))j^{-1})\hspace{2cm}j\in J, f\in\Hom(N,J).$$
Note that: 
\begin{align*}\gamma\cdot (j\cdot f)&=\gamma \cdot (x\mapsto (j)(f(x))j^{-1})\\&=\big(x\mapsto\gamma\big((j)(f(\gamma^{-1}\cdot x))j^{-1}\big)\big)\\
&=\big(x\mapsto\gamma(j)\gamma(f(\gamma^{-1}\cdot x))\gamma(j)^{-1}\big)\\
&=(\gamma\cdot j)\cdot (\gamma\cdot f),
\end{align*}hence, in particular the action of~$M$ preserves the orbits of~$J$. We deduce that there is a continuous $M$-action on $$\Conj(N,J):=\Hom(N, J)/J$$given by $$\gamma\cdot(fJ)=(\gamma\cdot f)J.$$ The set $\Conj(N,J)$ admits a distinguished element which is fixed by~$M$. The fiber of the distinguished element in $\Conj(N,J)$ is the distinguished element in $\Hom(N,J)$. If~$J$ is assumed commutative, one has that~$\Hom(N,J)=\Conj(N,J)$ admits a structure of an $M$-module. 

Note that if $N\to N'$ is a homomorphism of profinite groups, the canonical map $\Hom(N',J)\to\Hom(N,J)$ is $J$-equivariant and we have a map of pointed sets$$\Conj(N',J)\to\Conj(N, J).$$ If $N\to N'$ is a homomorphism of profinite $M$-groups, one has that the map $\Conj(N',J)\to\Conj(N, J)$ is a morphism of $M$-pointed sets. Similarly, if $J\to J'$ is a homomorphism of finite groups, the canonical map $\Hom(N,J)\to\Hom(N, J')$ is $(J\to J')$-equivariant and we have a map of pointed sets $$\Conj(N, J)\to\Conj(N, J').$$ If $J\to J'$ is a homomorphism of finite $M$-groups, the map $\Conj(N, J)\to\Conj(N, J')$ is a morphism of $M$-pointed sets. %Whenever~$J$ is a finite $M$-group and the group~$N$ is profinite, one has an equality $\Hom_{\cont}(N,J)=\Hom(N,J)$ and we will write $\Conj(N,J)$ for $\Conj_{\cont}(N,J)$.
\subsubsection{}\label{paragtwist} Let now~$M$ be a profinite group and let~$J$ be a finite $M$-group. Let $\sigma\in Z^1(M, J)$. We denote by~$\prescript{}{\sigma}J$ the twist of~$J$ by~$\sigma$ \cite[Section 5.4, Example 2]{Cohomologiegalois}; it is the abelian group~$J$ endowed with with the $\Gamma_F$-action $$\gamma\cdot_{\sigma} f:=(\sigma(\gamma))(\gamma(f))(\sigma(\gamma)^{-1}).$$ 
By \cite[Chapter II, Proposition 35 bis]{Cohomologiegalois}, one has a pointed bijection $$\Lambda_{\sigma}:Z^1(M, \sJ)\to Z^1(M, J),\hspace{1cm}\beta\mapsto \beta\cdot\sigma,$$ which induces a pointed bijection $$\lambda_{\sigma}:H^1(M, \sJ)\to H^1(M, J).$$
Let us recall several properties of twists.
\begin{lem}\label{listproptwist} Let~$\sigma\in Z^1(M,J)$ and let~$\tau\in Z^1(M,\prescript{}{\sigma}J)$.
\begin{enumerate}
\item One has that $\prescript{}{\tau}(\prescript{}{\sigma}J)=\prescript{}{\Lambda_{\sigma}(\tau)}(J).$
\item Let $\sigma'$ be cohomologous to~$\sigma$ and let $k\in J$ be such that $$\sigma'(\gamma)=(k^{-1})(\sigma(\gamma))(\gamma(k))$$ for every $\gamma\in M$. The isomorphism of abelian groups $$\prescript{}{\sigma'}G\to\prescript{}{\sigma} G\hspace{1cm}g\mapsto kgk^{-1}$$is $M$-equivariant (i.e. is an isomorphism of $M$-groups).
\item One has that $$ \Lambda_{\Lambda_{\sigma}(\tau)}=\Lambda_{\sigma}\circ\Lambda_{\tau}$$ and that  $$ \lambda_{\Lambda_{\sigma}(\tau)}=\lambda_{\sigma}\circ\lambda_{\tau}.$$
\item Suppose that $f:N\to M$ is a continuous homomorphism of topological groups. The diagrams  % https://q.uiver.app/?q=WzAsNCxbMCwwLCJaXjEoTSwgXFxzSikiXSxbMSwwLCJaXjEoTixcXHNKKSJdLFsxLDEsIlpeMShOLCBKKSJdLFswLDEsIlpeMShNLCBKKSJdLFswLDFdLFsxLDJdLFszLDJdLFswLDNdXQ==
\[\begin{tikzcd}
	{Z^1(M, \sJ)} & {Z^1(N,\sJ)=Z^1(N, \prescript{}{\sigma\circ f}J)} \\
	{Z^1(M, J)} & {Z^1(N, J)}
	\arrow[from=1-1, to=1-2]
	\arrow["{\Lambda_{\sigma}}"', from=1-1, to=2-1]
	\arrow["{}"', from=2-1, to=2-2]
	\arrow["{\Lambda_{\sigma\circ f}}", from=1-2, to=2-2]
\end{tikzcd}\]
and 
\[\begin{tikzcd}
	{H^1(M, \sJ)} & {H^1(N,\sJ)=H^1(N, \prescript{}{\sigma\circ f}J)} \\
	{H^1(M, J)} & {H^1(N, J)}
	\arrow[from=1-1, to=1-2]
	\arrow["{\lambda_{\sigma}}"', from=1-1, to=2-1]
	\arrow["{}"', from=2-1, to=2-2]
	\arrow["{\lambda_{\sigma\circ f}}", from=1-2, to=2-2]
\end{tikzcd}\]
are commutative.
%\item Suppose that~$R$ is an $M$-invariant subgroup of~$J$. 
\item Let~$K$ be another $M$-group and let $h:J\to K$ be an $M$-equivariant homomorphism. The canonical homomorphism of finite groups $\sJ=J\xrightarrow{h} K=\prescript{}{h\circ\sigma}K$ is $M$-equivariant. The diagrams
% https://q.uiver.app/?q=WzAsNCxbMCwwLCJaXjEoTSwgXFxzSikiXSxbMSwwLCJaXjEoTSxcXHByZXNjcmlwdHt9e2hcXGNpcmNcXHNpZ21hfUspIl0sWzAsMSwiWl4xKE0sSikiXSxbMSwxLCJaXjEoTSxLKSJdLFswLDFdLFswLDIsIlxcTGFtYmRhX3tcXHNpZ21hfSIsMl0sWzIsM10sWzEsMywiXFxMYW1iZGFfe2hcXGNpcmNcXHNpZ21hfSJdXQ==
\[\begin{tikzcd}
	{Z^1(M, \sJ)} & {Z^1(M,\prescript{}{h\circ\sigma}K)} \\
	{Z^1(M,J)} & {Z^1(M,K)}
	\arrow[from=1-1, to=1-2]
	\arrow["{\Lambda_{\sigma}}"', from=1-1, to=2-1]
	\arrow[from=2-1, to=2-2]
	\arrow["{\Lambda_{h\circ\sigma}}", from=1-2, to=2-2]
\end{tikzcd}\]
and
% https://q.uiver.app/?q=WzAsNCxbMCwwLCJaXjEoTSwgXFxzSikiXSxbMSwwLCJaXjEoTSxcXHByZXNjcmlwdHt9e2hcXGNpcmNcXHNpZ21hfUspIl0sWzAsMSwiWl4xKE0sSikiXSxbMSwxLCJaXjEoTSxLKSJdLFswLDFdLFswLDIsIlxcTGFtYmRhX3tcXHNpZ21hfSIsMl0sWzIsM10sWzEsMywiXFxMYW1iZGFfe2hcXGNpcmNcXHNpZ21hfSJdXQ==
\[\begin{tikzcd}
	{H^1(M, \sJ)} & {H^1(M,\prescript{}{h\circ\sigma}K)} \\
	{H^1(M,J)} & {H^1(M,K)}
	\arrow[from=1-1, to=1-2]
	\arrow["{\lambda_{\sigma}}"', from=1-1, to=2-1]
	\arrow[from=2-1, to=2-2]
	\arrow["{\lambda_{h\circ\sigma}}", from=1-2, to=2-2]
\end{tikzcd}\]
are commutative.
\end{enumerate}
\end{lem}
\begin{proof}
\begin{enumerate}
\item  Let $\gamma\in M$ and let $g\in \prescript{}{\tau}(\prescript{}{\sigma}J)=J$. For the canonical action of~$M$ on~$\prescript{}{\tau}(\prescript{}{\sigma}J),$ we have that 
\begin{align*}\gamma\cdot g&=(\tau(\gamma)) (\gamma\cdot_{\sigma}g)(\tau(\gamma))^{-1}\\&=(\tau(\gamma)) \cdot (\sigma(\gamma))\cdot (\gamma(g))\cdot(\sigma(g))^{-1}) \cdot (\tau(\gamma))^{-1}\\
&=(\Lambda_{\sigma}(\tau)(\gamma))(\gamma(g))(\Lambda_{\sigma}(\tau)(\gamma))^{-1},
\end{align*}
which is precisely the element $\gamma\cdot g$ for the canonical action of~$M$ on~$\prescript{}{\Lambda_{\sigma}(\tau)}(J)$. %(We have denote by~$\cdot_{\sigma}$ the canonical action of~$M$ on $\sJ$).
\item Let $\gamma\in M$ and let $g\in G$. We have that\begin{align*}(k)(\gamma\cdot_{\sigma'})(k^{-1})\hskip-2cm&\\&=(k)\big(\sigma'(\gamma)\big)(\gamma(g))\big(\sigma'(\gamma))^{-1}\big)(k^{-1})\\
&=(k)\bigg((k^{-1})(\sigma(\gamma))(\gamma(k))\bigg)(\gamma(g))\bigg((\gamma(k)^{-1})(\sigma(\gamma))^{-1}(k)\bigg)(k^{-1})\\
&=(\sigma(\gamma))(\gamma(k))(\gamma(g))(\gamma(k^{-1}))(\sigma(\gamma))^{-1}\\
&=\gamma\cdot_{\sigma}(kgk^{-1}),
\end{align*}
and the claim follows.
\item Let $\theta\in Z^1(M,\prescript{}{\tau}(\prescript{}{\sigma}J))$. We have that $$\Lambda_{\Lambda_{\sigma}(\tau)}(\theta)=\theta\cdot \Lambda_{\sigma}(\tau)=\theta\cdot \tau\cdot\sigma\in Z^1(M, J).$$ On the other side $$\Lambda_{\sigma}(\Lambda_{\tau}(\theta))=\Lambda_{\sigma}(\theta\cdot\tau)=\theta\cdot\tau\cdot\sigma.$$ The first claim follows. The second claim follows immediately from the first claim.
\item The first diagram is commutative, because the image of $\theta\in Z^1(M, \sJ)$ under the composite of left vertical and upper horizontal maps is $\Lambda_{\sigma\circ f}(\theta\circ f)=(\theta\circ f)\cdot (\sigma\circ f),$ which is precisely $(\theta\cdot \sigma)\circ f=\Lambda_{\sigma}(\theta)\circ f.$ The commutativity of the second diagram follows immediately from the commutativity of the first diagram.
\item Let $j\in J$ and let $\gamma\in M$. We have that 
\begin{align*}
h(\gamma\cdot_{\sigma}j)&=h\big((\sigma(\gamma))\cdot(\gamma(j))\cdot(\sigma(\gamma))^{-1}\big)\\
&=h(\sigma(\gamma))\cdot h(\gamma(j))\cdot (h(\sigma(\gamma)))^{-1}\\
&=((h\circ\sigma)(\gamma))\cdot\gamma(f(j))\cdot((h\circ\sigma)(\gamma))^{-1}\\
&=\gamma\cdot_{h\circ\sigma} f(j),
\end{align*}
and the first claim follows. If $\theta\in Z^1(M, \sJ)$, its image under the composite of the left vertical and upper horizontal maps is given by $\Lambda_{h\circ\sigma}(h\circ \theta)=(h\circ\theta)\cdot (h\circ \sigma).$ This is equal to $h\circ (\theta\cdot\sigma)=h\circ(\Lambda_{\sigma}(\theta))$ and the commutativity of the first diagram follows. The commutativity of the second diagram is now immediate.
\end{enumerate}
\end{proof}
%In situation of Part (5) of the lemma, we may write $\prescript{}{\sigma}K,$  for $\prescript{}{h\circ\sigma}K$. We may also write~$\Lambda_{\sigma}$ for~$\Lambda_{h\circ\sigma}$ and~$\lambda_{\sigma}$ for~$\lambda_{h\circ\sigma}$.
%\begin{lem} \label{doslemtwis}
%Suppose that $\phi_A:J\hookrightarrow A$ and $\phi_B:J\hookrightarrow B$ are $M$-equivariant inclusions to finite $M$-groups. We suppose that~$\phi_A(J)$ and~$\phi_B(J)$ are normal. Let $f:A\to B$ be an $M$-equivariant homomorphism such that $\phi_B=f\circ\phi_A$. Let $\sigma_A\in Z^1(M, A)$ and let $\sigma_B:=f\circ\sigma_A$. We define $\prescript{}{\sigma_A}J$ to be the group~$J$ endowed with the following $M$-action $$\gamma\cdot_{A} j=\phi_A^{-1}\big((\sigma_A(\gamma))(\phi_A(\gamma\cdot j)) (\sigma_A(\gamma))^{-1}\big).$$
%(here $\phi_A^{-1}:\phi_A(J)\to J$ is the inverse of~$\phi_A$.)
%\begin{enumerate}
%\item  One has that $\prescript{}{\sigma_A}J=\prescript{}{\sigma_B}J$ as $M$-groups.
%\item The map $\prescript{}{\sigma_A}J=J\xrightarrow{\phi_A}A=\prescript{}{\sigma_A}A$ of finite groups is $M$-equivariant. 
%\item One has that $\Lambda_{\sigma_A}=\Lambda_{\sigma_B}$ and $\lambda_{\sigma_A}=\lambda_{\sigma_B}.$
%\end{enumerate} 
%\end{lem}
%\begin{proof} Let $\gamma\in M$ and $j\in J.$
%\begin{enumerate}
%\item One has that
%\begin{align*}
%\gamma\cdot_Bj&=\phi_B^{-1}\big((\sigma_B(\gamma))(\phi_B(\gamma\cdot j)) (\sigma_B(\gamma))^{-1}\big)\\&=(f\circ \phi_A)^{-1}\big((f(\sigma_A(\gamma)))(f(\phi_A(\gamma\cdot j))) (f(\sigma_A(\gamma)))^{-1}\big)\\
%&=\phi_A^{-1}\big((\sigma_A(\gamma))(\phi_A(\gamma\cdot j))(\sigma_A(\gamma))^{-1}\big)\\
%&=\gamma\cdot_Aj.
%\end{align*}
%\item One has that
%\begin{align*}
%\phi_A(\gamma\cdot _A j)&=\phi_A\big(\phi_A^{-1}\big((\sigma_A(\gamma))(\phi_A(\gamma\cdot j))(\sigma_A(\gamma))^{-1}\big)\big)\\
%&=(\sigma_A(\gamma))(\gamma\cdot \phi_A(j)) (\sigma_A(\gamma))^{-1}\\
%&=\gamma\cdot_{\sigma_A}\phi_A(j).
%\end{align*}
%\end{enumerate}
%\end{proof}
\subsubsection{}We now turn to~$\Gamma_F$-groups. %Thanks to Subsection \ref{equivalenceofcat}, this is easily translated into the language of \' etale group schemes. 
When $k\geq 1$ is an integer, we denote by~$\muu_k$ the group of~$k$-th roots of unity in~$\overline{F}$. It has a structure of~$\Gamma_F$-module: $$\Gamma_F\times\muu_k\to\muu_k\hspace{1cm}(\gamma,\xi)\mapsto \xi^{\chi_{\cycl}(\gamma)},$$where $\chi_{\cycl}:\Gamma_F\to\prod_{\ell\neq \charr(F)}\ZZ_{\ell}^{\times}$ is the cyclotomic character \cite[Definition 7.3.6]{CohomologyNF}. Let~$J$ be a finite $\Gamma_F$-group.
\begin{mydef}
Let $e=e(J)$ be the exponent of~$J$. %The group $\Gamma_F$ acts on the group $\muu_e$ via the cyclotomic character $\chi_{\cycl}:\Gamma_F\to\widehat{\ZZ} $: $$\Gamma_F\times\muu_e\to\muu_e \hspace{1cm}(\gamma,\xi_e)\mapsto \xi_e^{\chi_{\cycl}(\gamma)}.$$
We set $$J_*:=\Conj(\muu_{e},J)$$ and we endow it with the continuous $\Gamma_F$-action as above. If~$J$ is abelian, then~$J_*$ is understood to be endowed with a structure of a commutative $\Gamma_F$-module.
\end{mydef}
\begin{rem}
\normalfont
When~$J$ is abelian, the $\Gamma_F$-module~$J_*$ is precisely the Tate twist $J(-1)$ \cite[Section 7.8]{cohomologicalinvariants}.
\end{rem}
\begin{lem}\label{slabalema}
\begin{enumerate}
\item Suppose that $J$ is endowed with a trivial~$\Gamma_F$-action.  Let $\Conj_F(J)$ be the set of $F$-conjugacy classes of~$J$ (\cite[Section 2]{Malle})\footnote{In \cite{Malle}, Malle works only for the case~$F$ is a number field. The definitions and results that we use, are generalized to the case of a global field in an evident way.}. A choice of be a primitive $e=e(J)$-th root of unity $\xi$ in $\overline F$ induces a map of pointed sets $$\lambda_{\xi}:\Hom(\muu_e, J)\to J\to\Conj_F(J),\hspace{1cm} f\mapsto [f(\xi)]$$ where $[g]$ denotes the $F$-conjugacy class of~$g$. The map $\lambda_{\xi}$ is surjective, of trivial kernel, $J$-invariant, $\Gamma_F$-invariant and induces a bijection $$J_*/\Gamma_F\xrightarrow{\sim}\Conj_F(J).$$ If~$r\geq 1$ is an integer such that $\gcd(r,e)=1,$ then for $y\in \Hom(\muu_e, J)$ one has that $\lambda_{\xi}(y)=\lambda_{\xi^r}(y)^r,$ where the map of pointed sets $\Conj_F(J)\to\Conj_F(J)$ denoted by $x\mapsto x^r$ is the map induced from $J$-equivariant and $\Gamma_F$-equivariant automorphism $J\to J$ given by $j\mapsto j^r$.  %The induced map $\lambda_\xi:J_*/\Gamma_F\to\Conj_F(J)$ is a bijection of pointed sets. 
%\item Suppose again that $J$ is constant. Let $r\geq 1$ be such that $\gcd(r,n)=1$. One has that $\lambda_{\xi^r}(x)=\lambda_{\xi}(x) ^r$, where $y\mapsto y^r$ is the bijection $\Conj_F(J)\to \Conj_F(J)$ induced from $J$-equivariant and $\Gamma_F$-equivariant automorphism $J\to J$ which is given by $h\mapsto h^r$.
\item Suppose that $J=~\muu_k$ for some integer~$k\geq 1$. The~$\Gamma_F$-action on~$(\muu_k)_*$ is trivial.
\end{enumerate}
\end{lem}
\begin{proof}
\begin{enumerate}
\item Note that $f\mapsto f(\xi)$ is a group isomorphism $$\Hom(\muu_e, J)\xrightarrow{\sim} J.$$ (It is an injective group homomorphism, which is surjective, because an element $g\in J$ of order~$k$ is the image of the composite homomorphism $\muu_e\to\muu_k\to J$ which is given by $\xi\mapsto \xi^{e/k}\mapsto g$). We endow~$J$ with a $\Gamma_F$-action using the isomorphism. For this action, one has that $\gamma\cdot g=g^{\chi_{\cycl}(\gamma)^{-1}}$, thus, this action is the opposite action of the one given in \cite[Section 2]{Malle}. Hence, the~$J$-invariant, surjective map of pointed sets $J\to \Conj_F(J)$ with the trivial kernel is~$\Gamma_F$-invariant.  Let us prove the last claim.  The map $f\mapsto f(\xi^r)$ factorizes as the composite of the maps $\Hom(\muu_e,J)$ given by $f\mapsto f(\xi)$ and the map $J\to J$ given by $x\mapsto x^r$. Hence, $\lambda_{\xi^r}(f)=[f(\xi^r)]=[f(\xi)^r]=[f(\xi)]^r=\lambda_\xi(f)^r,$ as claimed
\item This is a well known fact, but we include proof for completeness. Let~$\xi$ be a primitive~$k$-th root of unity in~$\oF$. A homomorphism $f:\muu_k=\muu_{e(\muu_k)}\to \muu_k$ is determined by the image on~$\xi$. We have that $$(\gamma\cdot f)(\xi)=\gamma(f(\gamma^{-1}(\xi)))=f(\xi^{\chi_{\cycl}(\gamma)^{-1}})^{\chi_{\cycl}(\gamma)}=f(\xi),$$i.e. $\gamma\cdot f=f$.
\end{enumerate}
\end{proof}
For positive integers $d,k$ with $d|k$, one has a canonical $\Gamma_F$-equivariant surjective homomorphism $$q_{k,d}:\muu_k\to\muu_d\hspace{1cm}x\mapsto x^{k/d}.$$
\begin{lem}\label{funcinj}
Let $\iota:N\hookrightarrow J$ be an injective $\Gamma_F$-equivariant homomorphism (in particular $e(N)|e(J)$). The map $$\Hom(\muu_{e(N)}, N)\to\Hom(\muu_{e(J)},J)\hspace{1cm}f\mapsto \iota\circ f\circ q_{e(J), e(N)}$$ is an injective homomorphism which is $\iota$-equivariant and $\Gamma_F$-equivariant. The induced map of pointed sets $\iota_*:N_*\to J_*$ is $\Gamma_F$-equivariant and of trivial kernel.
\end{lem}
\begin{proof}
The map $\Hom(\muu_{e(N)}, N)\to\Hom(\muu_{e(J)},J)$ is the composite of the injective homomorphisms \begin{equation}\label{compnj}\tag{*}\Hom(\muu_{e(N)}, N)\to\Hom(\muu_{e(J)}, N)\to\Hom(\muu_{e(J)}, J), \end{equation} hence is an injective homomorphism. The first map in (\ref{compnj}) is $N$-invariant. The image of $f\in\Hom(\muu_{e(J)},N)$ for the second map in (\ref{compnj}) is $\iota\circ f$, while for $g\in N$, the image of $g\cdot f=(x\mapsto (g)f(x)(g)^{-1})$ is given by $$\big(x\mapsto \iota((g)f(x)(g)^{-1})\big)=\iota(g)\cdot (\iota\circ f).$$ It follows that the second map in (\ref{compnj}) is $\iota$-equivariant. Hence, the canonical map $\Hom(\muu_{e(N)},N)\to\Hom(\muu_{e(J)},J)$ is $\iota$-equivariant. By the fact that $\muu_{e(J)}\to\muu_{e(N)}$ is $\Gamma_F$-equivariant, it follows that the first map in (\ref{compnj}) is $\Gamma_F$-equivariant and by the fact that $\iota$ is $\Gamma_F$-equivariant, it follows that the second map in (\ref{compnj}) is $\Gamma_F$-equivariant. Hence, $\Hom(\muu_{e(N)}, N)\to\Hom(\muu_{e(J)}, J)$ is $\Gamma_F$-equivariant. It follows that the map $\iota_*:N_*\to J_*$ is $\Gamma_F$-equivariant. The map $\Hom(\muu_{e(N)}, N)\to\Hom(\muu_{e(J)}, J)\to J_*$ is of the trivial kernel as the composite of such maps, hence is such $N_*\to J_*$. The statement is proven. 
\end{proof}
We now verify that twisting does not change~$J_*$ and some other properties. 
\begin{lem}\label{twistsandstar} Let $\sigma\in Z^1(\Gamma_F, J)$.
\begin{enumerate}
\item One has an equality of $\Gamma_F$-pointed sets $(\prescript{}{\sigma}J)_*=J_*.$
\item Let $R\subset J$ be a $\Gamma_F$-invariant subgroup. Let $\tau\in Z^1(\Gamma_F, R)$. The canonical map $R_*\to J_*$ identifies with the map $(\prescript{}{\tau}R)_*\to (\prescript{}{\tau}J)_*.$
\item Let~$\sigma'\in  Z^1(\Gamma_F, J)$ be cohomologous to~$\sigma$ and let $k\in J$ be such that for every $\gamma\in\Gamma_F$ one has that $\sigma'(\gamma)=(k^{-1})(\sigma(\gamma))(k).$ The isomorphism $$f_k:\prescript{}{\sigma'}J\to\sJ\hspace{1cm}j\mapsto kjk^{-1},$$which is $\Gamma_F$-equivariant by Part (2) of Lemma \ref{listproptwist}, induces the identity map $(f_k)_*:J_*=(\prescript{}{\sigma'}J)_*\xrightarrow{\sim} (\sJ)_*=J_*$. If~$N$ is a $\Gamma_F$-invariant subgroup of~$\prescript{}{\sigma'}J$, the diagram % https://q.uiver.app/?q=WzAsMyxbMCwwLCJNXyoiXSxbMCwxLCIoZl9rKE0pKV8qIl0sWzEsMSwiSl8qIl0sWzAsMV0sWzAsMl0sWzEsMl1d
\[\begin{tikzcd}
	{N_*} \\
	{(f_k(N))_*} & {J_*}
	\arrow[from=1-1, to=2-1]
	\arrow[from=1-1, to=2-2]
	\arrow[from=2-1, to=2-2]
\end{tikzcd}\] is commutative in which the vertical map is $\Gamma_F$-equivariant bijection of pointed sets.
\end{enumerate}
\end{lem}
\begin{proof}
\begin{enumerate}
\item Clearly, one has an equality of pointed sets $\Hom(\muu_e, \prescript{}{\sigma}J)=\Hom(\muu_e, J)$ and hence of pointed sets $(\prescript{}{\sigma}J)_*=J_*$. We prove that the Galois actions are identical. Let~$\gamma\in\Gamma_F$ and let~$j\in (\prescript{}{\sigma}J)_*$. Let~$\widetilde j\in \Hom(\muu_e, \prescript{}{\sigma}J)$ be a lift of~$j$. We have that $$\gamma\cdot_{\sigma} j=[(\sigma(\gamma))(\gamma(\widetilde j))(\sigma(\gamma))^{-1}]=[\gamma(\widetilde j)]=\gamma\cdot j,$$where %here~$\cdot^{\sigma}$ stands for the $\Gamma_F$-action on~$(\prescript{}{\sigma}J)_*$ and 
$[-]$ stands for the equivalence class in~$(\prescript{}{\sigma}J)_*=J_*$ of an element. The statement is proven.
\item This follows from Part (1) and the fact that the canonical map $R_*\to J_*$ is defined purely in the terms of groups (and does not depend on Galois actions). %Let $r\in R_*=(\prescript{}{\tau}R)_*$. Its image for the canonical map $R_*\to J_*$ is given by the image of a lift $\widetilde r\in \Hom(\muu_{e(R)}, R)$ for the canonical map $\Hom(\muu_{e(R)}, R)\to \Hom(\muu_{e(J)}, J)\to J_*.$
\item Let $j\in (\prescript{}{\sigma'}J)_*=J_*$ and let $\widetilde j\in \Hom(\muu_e,J)$ be its lift. One has that~$(f_k)_*(j)$ is the image of~$(k)(\widetilde j)(k)^{-1}$ under the quotient map $\Hom(\muu_e, J)\to J_*=(\prescript{}{\sigma}J)_*$ and this equals to the image of~$\widetilde j$ under the quotient map, hence $(f_k)_*(j)=j$. Let us prove the second claim. Let $g\in N_*$ and let~$\widetilde g\in N$ be its lift. The image of~$g$ in~$J_*$ equals to the class of~$\widetilde g\in J$. On the other side, the image of~$g$ in~$(f_k(N))_*$ equals to the image of $k\widetilde gk^{-1}$ in~$f_k(N)$. The map $(f_k(N))_*\to J_*$ takes the class of $k\widetilde gk^{-1}$ to the class of~$k\widetilde gk^{-1}$ in~$N_*$, which coincides with the class of $\widetilde g$. The second claim is proven.
\end{enumerate}
\end{proof}
\subsubsection{}Let~$G$ be a finite \' etale $F$-group scheme. Set $J=G(\oF)$ and endow it with the usual $\Gamma_F$-action. The~$\Gamma_F$-pointed set~$J_*$ defines a pointed \' etale scheme~$G_*$. If~$G$ is commutative, then $G_*$ has a structure of a commutative group scheme.  If $\sigma\in Z^1(\Gamma_F, G(\oF))$, we denote by~$\sG$ the finite \' etale tame group scheme defined by the $\Gamma_F$-group $\prescript{}{\sigma}(G(\oF)).$
\subsection{Counting function}
%Let $F$ be a global field. Let $G$ be a commutative finite \' etale group scheme over $F$ of order prime to the order of the field $F$.
%Let $e$ be the exponent of $G$. By $\muu_e$ we denote the group of $e$-th roots of $1$ in $\overline{F}$. The absolute Galois group $\Gamma_F$ acts on $\muu_e$ via the cyclotomic character $\chi_{\cycl}:\Gamma_F\to\widehat{\ZZ} $: $$\Gamma_F\times\muu_e\to\muu_e \hspace{1cm}(\gamma,\xi_e)\mapsto \xi_e^{\chi_{\cycl}(\gamma)}.$$ Moreover, the group $\Gamma_F$ acts (to the left) on the set of $\overline{F}$-points $G(\overline{F})$. 
%\begin{mydef}
%We define a $\Gamma_F$-module $G_*(\overline F)$ by setting it to be $$G_*(\overline F):=\Hom(\muu_e,G(\overline{F}))$$ where the action of $\Gamma_F$ is given by $$\gamma\cdot f:= \gamma\circ f\circ \gamma^{-1}. $$
%\end{mydef}
%Note that as abelian groups, $G(\overline{F})$ and $G_*(\overline F)$ are isomorphic (an isomorphism is provided by $g\mapsto (\xi_e\mapsto g)$, where $(\xi_e\mapsto g)$ denotes the unique element of $\Hom (\muu_e, G(\overline{F}))$ defined by mapping $\xi_e$ to $g$). 
In this paragraph we define counting functions. They are used to define heights. When~$F=\QQ$ and~$G$ is constant, the definition coincides with the one of \cite[Section 4.2]{Gmethodant}. 
\subsubsection{}In this paragraph, we will define {\it discriminant counting function}. It ``induces" discriminants of $G$-torsors as will see in Subsection~\ref{discofgtorsor}.
\begin{mydef}
Let $J$ be a non-trivial finite $\Gamma_F$-group. A counting function is $\Gamma_F$-invariant function $$c:J_*\to\RR_{\geq 0},$$ which satisfies $c(g)=0$ if and only if~$g=1_{J_*}$ is distinguished in~$J_*.$ We define $$a(c):=\big(\min_{g\in J_*-\{1_{J_*}\}}\{c(g)\}\big)^{-1}\in\RR_{>0}.$$ The function~$c$ is said to be normalized if $a(c)=1.$
\end{mydef}
For any counting function $c:J_*\to\RR_{\geq 0}$, the function %its {\it normalization}~$c^N$ is defined by 
$$J_*\to\RR_{\geq 0},\hspace{1cm} g\mapsto \frac{c(g)}{a(c)}$$ is a normalized counting function.
\begin{mydef}
For $q\in\RR_{\geq 0}$, we denote by $(J_*)_q$ the $\Gamma_F$-invariant subset of the pointed set~$J_*$ given by $(J_*)_q:=c^{-1}(q).$ We define $$b(c):=\#  \big(((J_*)_{a(c)^{-1}})/\Gamma_F\big).$$
\end{mydef}
In the rest of paragraph we present several examples. %For $f\in\Hom(\muu_e,J)$, we denote by $\ord(f)$ its order.
 We have a function 
\begin{equation}
\Hom(\muu_e,J)\to\QQ_{\geq 0},\hspace{1cm} f\mapsto [J:f(\muu_e)]\cdot (\# f(\muu_e)-1).
%\#J\cdot \bigg(1-\frac{1}{\ord(f)}\bigg)
\label{discount}\tag{*}
\end{equation}The function is $J$-invariant and $\Gamma_F$-invariant as the orders and indices of subgroups are preserved under group automorphisms. Moreover, it takes the value~$0$ at the trivial homomorphism $\muu_e\to J$ and takes strictly positive values elsewhere.
\begin{mydef}
Let $c_{\Delta}:J_*\to\QQ_{\geq 0}$ be the counting function induced by the $J$-invariant function (\ref{discount}). We call it the discriminant counting function.
\end{mydef}
\subsubsection{} In this paragraph, we give another example of a counting function. The corresponding height is used in the statement of the Malle conjecture.

Let~$J$ be a non-trivial constant $\Gamma_F$-group. We suppose that there is a transitive embedding $\iota:J\hookrightarrow \mathfrak S_n$ for certain $n\geq 1$, where~$\mathfrak S_n$ stands for the group of permutations of the set $\{1\doots n\}$. For $g\in \Conj_F(J),$ let $\widetilde g\in J$ be a lift of $g$. We set
\begin{align*}\ind^{\iota}:\Conj_F(J)&\to\ZZ_{\geq 0}\\g&\mapsto\ind(\iota(\widetilde g)):=n-\# \{\textnormal{orbits of }\iota(\widetilde g)\textnormal{ in \{1\doots n\}}\}.
\end{align*} The definition does not depend on the choice of $\widetilde g$ and for $g\neq 1\in\Conj_F(J)$ one has $\ind(g)\neq 1$, as stated in \cite[Section 2]{Malle}. Note that if $r\geq 1$ is an integer such that $\gcd(r,e)=1$, then $\ind^{\iota} (g)=\ind^{\iota}(g^r)$. For a primitive $e=e(J)$-th root of unity~$\xi,$ we have a $\Gamma_F$-invariant surjection of pointed sets $\overline\lambda_\xi:J_*\to \Conj_F(J)$ induced from $J$-invariant and $\Gamma_F$-invariant surjection $\lambda_\xi:\Hom(\muu_e,J)\to \Conj_F(J)$ which is defined in Lemma \ref{slabalema}. Moreover, one has that $\overline\lambda_{\xi^r}=\overline\lambda_{\xi}^r$ and that $\overline{\lambda_{\xi}}$ is of trivial kernel. It follows from above that the map $$c_{{\iota}}:J_*\to\QQ_{\geq 0}\hspace{1cm}c_{{\iota}}:=\ind^{\iota}\circ\overline{\lambda}_\xi$$ is a counting function and that does not depend on the choice of $\xi$. 
%\begin{mydef}  We define a counting function $$c_{\ind}:J_*\to \ZZ_{\geq 0}$$ by setting it to be the composite of the $\Gamma_F$-invariant map of pointed sets $J_*\to \Conj_F(J)$ and the function $\ind$. 
%\end{mydef}
It is immediate that one has that $$a(c_{{\iota}})=a'(J\leq \mathfrak S_n),$$ and $$b(c_{{\iota}})=b'(F,J\leq \mathfrak S_n),$$where $a'(-)$ and $b'(-,-)$ are the standard invariants in the Malle conjecture \cite[Section 2]{Malle}.
%\begin{proof}Let $g\in \Hom(\muu_e,J)$. For $j\in J$, the permutations $(j\cdot g)(\xi)=(j) (g(\xi)) j^{-1}$ have the same cycle structure, as cycle structures are preserved under conjugation in~$\mathfrak S_n$. It follows that $\ind(g)=\ind(j\cdot g)$, i.e. that the function~$\ind$ is $J$-invariant. Moreover, for $\gamma\in\Gamma_F$, one has that $$(\gamma\cdot g)(\xi)=g(\xi^{\chi_{\cycl}(\gamma)^{-1}})=g(\xi)^{\chi_{\cycl}(\gamma)^{-1}}.$$
%\end{proof}
%AGE IS DEFINED HERE\begin{mydef}\label{defage} Let $\xi$ be an $e$-th primitive root of unity. Let $\rho: G_*(\overline F)\to \GLn(\overline F)$ be a representation of $G_*(\overline F).$  For $g\in G_*(\overline F)$, let $q_{g,1}\doots q_{g,n}$ be integers in the domain $[0, e-1]$ such that $\xi^{q_{g,1}}\doots \xi^{q_{g, n}}$ is the set of eigenvalues (with repetitions) of $\rho (g)$. We define $$\age^{\xi}_{\rho}:G_*(\overline F)\to\QQ_{\geq 0}\hspace{1cm}\age^{\xi}_{\rho}(g):=\frac{1}{e}\sum_{i=1}^nq_{g,i}.$$
%\end{mydef}
%IT IS PROVEN THAT AGE IS A COUNTING FUNCTION\begin{lem} Let $\xi$ be an $e$-th primitive root of unity. Let $\rho:G_*(\overline F)\to \GLn(\overline F)$ be a faithful $\overline F$-representation of $G_*(\overline F)$ which is $\Gamma_F$-equivariant, where $\Gamma_F$ acts on $\GLn(\overline F)$ via its action on $\overline F^{n^2}$. The function $\age^{\xi}_{\rho}:G_*(\overline F)\to\QQ_{\geq 0}$ is a counting function. 
%\end{lem}
%\begin{proof}Note that $\age^\xi_{\rho}(g)=0$ if and only if every $q_{i,g}=0$ (where $q_{i,g}$ are integers as in Definition \ref{defage}). This is true if and only if $\rho(g)=1\in\GLn(\overline F)$, and by faithfulness, if and only if $g=1$. Let us now verify that $\age^\xi_\rho$ is $\Gamma_F$-invariant. Let $g\in G_*(\overline F)$, by the definition, the element~$g$ is a homomorphism $\muu_e\to G(\overline F)$ such that $\xi\mapsto g$. If we denote by $\lambda:\Gamma_F\to\prod_{\ell\neq\charr(F)}\ZZ_{\ell}$ the cyclotomic character, then \begin{align*}
%\rho(\gamma\cdot g)&=\rho((\gamma \circ g\circ \gamma^{-1} )(\xi))\\&=(\rho\circ \gamma) (g \circ (\gamma^{-1}(\xi)))\\&=(\gamma\circ\rho)(g(\xi^{\lambda(\gamma)^{-1}}))\\&=\gamma (\rho( g(\xi^{\lambda(\gamma)^{-1}})))\\&=\gamma (\rho(g(\xi)^{\lambda(\gamma)^{-1}}))\\
%&=\gamma(\rho(g(\xi)))^{\lambda(\gamma)^{-1}}).
%\end{align*}
%The eigenvalues of $\rho(g(\xi))$ are given by $\{\xi^{\lambda(\gamma)^{-1}q_{i,g}}\}_i$ and it follows that the eigenvalues of $\rho((\gamma\cdot g))$ are given by $$\gamma(\{\xi^{\lambda(\gamma)^{-1}q_{i,g}}\}_i)=\{\gamma(\xi^{\lambda(\gamma)^{-1}q_{i,g}})\}_i=\{\xi^{q_{i,g}}\}_i.$$ In particular, $$\age^{\xi}_{\rho}(\gamma\cdot g)=\age^{\xi}_{\rho}(g).$$The statement is proven.
%\end{proof}
%\begin{rem}
%\normalfont
%By SGA, the $\Gamma_F$-module $G_*(\overline F)$ defines a finite \' etale group scheme $G_*$ over $F$. The condition that homomorphism $\rho: G_*(\overline F)\to\GLn(\overline F)$ in \ref{} is $\Gamma_F$-invariant, is equivalent to the condition that it is the  induced map on $\overline F$-points of a base change of an $F$-homomorphism $G_*\to \GLn$.
%\end{rem}
\subsection{Residue map} Let~$G$ be a non-trivial finite \' etale $F$-group scheme. Let~$v$ be a finite place of~$F$ such that $\Gamma_v$-group $G(\oF)=G(\overline{F_v})$ is unramified (equivalently, the $\Fv$-group scheme~$G_{\Fv}$ is the base change of an \' etale $\Ov$-group scheme) and such that $G$ is tame at~$v$ (by this we mean that $\gcd(\#G(\oF),q_v)=1$, where $q_v=~\#\kappa_v$).  We will define a map of pointed sets $\psi_v:H^1(\Gamma_v, G(\oF))\to G_*(\Fv)$. If~$G$ is commutative, the map is a group homomorphism. The maps are used to define heights in Subsection~\ref{secdefheight}.

We denote by~$\Fv^{\tame}$ %and by $\Fv^{\un}$ 
the maximal tame extension %and the maximal unramified extension of $\Fv,$ respectively, 
contained in $\overline{F_v}$ and by $\Gamma_v^{\tame}$ its Galois group over~$F_v$. %We denote by~$\Gamma_v:=\Gal(\overline{F_v}/F_v)$ the absolute Galois group of $\Fv$, by $\Gamma_v^{\un}:=\Gal(F_v^{\un}/F_v)$ the Galois group of the maximal unramified extension of $\Fv$ and by $\Gamma_v^{\tame}:=\Gal(F_v^{\tame}/F_v)$ the Galois group of the maximal tame extension of $F_v$.
Let~$T_v$ be the ramification group $$T_v:=\Gal(\overline{F_v}/\Fv^{\tame})=\ker(\Gamma_v\to\Gamma_v^{\tame})$$ and $I_v$ the inertia group $$I_v:=\Gal(\overline{\Fv}/\Fv^{\un})=\ker(\Gamma_v\to\Gamma_v^{\un}).$$ Let $I_v^{\tame}$ be the tame inertia group, that is $$I_v^{\tame}:=\Gal(F_v^{\tame}/F_v^{\un})=\coker(T_v\to I_v).$$
The sets %$H^1(F_v, G)=
$H^1(\Gamma_v, G(\oF))$ are finite (by \cite[Theorem 7.5.14]{CohomologyNF}, the group~$\Gamma_v$ is a topologically finitely generated group, and thus by \cite[Chapter III, \S 4, Proposition 9]{Cohomologiegalois}, the pointed set $H^1(\Gamma_v, G(\oF))$ is finite).
\subsubsection{} We will prove and recall several auxiliary results. As~$\Gfs$ is unramified at~$v$, the group~$I_v$, and hence the group~$T_v,$ act trivially on it. We deduce, in particular, that~$\Gamma_v^{\tame}$ acts  on~$\Gfs$. We let~$p_v$ be the characteristic of the residue field~$\kappa_v$.
\begin{lem}
%Let $\vMFz$ be a place such that $\Gfs$ is unramified at~$v$.
The inflation map $H^1(\Gamma_v^{\tame}, \Gfs)\to H^1(\Gamma_v, \Gfs)$ is a bijection of pointed sets. If $G$ is commutative, then the bijection is an isomorphism of abelian groups.
\end{lem}
\begin{proof}We denote by~$p_v$ the characteristic of~$\kappa_v$. By using that $\gcd(\# \Gfs,p_v)=1$ and the fact that %, where $p$ is the characteristic of the residue field $\kappa_v$. 
the group $T_v$ is a pro-$p_v$-group \cite[Proposition 7.5.1]{CohomologyNF}, we obtain that $\Hom_{\cont}(T_v,\Gfs)=\{1\}.$ %Moreover, by using again that $\Gfs$ is unramified at $v$, we have that the group $T_v\subset I_v$ acts trivially on $\Gfs$ and we obtain that 
We deduce that one has $H^1(T_v, \Gfs)=~\{1\}.$ %where $\Hom_{\cont}$ stands for the set of continuous homomorphisms.
Now, the inflation-restriction exact sequence \cite[Chapter I, \S 5.8 (a)]{Cohomologiegalois}, applied to the inclusion $T_v\subset~\Gamma_v,$ reads
$$1\to H^1(\Gamma_v^{\tame}, \Gfs)\to H^1(\Gamma_v, \Gfs)\to 1.$$ If $G$ is commutative, then by \cite[Chapter I, \S 2.6 (b)]{Cohomologiegalois}, the exact sequence is an exact sequence of abelian groups and the bijection is an isomorphism of abelian groups.
\end{proof}
The conjugation in~$\Gamma_v^{\tame}$ defines an action on its abelian normal subgroup~$I_v^{\tame}$ of the group~$\Gamma_v^{\un}=\Gamma_v^{\tame}/I_v^{\tame}$. If we denote by~$\sigma$ the canonical topological generator of~$\Gamma_v^{\un}$, by \cite[Page 410]{CohomologyNF}, the action is determined by $\sigma\cdot \gamma=~\gamma^{q_v},$ where $\gamma\in\Gamma_v^{\un}$. For $k\geq 1$ with $\gcd(k, q_v)=1$, there is a canonical identification of the group of $k$-th roots of unity~$\muu_k$ in~$\overline {F_v}$, with the Galois group of the unique tame totally ramified extension of degree~$k$ of~$F_v^{\un}$ (it is obtained by mapping an automorphism $\gamma$ to $\gamma(\piv^{1/k})/\piv^{1/k}$ for a choice $\piv^{1/k}$ of $k$-th root of an uniformizer~$\piv$, the result is independent of the choices). The group~$\Gamma_v^{\un}$ acts on~$\muu_k$ by $\sigma\cdot \xi=\xi^{q_v}$ and it follows that the canonical homomorphism $I_v^{\tame}\to \muu_k$ is $\Gamma_v^{\un}$-equivariant. 
  
%There exists a canonical I AM NOT SURE isomorphism of $\Gamma_v^{\un}$-modules $$I_v^{\tame}\xrightarrow{\sim}\widehat{\muu}_v^{(p)},$$where $\widehat{\muu}_v^{(p)}$ is the projective limit of $\Gamma_v$-modules $$\widehat{\muu}_v^{(p)}=\varprojlim _{\gcd(p,n)=1}\muu_n.$$
\begin{lem}\label{indmue} The canonical map of $\Gamma_v^{\un}$-groups and $G(\oF)$-groups
\begin{equation}
\label{oglic}
\Hom(\muu_e, G(\oF))\to\Hom (I_v^{\tame}, G(\oF))
\end{equation} %$$\Hom(I_v^{\tame}, G(\overline F))\to\Hom(\muu_e, G(\overline F))\hspace{1cm} f\mapsto f_e$$
is a bijection. The induced map of $\Gamma_v^{\un}$-pointed sets $$G_*(\oF)=\Conj(\muu_e, G(\oF))\to\Conj(I_v^{\tame}, G(\oF))=H^1(I_v^{\tame}, G(\oF))$$ is a bijection. %if $\Hom(I_v^{\tame}, G(\overline F))$ is endowed with the following $\Gamma_v^{\un}$-action: $$\Gamma^{\un}_v\times \Hom(I_v^{\tame}, G(\overline F))\to \Hom(I_v^{\tame}, G(\overline F))\hspace{1cm}(\gamma, f)\mapsto \gamma\circ f\circ\gamma^{-1}.$$
If~$G$ is commutative, then the bijections are isomorphisms of abelian groups.
\end{lem}
\begin{proof}
%By \cite[Chapter II, \S 1,\no 9, Corollary 2 of Proposition 11]{TSpectrale}, the Pontryagin dual of $\muup=\varprojlim_{\gcd(p,n)=1}\mu_{n}$ identifies with the group $$\bigcup_{\gcd(n,p)=1}\mu_n.$$ In particular, all of its finite subgroups are cyclic.
The canonical homomorphism $I_v^{\tame}\to\muu_e$ is surjective, hence, the map (\ref{oglic}) is injective. Let us prove that a homomorphism $f:I_v^{\tame}\to G(\overline F)$ factorizes through the surjection $I_v^{\tame}\to\muu_e$.  By \cite[Page 410]{CohomologyNF}, the group~$I_v^{\tame}$ is isomorphic to the group $\prod_{\ell\neq p_v}\ZZ_{\ell}$. In particular, for every~$n$ coprime to~$p_v$, there exists up to isomorphism a unique quotient of order~$n$ of~$I_v^{\tame}$ which is moreover a cyclic group (this fact is equivalent to its more known ``dual" claim, that is, for every~$n$ coprime to~$p_v$ there exists a unique subgroup of $\oplus_{\ell\neq p_v}\QQ_{\ell}/\ZZ_{\ell}\subset\QQ/\ZZ$ of order~$n$ and it is moreover cyclic). %Therefore,  there exists $k\geq 1$ and a homomorphism $f_k:\muu_k\to~G(\oF)$ such that $f_e\circ (I_v^{\tame}\to \muu_k)=~f$. Moreover,  %Before proving the existence, we establish that up to isomorphism, for every positive integer~$n,$ there exists at most one quotient of order~$n$ of~$\muup$, which is, moreover, a cyclic group. As abelian groups, one has that $$\muup= \prod_{\ell\neq p}\ZZ_{\ell}.$$ Using the fact that the Pontryagin dual of a product of compact groups is the direct sum of their Pontryagin duals \cite[Chapter II, \S 1, \no 9, Corollary 3 of Proposition 11]{TSpectrale} and the fact that for a prime~$\ell,$ the group $\ZZ_{\ell}^*$ identifies with the group $\QQ_{\ell}/\ZZ_{\ell}$ from \cite[Chapter II, \S 1, \no 9, Corollary 4 of Proposition 9]{TSpectrale}, we deduce that $$(\muup)^*=\oplus_{\ell\neq p}(\QQ_{\ell}/\ZZ_{\ell}).$$ In particular, $(\muup)^*$ is a subgroup of $\QQ/\ZZ$ and any of its subgroups of order $n$ if it exists, is the cyclic group $\{\frac{k}{n}|k\in\{0\doots n-1\}\}$. It follows from \cite[Chapter II, \S 1, \no 7, Theorem 4]{TSpectrale} that, up to isomorphism, for every integer $n$, there exists at most one finite quotient of $\muup$ of order $n$ and that is a cyclic group of order $n$.
The homomorphism~$f$ can thus be factorized as $$I_v^{\tame} \to \muu_{\#\Imm(f)}\xrightarrow{\sim}\Imm(f)\hookrightarrow G(\overline F).$$ %If~$m$ is the order of $\Imm(f)$, it follows that $I_v^{\tame}\to \Imm(f)\hookrightarrow G(\overline F)$ factorizes as $I_v^{\tame}\to\muu_m\xrightarrow{\sim}\Imm(f)\hookrightarrow G(\overline F)$. 
Now, as $\#\Imm(f)|e$, the surjection $I_v^{\tame}\to~\muu_{\#\Imm(f)}$ factorizes as $I_v^{\tame}\to\muu_e\to\muu_{\#\Imm(f)}$ and we deduce the factorization of~$f$ $$I_v^{\tame}\to\muu_{e}\to\muu_{\#\Imm(f)}\xrightarrow{\sim}\Imm(f)\hookrightarrow G(\overline F).$$
We deduce that the first map in the statement, and hence the second map in the statement, are bijections. If~$G$ is commutative, then both maps are isomorphisms of abelian groups. 
%We have, thus, proven the existence and let us prove the last claim. %is a cyclic group and let us denote by $m$ its order. The dual homomorphism $\phi^*:G(\overline F)\to (\muup)^*$ factorizes through the inclusion $(\Imm(\phi))^*\hookrightarrow (\muup)^*$. It follows that the subgroup $(\Imm(\phi))^*\subset (\muup)^*$ is cyclic. Moreover, its order is equal to the order of $\Imm(\phi)$ and hence divides~$e$. We deduce that $(\Imm(\phi))^*$ canonically identifies with $$\{\frac{k}{e}|k\in\{0\doots e-1\}\}\cong (\muu_e)^*.$$ The homomorphism $\phi^*:(G(\overline F))^*\to (\muup)^*,$ hence factorizes as $$(G(\overline F))^*\to (\muu_e)^*\to(\muup)^*$$ and we deduce that its dual $\phi:(G(\overline F))\to \muup$ factorizes as $$\muup\to \muu_e\to G(\overline F).$$ The uniqueness follows from the facts that the homomorphisms $\muup\to\muu_e$ and $\phi:\muup\to \Imm(\phi)$ are surjective. The claim is proven.
%Moreover, for every prime $\ell\neq p$, one has that $\QQ_{\ell}/\ZZ_{\ell}$ is a, hence the order of a finite subgroup of $(\muup)^*$ is coprime to $p$. It follows that we have an isomorphism $$\ZZ_p\times \muup\xrightarrow{\sim}\widehat{\ZZ}.$$By \ref{}, one has an isomorphism $\QQ/\ZZ\xrightarrow{\sim}(\widehat{\ZZ})^*$ and \cite[Chapter II, \S 1, \no 7, Corollary 5 of Theorem 4]{TSpectrale} produces an isomorphism $$(\QQ/\ZZ)\xrightarrow{\sim}(\ZZ_p)^*\times \big(\prod_{p\neq \ell}\ZZ_{\ell}\big)^*.$$ and an identification of the Pont
%
%The abelian group $\muup$ identifies with the quotient of $\widehat{\ZZ}=\prod_{\ell}\ZZ_{\ell}$ by its $p$-subgroup $\ZZ_p=\ZZ_p\times\{0\}_{\ell\neq p}$. Applying Pontryagin duality and the fact that , we obtain an exact sequence $$ 0\to (\muup)^*\to(\widehat{\ZZ})^*\to (\ZZ_p)^*\to 0.$$ It follows that
%By using that the canonical map $I_v^{\tame}\to \muu_e$ is $\Gamma_v^{\un}$-equivariant, for $\gamma\in\Gamma^{\un}_v$ and $f\in\Hom(I_v^{\tame}, G(\overline F))$ we get that:
%\begin{align*}
%\gamma\cdot f&= \gamma\circ f\circ \gamma^{-1}\\
%&=\gamma\circ(f_e\circ (I_v^{\tame}\to\muu_e))\circ\gamma^{-1}\\
%&=\gamma\circ f_e\circ ((I_v^{\tame}\to\muu_e)\circ\gamma^{-1})\\
%&=(\gamma\circ f_e\circ\gamma^{-1})\circ(I_v^{\tame}\to \muu_e)\\
%&=(\gamma\cdot f_e)\circ(I_v^{\tame}\to\muu_e).
%\end{align*}
%By the uniqueness of the induced homomorphism $\muu_e\to G(\oF)$, we deduce that $\gamma\cdot f_e= (\gamma\cdot f)_e$. The statement is proven.
\end{proof}
%The isomorphism of groups $\Hom(I_v^{\tame}, G(\oF))\xrightarrow{\sim}\Hom(\muu_e, G(\oF))$ is $G(\oF)$-equivariant, if the group~$G(\oF)$ acts on both groups by conjugation. The $G(\oF)$-action commutes with~$\Gamma_v^{\un}$-action and we deduce a $\Gamma_v^{\un}$-bijection of pointed sets 
%\begin{align*}H^1(I_v^{\tame}, G(\oF))&=\Hom(I_v^{\tame}, G(\oF))/G(\oF)\\&\xrightarrow{\sim}\Hom(\muu_e, G(\oF))/G(\oF)\\&=G_*(\oF).
%\end{align*}
\begin{lem} \label{funddefh} %Suppose $\Gfs$ is unramified at $v$ (in particular $\gcd (e,p)=1,$ where $p$ is the characteristic of $\kappa_v$). 
The sequence of pointed sets $$0\to H^1(\Gamma_v^{\un}, \Gfs)\to H^1(\Gamma_v, \Gfs)\xrightarrow{\psi^G_v} G_*(F_v)$$ is exact, where the map $H^1(\Gamma_v^{\un}, \Gfs)\to H^1(\Gamma_v, \Gfs)$ is the inflation map (hence, by  \cite[Chapter I, Section 5.8 (a)]{Cohomologiegalois} is injective) and the map $$\psi^G_v:H^1(\Gamma_v, \Gfs)\to G_*(F_v)$$ is the restriction map $$H^1(\Gamma_v, \Gfs)=H^1(\Gamma_v^{\tame}, \Gfs)\to H^1(I_v^{\tame}, \Gfs)^{\Gamma_v^{\un}}=G_*(\Fv).$$ %composed with the map \begin{align*}H^1(I_v^{\tame}, \Gfs)^{\Gamma_v^{\un}}&=H^1(I_v^{\tame}, \Gfs)^{\Gamma_v^{\un}}\\&\twoheadrightarrow\Hom_{\cont} (I_v^{\tame}, \Gfs)^{\Gamma_v^{\un}}\\&=\Hom (I_v^{\tame}, G(\overline{F}))^{\Gamma_v^{\un}}\\&\xrightarrow{\sim}\Hom(\muu_e, G(\overline{F}))^{\Gamma_v^{\un}}\\
%&\twoheadrightarrow G_*(\Fv).
%\end{align*}
Suppose that~$G$ is commutative. Then the sequence is an exact sequence of abelian groups and the map $H^1(\Gamma_v, G(\overline F))\to G_*(F_v)$ is surjective.
\end{lem}
\begin{proof}
The sequence is the inflation-restriction exact sequence \cite[Chapter I, \S 5.8 (a)]{Cohomologiegalois}, applied to the inclusion $\Gamma_v^{\un}\subset~\Gamma_v$ and the first claim follows. When~$G$ is commutative, 
%, reads: 
%\begin{align*}0\to H^1(\Gamma_v^{\un}, G(\overline F))\to H^1(\Gamma^{\tame}_v, G(\overline F))\to H^1(I_v^{\tame}, G(\overline F))^{\Gamma^{\un}_v}.
%\end{align*}
%The kernel of the surjection of pointed sets $$H^1(I_v^{\tame}, G(\overline F))^{\Gamma^{\un}_v}=\Hom(\muu_e, G(\overline F))^{\Frob_v}\to G_*(F_v)$$ is trivial. Hence, the first claim follows. Moreover, if $G$ is commutative, then $\Hom(\muu_e, G(\overline F))^{\Gamma_v^{\un}}= G_*(\overline {\Fv})$ and 
by \cite[Chapter I, \S 2.6 (b)]{Cohomologiegalois}, the sequence is an exact sequence of abelian groups and there is an exact sequence $$H^1(\Gamma_v, G(\overline F))=H^1(\Gamma_v^{\tame}, G(\overline F))\to G_*(F_v)\to H^2(\Gamma_v^{\un}, G(\overline F)).$$ By \cite[Page 387]{CohomologyNF}, one has that $H^2(\Gamma_v^{\un}, G(\overline F))=0.$ The second claim follows. 
\end{proof}
\begin{mydef}
The map $\psi^G_v:H^1(\Gamma_v, G(\oF))\to G_*(\oF)$ provided by Lemma \ref{funddefh} will be called the residue map (at the place~$v$).
\end{mydef}
We have the following version of compatibility with closed immersions.
\begin{lem}\label{residueandsubgroup}
Let~$\iota:R\hookrightarrow G$ be a closed subgroup (in particular, $R(\oF)$ is unramified and tame at~$v$). The diagram
% https://q.uiver.app/?q=WzAsNCxbMCwwLCJIXjEoXFxHYW1tYV92LFIoXFxvRikpIl0sWzEsMCwiSF4xKFxcR2FtbWFfdiwgRyhcXG9GKSkiXSxbMCwxLCJSXyooXFxvRikiXSxbMSwxLCJHXyooXFxvRikiXSxbMCwxXSxbMCwyLCJcXHBzaV5SX3YiLDJdLFsyLDMsIlxcaW90YV8qIiwyXSxbMSwzLCJcXHBzaV5HX3YiXV0=
\[\begin{tikzcd}
	{H^1(\Gamma^{\tame}_v,R(\oF))} & {H^1(\Gamma^{\tame}_v, G(\oF))} \\
	{R_*(\oF)} & {G_*(\oF)}
	\arrow[from=1-1, to=1-2]
	\arrow["{\psi^R_v}"', from=1-1, to=2-1]
	\arrow["{\iota_*}"', from=2-1, to=2-2]
	\arrow["{\psi^G_v}", from=1-2, to=2-2]
\end{tikzcd}\]
is commutative.
\end{lem}
\begin{proof}
The diagram 
% https://q.uiver.app/?q=WzAsNCxbMCwxLCJSXyooXFxvdmVybGluZSBGKSJdLFsxLDEsIkdfKihcXG92ZXJsaW5lIEYpIl0sWzEsMCwiR18qKEZfdikiXSxbMCwwLCJSXyooRl92KSJdLFswLDFdLFsyLDFdLFszLDJdLFszLDBdXQ==
% https://q.uiver.app/?q=WzAsNCxbMCwxLCJcXEhvbShJX3Zee1xcdGFtZX0sIFIoXFxvRikpIl0sWzEsMSwiXFxIb20oSV92XntcXHRhbWV9LCBHKFxcb0YpKSJdLFsxLDAsIlpeMShcXEdhbW1hX3Zee1xcdGFtZX0sRyhcXG9GKSkiXSxbMCwwLCJaXjEoXFxHYW1tYV92XntcXHRhbWV9LCBSKFxcb0YpKSJdLFswLDFdLFsyLDFdLFszLDJdLFszLDBdXQ==
\[\begin{tikzcd}
	{Z^1(\Gamma_v^{\tame}, R(\oF))} & {Z^1(\Gamma_v^{\tame},G(\oF))} \\
	{\Hom(I_v^{\tame}, R(\oF))} & {\Hom(I_v^{\tame}, G(\oF))}
	\arrow[from=2-1, to=2-2]
	\arrow[from=1-2, to=2-2]
	\arrow[from=1-1, to=1-2]
	\arrow[from=1-1, to=2-1]
\end{tikzcd}\]
is commutative, hence is such the diagram
% https://q.uiver.app/?q=WzAsNCxbMCwxLCJSXyooIEZfdikiXSxbMSwxLCJHXyooRl92KSJdLFsxLDAsIkheMShcXEdhbW1hX3YsIEcoXFxvdmVybGluZSBGKSkiXSxbMCwwLCJIXjEoXFxHYW1tYV92LFIoXFxvdmVybGluZSBGKSkiXSxbMCwxLCJcXGlvdGFfKiIsMl0sWzIsMSwiXFxwc2leR192Il0sWzMsMl0sWzMsMCwiXFxwc2leUl92IiwyXV0=
\[\begin{tikzcd}
	{H^1(\Gamma^{\tame}_v,R(\overline F))} & {H^1(\Gamma^{\tame}_v, G(\overline F))} \\
	{H^1(I_v^{\tame}, R(\oF))^{\Gamma_v^{\un}}=R_*( F_v)} & {H^1(I_v^{\tame}, G(\oF))^{\Gamma_v^{\un}}=G_*(F_v).}
	\arrow["{}"', from=2-1, to=2-2]
	\arrow["{\psi^G_v}", from=1-2, to=2-2]
	\arrow[from=1-1, to=1-2]
	\arrow["{\psi^R_v}"', from=1-1, to=2-1]
\end{tikzcd}\]
The diagram of the lemma is commutative, because it is the composite of the latter diagram and the commutative diagram
% https://q.uiver.app/?q=WzAsNCxbMCwxLCJSXyooXFxvdmVybGluZSBGKSJdLFsxLDEsIkdfKihcXG92ZXJsaW5lIEYpIl0sWzEsMCwiR18qKEZfdikiXSxbMCwwLCJSXyooRl92KSJdLFswLDFdLFsyLDFdLFszLDJdLFszLDBdXQ==
\[\begin{tikzcd}
	{R_*(F_v)} & {G_*(F_v)} \\
	{R_*(\overline F)} & {G_*(\overline F).}
	\arrow[from=2-1, to=2-2]
	\arrow[from=1-2, to=2-2]
	\arrow[from=1-1, to=1-2]
	\arrow[from=1-1, to=2-1]
\end{tikzcd}\]
\end{proof}
We are going to prove a compatibility property with twists. %We choose a counting function $c:G_*(\oF)\to\RR_{\geq 0}$. 
Let $\sigma\in Z^1(\Gamma_v^{\tame}, G(\oF))$. %For almost all~$v$, one has that $$[\sigma]\in H^1(\Gamma_F, G(\oF)).$$
%We denote by~$\prescript{}{\sigma}G$ the twist of~$G$ given by~$\sigma$, that is, the group scheme defined by the $\Gamma_F$-group~$\prescript{}{\sigma}(G(\oF))$. By \cite[Chapter II, Proposition 35 bis]{Cohomologiegalois}, one has a pointed bijection $$\Lambda_{\sigma}:Z^1(\Gamma_v^{\tame}, G^{\sigma}(\oF))\to Z^1(\Gamma^{\tame}_v, G(\oF)),\hspace{1cm}\beta\mapsto \beta\cdot\sigma,$$ which induces a pointed bijection $$\lambda_{\sigma}:H^1(\Gamma^{\tame}_v, G^{\sigma}(\oF))\to H^1(\Gamma^{\tame}_v, G(\oF)).$$
%Analogously, one has maps $\Lambda_\sigma:Z^1(\Gamma_F, \sG(\oF))\to Z^1(\Gamma_F, \sG(\oF))$ and $\lambda_{\sigma}:H^1(\Gamma_F, G^{\sigma}(\oF))\to H^1(\Gamma_F, G(\oF)),$ which by abuse notation, are denoted with the same letters.
\begin{lem}\label{residueandtwist}
Suppose that $[\sigma]\in H^1(\Gamma_v^{\un}, G(\oF))$. The diagram
\[\begin{tikzcd}
	{H^1(\Gamma^{\tame}_v, \sG(\oF))} && {H^1(\Gamma^{\tame}_v, G(\oF))} \\
	& {G_*(\oF)}
	\arrow["{\lambda_{\sigma}}", from=1-1, to=1-3]
	\arrow["{\psi^{G^{\sigma}}_v}"', from=1-1, to=2-2]
	\arrow["{\psi^G_v}", from=1-3, to=2-2]
\end{tikzcd}\]
is commutative.
\end{lem}
\begin{proof}
The condition $[\sigma]\in H^1(\Gamma_v^{\un}, G(\oF))$ is by Lemma \ref{funddefh} equivalent to the condition that the image of $\sigma|_{I_v^{\tame}}\in \Hom(I_v^{\tame}, G(\oF))^{\Gamma_v^{\un}}$ in $G_*(\Fv)=\Conj(I_v^{\tame}, G(\oF))^{\Gamma_v^{\un}}$ is the trivial element, or in other words, $\sigma|_{I_v^{\tame}}$ is conjugate to the trivial map $I_v^{\tame}\to G(\oF)$. But this means that $\sigma|_{I_v^{\tame}}$ itself is the trivial map. Let now $\widetilde\alpha\in Z^1(\Gamma_v^{\tame}, \sG(\oF))$ be a lift of $\alpha\in H^1(\Gamma_v^{\tame}, \sG(\oF)).$ We have that \begin{align*}\psi_v^G(\lambda_{\sigma}(\alpha))=\psi_v^G([\widetilde\alpha\cdot\sigma])=(\widetilde\alpha\cdot\sigma)|_{I_v^{\tame}}&=\widetilde\alpha|_{I_v^{\tame}}\cdot\sigma|_{I_v^{\tame}}\\&=\widetilde\alpha|_{I_v^{\tame}}\\&=\psi^{\sG}_v(\alpha)
\end{align*}and the claim follows. 
\end{proof}
\subsection{Heights}\label{secdefheight} In this subsection we will define heights and provide examples.
\subsubsection{} Let~$G$ be a non-trivial finite \' etale tame $F$-group scheme.  %which is tame (if the characteristic of~$F$ is positive, the order of~$G(\oF)$ is coprime to the characteristic). 
For every place~$v$ such that~$G(\oF)$ is unramified at~$v$ and such that $G(\oF)$ is tame at~$v$, we let $\psi^G_v:H^1(\Gamma_v, G(\oF))\to G_*(\Fv)\subset G_*(\oF)$ be the residue map. %given by Lemma \ref{funddefh}.
\begin{lem}\label{klasikatreba} Let $x\in H^1(\Gamma_F, G(\oF))$. For~$v\in M_F$, let us denote by~$x_v$ the image of$~x$ in $H^1(\Gamma_v, G(\oF))$. For almost every~$v$ such that~$G(\oF)$ is unramified at~$v$, one has that $$\psi^G_v(x_v)=1\in G_*(\oF).$$
\end{lem}
\begin{proof} By Lemma \ref{funddefh}, we need to verify that for almost all~$v$ such that~$G(\oF)$ is unramified at~$v$, the element~$x_v$ lies in the image of the map $H^1(\Gamma_v^{\un}, G(\oF))\to H^1(\Gamma_v, G(\oF))$. There exists a finite Galois extension $K/F$ such that~$x$ lies in the image of the map $H^1(\Gal(K/F), G(\oF))\to H^1(\Gamma_F, G(\oF)).$ Let~$v$ be a finite place such that~$K$ is unramified at~$v$ and let~$w$ be a place of~$K$ lying over~$v$. One has a commutative square 
% https://q.uiver.app/?q=WzAsNSxbMSwwLCJIXjEoXFxHYWwoSy9GKSwgRyhcXG9GKSkiXSxbMiwwLCJIXjEoXFxHYW1tYV9GLCBHKFxcb0YpKSJdLFsxLDEsIkheMShcXEdhbChLX3cvRl92KSwgRyhcXG9GKSkiXSxbMiwxLCJIXjEoXFxHYW1tYV92LCBHKFxcb0YpKSJdLFswLDFdLFswLDJdLFsyLDNdLFsxLDNdLFswLDFdXQ==
\[\begin{tikzcd}
	& {H^1(\Gal(K/F), G(\oF))} & {H^1(\Gamma_F, G(\oF))} \\
	{} & {H^1(\Gal(K_w/F_v), G(\oF))} & {H^1(\Gamma_v, G(\oF)),}
	\arrow[from=1-2, to=2-2]
	\arrow[from=2-2, to=2-3]
	\arrow[from=1-3, to=2-3]
	\arrow[from=1-2, to=1-3]
\end{tikzcd}\]
hence, the element $x_v$ is in the image of $H^1(\Gal(K_w/F_v), G(\oF))$. But, the map $H^1(\Gal(K_w/F_v), G(\oF))\to H^1(\Gamma_v, G(\oF))$ factorizes through the map $$H^1(\Gal(K_w/F_v), G(\oF))\to H^1(\Gamma_v^{\un}, G(\oF)).$$ The statement follows.
\end{proof}
Let~$\Sigma_G$ be the finite set of all finite places of~$F$ for which~$G$ is ramified at~$v$ or for which $G$ is wild at~$v$ (i.e. not tame). 
\begin{mydef}\label{definitionofheight}Let $c:G_*(\overline F)\to\RR_{\geq 0}$ be a counting function. Let $M_F^{\infty}\cup \Sigma_G\subset\Sigma\subset M_F$ be a finite set of places.  For $v\in\Sigma$, we let $c_v: BG(F_v)\to\RR_{\geq 0}$ be functions and for $v\in M_F-\Sigma$ let us set $$c_v=c\circ \psi^G_v:BG(F_v)\to \RR_{\geq 0}.$$ For $v\in M_F,$ we denote by~$H_v$ the function $$H_v:BG(F_v)\to\RR_{>0}\hspace{1cm} x\mapsto q_v^{c_v(x)}.$$ The function $$H=H((c_v)_v):BG(F)\to ~\RR_{>0}\hspace{1cm}x\mapsto\prod_{\vMF}H_v(x_v),$$(the product is finite by Lemma \ref{klasikatreba}) is called the height function defined by~$(c_v)_v$ (sometimes simply the height).
\end{mydef}
The counting function~$c$ will be said to be the {\it type} of the height~$H$. For a height~$H$, we may write~$\Sigma_H$ for the set~$\Sigma$ from the definition.
\begin{lem}\label{boundba}
Let $c:G_*(\overline F)\to\QQ_{\geq 0}$ be a counting function and let~$H_1$ and~$H_2$ be two heights having~$c$ for its type. There exists $C_1, C_2>0$ such that for every $x\in BG(F)$ one has that $$C_1\leq \frac{H_1(x)}{H_2(x)}\leq C_2.$$
\end{lem}
\begin{proof} There exists a finite set of places~$\Sigma$, such that for all $v\in M_F-\Sigma$, one has that $H_{1v}=H_{2v}$. As the sets~$BG(F_v)$ are finite, it follows that $$\frac{H_1(x)}{H_2(x)}\geq\prod_{v\in \Sigma}\frac{H_{1v}(x_v)}{H_{2v}(x_v)}\geq \min_{x\in\prod_{v\in\Sigma}(BG(F_v))}\bigg\{\frac{\prod_{v\in\Sigma}H_{1v}(x_v)}{\prod_{v\in\Sigma}H_{2v}(x_v)}\bigg\}.$$ Similarly, one establishes that $\frac{H_1}{H_2}$ is bounded above.
\end{proof}
\subsubsection{}\label{pullbackcountfunc} Let~$G$ be a non-trivial finite \' etale tame group scheme over~$F$ and let~$\sigma\in Z^1(\Gamma_F, G(\oF))$. Let $c:G_*(\oF)=(\sG)_*(\oF)\to\RR_{\geq 0}$ be a counting function. %By \cite[Chapter II, Proposition 35 bis]{Cohomologiegalois}, we have a canonical pointed bijection $$Z^1(\Gamma_F,G(\oF))\to Z^1(\Gamma_F, G(\oF))\hspace{1cm}\alpha\mapsto \alpha\cdot\sigma$$which induces a pointed bijection $$\lambda_\sigma:H^1(\Gamma_F, G^\sigma(\oF))\to H^1(\Gamma_F, G(\oF)).$$
Let $\iota:R\hookrightarrow \sG$ be a closed subgroup. By Lemma \ref{funcinj}, we have a pointed $\Gamma_F$-equivariant map $\iota_{*}:R_*(\oF)\to G_*(\oF)$ which is of trivial kernel. The map $$\iota^*c:R_*(\oF)\to\RR_{\geq 0}\hspace{1cm}x\mapsto c(\iota_*(x))$$ is thus a counting function.
\begin{lem}\label{composingheights} Let $\Sigma_G\cup\Sigma_{G^{\sigma}}\subset \Sigma\subset M_F^0$ be finite. For~$v\in M_F-\Sigma$, we choose functions $c_v:BG(F_v)\to\RR_{\geq 0}$ and set~$H_v$ and~$H$ as in  as in Definition \ref{definitionofheight}. The function $$BR(F)\xrightarrow{\iota} B(\sG)(F)\xrightarrow{\lambda_{\sigma}} BG(F)\xrightarrow{H}\RR_{>0}$$is the height $H((c_v\circ\iota)_v)$ (hence, having~$\iota^*c$ for its type).
\end{lem}
\begin{proof}
Let~$v\in M_F^0-(\Sigma_H\cup\Sigma_G\cup\Sigma_{G^{\sigma}})$ be such that $[\sigma]\in H^1(\Gamma_v^{\un}, G(\oF))$. Let~$x\in BR(F)$. By Lemma \ref{residueandtwist} and Lemma \ref{residueandsubgroup}, one has that \begin{align*}
H_v(\lambda_\sigma(\iota(x_v)))&=\exp\big(\log(q_v)\cdot {c(\psi^G_v(\lambda_{\sigma}(\iota(x_v))))}\big)\\&=\exp\big(\log(q_v)\cdot {c(\psi_v^{(\sG)}(\iota(x_v)))}\big)\\&=\exp\big(\log(q_v)\cdot c(\iota_*(\psi^R_v(x_v)))\big).
\end{align*}
The statement follows.
\end{proof}
%\subsection{Examples of heights} We present several examples of heights that one encounters in literature.
\subsubsection{} In this paragraph we provide an example of a height, which appears in the Malle conjecture.
%
%We recall several facts from \cite[\' Expos\' e V, Section 7 and Section 8]{SGAj}. For every field~$K,$ one can construct an anti-equivalence of categories of finite sets on which the absolute Galois group~$\Gamma_K$ acts and the finite \' etale $K$-algebras. Moreover, if a $K$-algebra~$A$ corresponds to a finite set~$X_A$ endowed with $\Gamma_K$-action, and~$L$ is an extension of~$K$, then the $L$-algebra $A\otimes_KL$ corresponds to the set~$X_A$ which is endowed with the action of~$\Gamma_L$ via the inclusion $\Gamma_L\to\Gamma_K$.

Let~$v$ be a finite place of~$F$. Let~$n$ be a positive integer and let~$G$ be a  non-trivial constant tame $\Fv$-group scheme (we may write~$G$ for $G(\oF)$). We suppose that~$G$ is embedded as a transitive subgroup $\iota:G\hookrightarrow \mathfrak S_n$. The embedding~$\iota$ defines an action of~$G$ on the set $\{1\doots n\}$. One has a canonical surjection $$\Hom_{\cont}(\Gamma_v, G)\to H^1(\Gamma_v, G)%=H^1(\Gamma_v^{\tame}, G)
.$$We may use the same notation for the torsor and the underlying algebra. For a $G$-torsor~$x$ over~$F_v$, we let $\widetilde x:\Gamma_v\to G$ be its lift. The composite homomorphism $\iota\circ {\widetilde x}$ defines a $\Gamma_v$-action on the set $\{1\doots n\}$, and hence by Subsection \ref{equivalenceofcat}, up to isomorphism, a degree~$n$ \' etale algebra~$E_x$ over~$F_v$. We set $$\Delta^{\iota}(x):=\Delta(E_x)\subset\Ov$$where~$\Delta$ denotes the discriminant ideal.
 % Let us set $$G_0:=\{g\in G: g\cdot 1=1\},$$where $\mathfrak S_n$ acts on $\{1\doots n\}$ in the usual way. For $x\in BG(F),$ we set $$\Delta^{\iota}(x):=\Delta(x/G_0).$$ 
\begin{lem}
Let $x\in H^1(\Gamma_v, G(\oF))$. One has that $$\ord_v(\Delta^{\iota}(x))=c_{{\iota}}(\psi_v(x)).$$
\end{lem}
\begin{proof}
%We denote by $\Delta^{\iota}_v$ the function $$\Delta^{\iota}_v:BG(F_v)\to\QQ_{\geq 0}\hspace{1cm}x\mapsto v(\Delta^{\iota}(x_v/G_0)).$$
%Let us first prove that $v(\Delta^{\iota}(x))=v(\Delta^{\iota}_v(x_v))$. 
Consider the base change $x_{\Fv^{\un}}=x\otimes _{\Fv}\Fv^{\un}$. The $G$-torsor~$x_{\Fv^{\un}}$ over~$\Fv^{\un}$ is given by the continuous homomorphism~$\widetilde x|_{\Gamma_{\Fv^{\un}}}$. By composing~$\widetilde x|_{\Gamma_{\Fv^{\un}}}$ with the inclusion~$\iota$, we obtain a continuous action of~$\Gamma_{\Fv^{\un}}$ on $\{1\doots n\}$, and the resulting $\Fv^{\un}$-algebra~$E_{x\otimes {\Fv^{\un}}}$ is isomorphic to $E_x\otimes_{\Fv}\Fv^{\un}$.  Moreover, the inclusion of rings $x\to x_{\Fv^{\un}}$ is unramified, hence, one has equality of the ideals of~$\Ov$:
$$\Delta(E_{x\otimes\Fv^{\un}})\cap\Ov=\Delta(E_x)=(\piv^{v(\Delta(E_x))})=(\piv^{v(\Delta^{\iota}(x))}).$$ 
Let $\sigma_1\doots \sigma_k$ be the orbits of~$\Gamma_{\Fv^{\un}}$ for its action on $\{1\doots n\}$. The indices of point stabilizers in~$\Gamma_{\Fv^{\un}}$ are, by point stabilizer lemma, equal to the cardinalities $\# \sigma_1, \doots \#\sigma_k$. A point stabilizer~$H_j$ in the orbit~$\sigma_j$ defines a degree $\#\sigma_j$ totally ramified extension $(\overline{F_v})^{H_j}$ of~$\Fv^{\un}$, hence its discriminant ideal is $(\piv^{\#\sigma_j -1})$. One has that $$(E_{x\otimes{\Fv^{\un}}})\cong\prod_{i=1}^k(\overline{F_v})^{H_k},$$ thus the discriminant ideal of~$E_{x\otimes{\Fv^{\un}}}$ is equal to $$(\piv^{\sum (\#\sigma_k-1)})=(\piv^{n-k}).$$ We deduce that $$v(\Delta^{\iota}(x))=n-k.$$ On the other side, the homomorphism $\widetilde x:\Gamma_v\to G$ factorizes through $\widetilde x':\Gamma_v^{\tame}\to G$ and one has that $\widetilde x'|_{I_v^{\tame}}\circ (\Gamma_{\Fv^{\un}}\to I_v^{\tame})=\widetilde x|_{\Gamma_{\Fv^{\un}}}.$ The image of $x\in H^1(\Gamma_v, G)=H^1(\Gamma_v^{\tame}, G)$ in~$G_*$ under~$\psi_v$ is the class of the homomorphism $x_e:\muu_e\to G$ for which $x_e\circ (I_v^{\tame}\to\muu_e)=\widetilde x'|_{I_v^{\tame}}$. Let~$\xi\in F_v^{\un}$ be a primitive $e$-th root of unity. The value of $c_{{\iota}}$ on the class of~$x_e$ is defined to be the index $$\ind(\iota(x_e(\xi))).$$ The cyclic group $\iota(x_e(\muu_e))=\iota (\widetilde x'(I^{\tame}_v)=\iota(\widetilde x(\Gamma_{\Fv^{\un}})$ has~$k$ orbits, thus the generator $\iota(x_e(\xi))$ is a product of~$k$ cycles. We obtain that $$\ord_v(\Delta^{\iota}(x))=n-k=\ind(\iota(x_e(\xi)))=c_{{\iota}}(\psi_v(x)),$$ as claimed.
\end{proof}
\subsubsection{}\label{discofgtorsor} We now establish that for~$G$ non-trivial finite \' etale tame group scheme over~$F$, one has that the function $$BG(F)\to\RR_{>0},\hspace{1cm}x\mapsto N(\Delta(x))$$ is a height. 

Let~$v$ be a place such that $G(\overline F)$ is unramified and tame at~$v$. Let $x\in H^1(\Gamma_v, G(\overline F))$ and let $\widetilde x:\Gamma_F\to G(\overline F)$ be a continuous crossed homomorphism lifting $x$. By the fact that $G(\overline F)$ is unramified at $v$, we deduce that the restriction $\widetilde x|_{I_v^{\tame}}$ is a homomorphism, and that the restriction $\widetilde x|_{I_v^{\tame}}$ does not depend on the choice of the lift of $\widetilde x$. We let $x_e:\muu_e\to G(\overline F)$ be  the homomorphism given by Lemma \ref{indmue}.
\begin{prop}Let $\vMFz$ be such that $G$ is unramified at $v$. For $x\in H^1(\Gamma_v, G(\oF))$, one has that $$\ord_v(\Delta(x))=c_{\Delta}(\psi_v(x)).$$
\end{prop}
\begin{proof}
Let $\widetilde x:\Gamma_v^{\tame}\to G(\overline F)$ be a continuous crossed homomorphism which lifts $x\in H^1(\Gamma_v, G(\overline F))=H^1(\Gamma_v^{\tame}, G(\overline F)).$ As~$G(\overline F)$ is unramified at~$v$, for $\gamma, \gamma'\in I_v^{\tame}$, one has that $$\widetilde x(\gamma\gamma')=\widetilde x(\gamma)(\gamma\cdot \widetilde x(\gamma'))=\widetilde x(\gamma)\widetilde x(\gamma'),$$i.e. $\widetilde x|_{I_v^{\tame}}$ is a homomorphism. The group scheme $G_{\Fv^{\un}}=G\times _{\Fv}\Fv^{\un}$ is constant, and the $G_{\Fv^{\un}}$-torsor~$x_{\Fv^{\un}}$ is represented by the continuous crossed homomorphism $\widetilde x|_{I_v^{\tame}}:I_v^{\tame}\to G(\overline F)$. We let $M:=\widetilde x(I_v^{\tame}).$ Note that the extension $L:=(\Fv^{\tame})^{\ker(\widetilde x|I^{\tame}_v)}/F_v^{\un}$ is a Galois totally ramified degree $\# M$ extension, hence its discriminant ideal equal to $(\piv^{\#M -1}).$ One has that $x_{\Fv^{\un}}\cong L^{[G(\oF):M]},$ thus, the discriminant ideal of $x_{\Fv^{\un}}$ is equal to $(\piv^{(\#M-1)\cdot [G(\oF):M] }).$ The morphism $x_{\Fv^{\un}}\to x$ is unramified, hence the discriminant ideal of~$x$ is equal to $(\piv^{(\# M-1)\cdot [G(\oF):M]}),$ i.e. $\ord_v(\Delta_v(x))=(\#M-1) \cdot [G(\oF):M].$

On the other side the image of $x\in BG(F_v)=H^1(\Gamma_v^{\tame}, G(\oF))$ in $G_*(\Fv)\subset G_*(\oF)$ is the class of the homomorphism~$x_e$. By definition, the value of~$c_{\Delta}$ at the class of~$x_e$ is equal to $$(\# x_e(\muu_e)-1)\cdot [G(\oF): x_e(\muu_e)]=(\# M-1)\cdot [G(\oF):M].$$ The statement is proven.
%By Lemma \ref{indmue}, 
%The homomorphism $\widetilde x|_{I_v^{\tame}}:I_v^{\tame}\to G(\oF)$ is $\Gamma_v$-equivariant VERIFY, and hence, the homomorphism  $\widetilde x_e:\muu_e\to G(\overline F)$ induced by Lemma \ref{indmue} is $\Gamma_v$-equivariant. The $\Gamma_v$-equivariant homomorphism $\widetilde x_e:\muu_e\to G(\oF)$ factorizes through a unique injective homomorphism $\widetilde x_k:\muu_k\to G(\overline F)$. The homomorphism $\muu_e\to\muu_k $ given by $x\mapsto x^{e/k}$ is $\Gamma_v$-equivariant, thus $\widetilde x_k:\muu_k\to G(\oF)$ is $\Gamma_v$-equivariant. We deduce a monomorphism of $\Fv$-schemes $\widetilde x_k:\mu_k\to G$. The quotient $x/\mu_k$ for the $\mu_k$-action induced by $\widetilde x_k$ is an affine \' etale $\Fv$-scheme $\Spec(R)$. Note that the $\Gamma_v$-module $\Spec(R)(\oF)$ is unramified WHY, hence~$R$ is unramifed. We write~$R$ as a product of finite unramified $F_v$-extensions $R=F_1\times\cdots F_{\ell}$.
%
%Now, $x\to R$ is a $\mu_k$-torsor. The decomposition $\Spec(R)=\coprod_{i=1}^{\ell}\Spec(F_i)$ gives a decomposition of $x$ into a union $x=\coprod_{i=1}^{\ell} x_i$, where each~$x_i$ is a $\mu_k$-torsor over~$F_i$. One has that $$\Delta(x/R)=\prod_{i=1}^{\ell}\Delta(x_i/F_i).$$As each $F_i$ is unramified over~$F_v$, we deduce that 
%\begin{align*}
%\Delta_v(x_i)&=N_{F_i/F_v}(\Delta(x_i/F_i)\cdot\Delta(F_i/F_v)^{[x_i:F_v]}\\
%&=\Delta(x_i/F_i)^{[F_i:F_v]}
%\end{align*}
%Finally, by \cite[Lemma 8.1.2.1]{darda2021rational}, one has that $\Delta(x_i/F_i)=\pi_i$
%
% By using transitivity of the discriminants we deduce that $$\Delta_v(x)=\Delta(x/R)^{[R:\Fv]},$$
%Discriminant of totally ramified extension $F_v(\sqrt[\ell]{\piv})$ is $\ell-1$ (see e.g. \cite[Lemma 8.1.2.1]{darda2021rational}).
\end{proof}
\subsubsection{}  We prove that the number of torsors of bounded height is finite. Let~$G$ be an \' etale group scheme and let $c:G_*(\oF)\to\RR_{\geq 0}$ be a counting function.
\begin{prop}\label{estimatenumber}
%Let $c:G_*(\overline F)\to\QQ_{\geq 0}$ be a counting function and l
Let $H$ be a height having $c$ for its type.  There exists $C,m>0$, such that for every $A>0,$ one has that $$\#\{x\in BG(F)|H(x)\leq A\}\leq CA^m.$$
\end{prop}
\begin{proof}
Let $c_{\Delta}:G_*(\overline F)\to\QQ_{\geq 0}$ be the discriminant counting function. Let $M>\min_{1\neq x\in G_*(\overline F)}\frac{c(x)}{c_\Delta(x)}$ be an integer. The function $M\cdot c$ is a counting function and let us set $$H_M:=H((M\cdot c_v)_v)=H((c_v)_v)^M.$$ By Lemma \ref{boundba}, there exists $C'>0$ such that for every $x\in BG(F)$ one has that  $$C'\cdot H_M(x)>N(\Delta(x)).$$ One obtains for~$A>0$ that
\begin{align*}\#\{x\in BG(F)|H(x)\leq A\}&=\#\{x\in BG(F)|H_M(x)\leq A^M\}\\&\leq \frac{\#\{x\in BG(F)|N(\Delta(x))\leq A^M\}}{C'}.
\end{align*}
It suffices, hence, to prove the statement for the norm of discriminant~$N\circ\Delta$. % $H=N\circ\Delta$. 
For a positive integer $k$, let us denote for $A>0$ by $\mathcal N^k(A)$ the number of extensions $$\mathcal N^k(A):=\#\{L/F|[L:F]=k, N(\Delta(L))<A\}.$$ By the main theorem of \cite{SchmidtNF} for the case when $F$ is a number field and by \cite[Lemma 2.4]{Gmethodant} for the case when $F$ is a function field, for every $k>0$ one has that there exists integer $r(k)>0$ such that $$\mathcal N^k(A)\leq C(k)\cdot A^{r(k)}$$ for some $C(k)>0$. Let $W(n)$ denotes the set of partitions of $n$. For a partition $\mathcal P\in W(n)$ of $n$, let us set  $$d(\mathcal P)[X]=\prod_{k\in\mathcal P}C(k)X^{r(k)}\in \RR[X].$$ To give an $F$-\' etale algebra of degree $n$ up to isomorphism amounts to give a partition $\mathcal P\in W(n)$ and for every $k\in\mathcal P$, an $F$-extension of degree $k$. The discriminant of an $F$-extension is bounded below by $1$, and thus the number of \' etale algebras of degree $n$ up to isomorphism, which are of discriminant at most $A$ (where $A>0$), is bounded above by $$\sum_{\mathcal P\in W(n)}\prod_{k\in\mathcal P}\mathcal N^k(A)\leq\sum_{\mathcal P\in W(n)}d(\mathcal P)(A).$$ In particular, one has that $$\#\{x\in BG(F)|N(\Delta(x))\leq A\}\leq d(\mathcal P)(A).$$The statement follows.
\end{proof}
\begin{rem}
\normalfont
The bounds used in the proof of Proposition~\ref{estimatenumber} are clearly not optimal. In Subsection~\ref{secequid} for~$G$ commutative, we will establish the precise asymptotic behaviour of $\#\{x\in BG(F)|H(x)\leq A\}$ when $A\to\infty$. 
\end{rem}
\subsection{Conjecture}
Let~$G$ be a non-trivial finite \' etale tame group scheme over~$F$ and let $c:G_*(\oF)\to \RR_{\geq 0}$ be a counting function. %We let~$BG$ denotes the quotient stack $(\Spec(F))/G$, where~$G$ acts on $\Spec(F)$ trivially. The stack $BG$ is a separated Deligne-Mumford stack by ???. For an $F$-point~$x$ of~$BG$, let $\underline{\Aut}_F(x)$ be the algebraic group space of~$x$, which is a group scheme by .?? By \cite[\href{https://stacks.math.columbia.edu/tag/06QG}{Lemma 06QG}]{stacks-project}, there exists a canonical isomorphism $B\underline{\Aut}_F(x)\to BG$.
\subsubsection{}\label{parbreakthin} In this paragraph, we will define breaking thin maps. For $\sigma\in Z^1(\Gamma_F, G(\oF))$, let~$\sG$ be the twist of~$G$ by~$\sigma$.  For a closed subgroup $\iota:R\hookrightarrow \sG,$ we have a counting function $\iota^*c:R_*(\oF)\to \RR_{\geq 0}$ defined in Paragraph~\ref{pullbackcountfunc}. We have a canonical map $BR(F)\to B(\sG)(F)$, which when composed with the pointed bijection $\lambda_{\sigma}:B(\sG)(F)\to BG(F)$ gives a canonical map $BR(F)\to BG(F)$.
\begin{lem}\label{finitefibers}
The fibers of the map $BR(F)\to BG(F)$ have cardinality at most~$\# G(\oF)$.
\end{lem}
\begin{proof}
Because $\lambda_{\sigma}:B(\sG)(F)\to BG(F)$ is a bijection, we can suppose that~$\sG=G$. In the commutative diagram
% https://q.uiver.app/?q=WzAsNCxbMCwwLCJaXjEoXFxHYW1tYV9GLCBSKFxcb3ZlcmxpbmUgRikpIl0sWzEsMCwiWl4xKFxcR2FtbWFfRiwgRyhcXG92ZXJsaW5lIEYpKSJdLFsxLDEsIkheMShcXEdhbW1hX0YsIEcoXFxvdmVybGluZSBGKSkiXSxbMCwxLCJIXjEoXFxHYW1tYV9GLCBSKFxcb3ZlcmxpbmUgRikpIl0sWzAsMV0sWzEsMl0sWzAsM10sWzMsMl1d
\[\begin{tikzcd}
	{Z^1(\Gamma_F, R(\overline F))} & {Z^1(\Gamma_F, G(\overline F))} \\
	{BR(F)} & {BG(F),}
	\arrow[from=1-1, to=1-2]
	\arrow[from=1-2, to=2-2]
	\arrow[from=1-1, to=2-1]
	\arrow[from=2-1, to=2-2]
\end{tikzcd}\]
the upper horizontal map is injective, the left vertical map is surjective and the right horizontal map has fibers of size at most~$\# G(\oF))$ (the equivalence class of $\gamma\mapsto f(\gamma)\in Z^1(\Gamma_F, G(\oF))$ is given by maps $\gamma\mapsto g^{-1}f(\gamma)(\gamma\cdot g)$ for $g\in G(\oF)$). We deduce that the lower horizontal map has fibers of size at most~$\# G(\oF)$.
\end{proof}
Lemma \ref{finitefibers} gives that if the number of $x\in BR(F)$ such that $(H\circ (BR(F)\to BG(F))(x)\leq B$ is at least $C_1 B^{r_1}\log(B)^{r_2}$ for $B\gg 0$, where $C_1, r_1, r_2$ are positive constants, then the number of $x\in BG(F)$ such that $H(x)\leq  B$ is at least $(C_1/\# G(F)) B^{r_1}\log(B)^{r_2}$ for $B\gg 0$.
\begin{mydef}
%Let $\sigma\in Z^1(\Gamma_F, G(\oF))$ and let~$G^{\sigma}$ be the twist of~$G$ by~$\sigma$.  Let~$\iota:R\hookrightarrow G^\sigma$ be a closed subgroup of~$G^{\sigma}$.
The canonical map $$BR(F)\to B(\sG)(F)\xrightarrow{\lambda_{\sigma}}BG(F)$$ is said to be breaking thin, if one has that $$(a(\iota^*c), b(\iota^*c))> (a(c), b(c)).$$
in the lexicographic order.  
\end{mydef}
It is immediate that one always has $a(\iota^*c)\leq a(c)$, hence the condition in the definition is equivalent to the condition $a(\iota^*(c))=a(c)$ and $b(\iota^*c)>b(c)$. Let us note that if~$G$ is commutative, then the canonical map $R_*(\oF)\to (\sG)_*(\oF)=G_*(\oF)$ is an injection by Lemma~\ref{funcinj} and so $b(\iota^*c)\leq b(c)$, hence there are no breaking thin maps $BR(F)\to BG(F)$. %$H^1(\Gamma_F, R(\oF))\to H^1(\Gamma_F, G(\oF))$.
%\begin{mydef}
%We say that a $1$-morphism of stacks $f:BH\to BG$ is breaking thin, if the following conditions are satisfied:
%\begin{enumerate}
%\item There exists a point $x\in BG(F)$ such that~$f$  is the composite morphism $BH\to B\underline{\Aut}_F(x)\xrightarrow{\sim} BG,$ where $BH\to B\underline{\Aut}_F(x)$ is the $1$-morphism induced by a closed embedding $H\hookrightarrow \underline{\Aut}_F(x)$ and $B\underline{\Aut}_F(x)\xrightarrow{\sim}BG$ is an isomorphism as above.
%\item One has that $a(f^*c)\geq a(c)$ and that $b(f^*c)\geq b(c)$ with at least one strict inequality.
%\end{enumerate}
%\end{mydef}
%\subsubsection{} We present a weaker variant of Malle conjecture.
%\begin{\label{weakmalleconjecture} Let~$H$ be a height having~$c$ for its type.  For every~$\epsilon>0,$  there exists~$A_1, A_2>0$ such that for every~$B>0$ one has that
%$$A_1 B^{a(c)}\leq \#\{x\in BG(F))| H(x)\leq B\}\leq A_2 B^{a(c)+\epsilon}.$$
%\end{conj}
%RELATE WITH ELLENBERG
%Clearly, if the conjecture (or one of the two inequalities) is verified for one height having~$c$ for its type, then it is verified for any height having~$c$ for its type, and in this case we may say that the pair~$(G,c)$ verifies Conjecture \ref{weakmalleconjecture} (or one of the two inequalities).
%\begin{rem}\label{noweights}
%\normalfont Note that as the cardinalities of~$\Aut(x)$ for $x\in BG(F)$ are bounded by~$\#G(\oF)$, Conjecture \ref{weakmalleconjecture} is equivalent to the version where a point $x\in BG(F)$ is counted with the multiplicity $\#\Aut(x)$.
%\end{rem}
%\begin{rem}
%\normalfont
%For~$G$ constant, one usually counts fields (connected torsors) and one has ``Weak Malle conjecture" \cite{mallefirst}, to which no counterexamples are known. We do not find it intuitive, at this stage, to remove some of the torsors from counting (contrary, to Conjecture~\ref{malletale}, where to have the correct exponent of the logarithmic factor, some torsors are removed). Another difference with \cite{mallefirst} is that we allow counting functions different from the discriminant counting function. Using our language, we can generalize \cite{mallefirst} for any pair~$(G,c)$ in an evident way. We have following implications. If the left inequality in the generalization of \cite{mallefirst} is valid for~$(G,c)$, then, clearly, the left inequality is verified for~$(G,c)$. If for every subgroup~$\iota:R\hookrightarrow G$, one has that the pair $(R,\iota^*c)$ verifies the right inequality of the generalization of Conjecture \cite{mallefirst}, then the pair $(G, c)$ verifies the right inequality of Conjecture \ref{weakmalleconjecture}.
%\end{rem}
%We have following results.
%\begin{lem}
%\begin{enumerate}
%\item Let~$\sigma\in Z^1(\Gamma_F, G(\oF))$. We have that~$(G,c)$ verifies the left (respectively, the right) inequality of Conjecture~\ref{weakmalleconjecture}, then the pair~$(G^{\sigma}, c)$ verifies the left (respectively, the right) inequality of Conjecture \ref{weakmalleconjecture}. 
%\item Suppose that pairs~$(G_1,c_1)$ and~$(G_2, c_2)$ of an \' etale and tame $F$-group scheme verify the left (respectively, the right) inequality of Conjecture \ref{weakmalleconjecture}. Then the pair~$(G_1\times G_2, c_1\oplus c_2sssss)$ verifies the left (respectively, the righ) inequality of Conjecture \ref{weakmalleconjecture}.
%\end{enumerate}
%\end{lem}
%\begin{proof}++++?
%The first claim follows from Lemma \ref{composingheights}. The second claim follows from \ref{}.
%\end{proof}
\subsubsection{}We present a conjecture which predicts the number of $G$-torsors of bounded height. %exponent of $\log(B)$ when~$F$ is a number field (respectively, the order of the pole at~$s=a(c)$, when~$F$ is a function field).
\begin{mydef} %Let $G$ be an \' etale group scheme and let $c:G_*(\oF)\to\QQ_{\geq 0}$ be a counting function. 
An element $x\in BG(F)$ is said to be secure if and only if it is not in the image of a breaking thin map. %Otherwise, it is said to be unsecure.
\end{mydef}
\begin{prop}\label{usefulcriterion}
Let $x\in BG(F)$ and let $\widetilde x\in Z^1(\Gamma_F, G(\oF))$ be a lift of~$x$. The element~$x$ is secure if and only there are no closed subgroups $\iota:R\hookrightarrow \prescript{}{\widetilde x}G$ such that $(a(\iota^*c), b(\iota^*c))> (a(c), b(c)).$ 
\end{prop}
\begin{proof}
If there exists such a subgroup~$R\subset \prescript{}{\widetilde x}G$, then~$x$ is the image of the element~$1\in BR(F)$  for the canonical map $BR(F)\to B(\prescript{}{\widetilde x}G)(F)\xrightarrow{\lambda_{\widetilde x}}BG(F),$ and hence~$x$ is not secure. Suppose now that~$x$ is not secure. Then there exists $\sigma\in Z^1(\Gamma_F, G(\oF)),$  a closed subgroup $\kappa:M\hookrightarrow \prescript{}{\sigma}G$ for which $(a(\kappa^*c),b(\kappa^*c))>(a(c),b(c))$ and an element $\tau\in BM(F)$ such that~$x$ is in the image of~$\tau$ for the usual map. Let~$\widetilde\tau\in Z^1(\Gamma_F, M(\oF))\subset Z^1(\Gamma_F, \prescript{}{\sigma}G(\oF))$ be a lift of~$\tau$. Part~(5) of Lemma~\ref{listproptwist} gives that the canonical map $\prescript{}{\widetilde{\tau}}\kappa:\prescript{}{\widetilde \tau}M\to \prescript{}{\widetilde \tau}(\sG)$ is a closed immersion. By using Parts (1), (3) and (4) of Lemma~\ref{listproptwist}, we obtain that the diagram
\[\begin{tikzcd}
	{B(\prescript{}{\widetilde\tau}M)(F)} & {B(\prescript{}{\widetilde\tau}(\prescript{}{\sigma}G))(F)=B(\prescript{}{\Lambda_{\sigma}(\widetilde\tau)}G)(F)} \\
	{BM(F)} & {B(\sG)(F)} \\
	& {} & {BG(F)}
	\arrow[from=1-1, to=1-2]
	\arrow["{\lambda_{\widetilde{\tau}}}"', from=1-1, to=2-1]
	\arrow[from=2-1, to=2-2]
	\arrow["{\lambda_{\widetilde\tau}}", from=1-2, to=2-2]
	\arrow["{\lambda_{\sigma}}"', from=2-2, to=3-3]
	\arrow["{\lambda_{\Lambda_{\sigma}(\widetilde\tau)}}", from=1-2, to=3-3]
\end{tikzcd}\]
is commutative. As the image of $1\in B(\prescript{}{\widetilde\tau}M)(F)$ for the left vertical map is~$\tau$, hence, one has that~$x$ is the image of~$1\in B(\prescript{}{\widetilde\tau}M)(F)$ for the map given by composing $\lambda_{\Lambda_{\sigma}(\widetilde\tau)}$ with the upper horizontal map. Now, by Lemma \ref{twistsandstar}, the canonical map $(\prescript{}{\widetilde\tau}M)_{*}\to (\prescript{}{\widetilde\tau}(\sG)_{*}$ identifies with the map $M_{*}\to (\sG)_{*}=G_{*}$, and thus $ (a(\prescript{}{\widetilde\tau}\kappa^*c),b(\prescript{}{\widetilde\tau}\kappa^*c))>(a(c),b(c)).$ The image of $1\in B(\prescript{}{\Lambda_{\sigma}(\widetilde\tau)}G)(F)$ for the map $\lambda_{\Lambda_{\sigma}(\widetilde\tau)}$ is~$x$, it follows that the cocycles $\Lambda_{\sigma}(\widetilde\tau)$ and~$\widetilde x$ are cohomologous. We let~$k\in G(\oF)$ be such that $(\Lambda_{\sigma}(\widetilde\tau))(\gamma)=(k^{-1})(\widetilde x(k))(\gamma(k))$ for every~$\gamma\in \Gamma_F$. We let~$\iota:R\hookrightarrow \prescript{}{\widetilde x}G$ be the closed subscheme corresponding to the $\Gamma_F$-group which is the image of~$\prescript{}{\widetilde\tau}M(\oF)$ for the map $\prescript{}{\Lambda_{\sigma}(\widetilde\tau)}G(\oF)\to\prescript{}{\widetilde x}G(\oF)$ given by $g\mapsto kgk^{-1}.$ It follows from Lemma \ref{twistsandstar} and from above that $$(a(\iota^*c),b(\iota^*c))=(a(\prescript{}{\widetilde\tau}\kappa^*c),b(\prescript{}{\widetilde\tau}\kappa^*c))>(a(c),b(c)).$$The statement is proven. 
\end{proof}
\begin{conj}\label{malletale} %Let $G$ be an \' etale group scheme and let $c:G_*(\oF)\to\QQ_{\geq 0}$ be a counting function. 
Let $H$ be a height on $BG(F)$ having~$c$ for its type.
\begin{enumerate}
\item ($F$ is a number field) There exists $A>0$ such that 
\begin{align*}
\sum_{\substack{x\in BG(F)\text{ is secure}\\H(x)\leq B}}\frac{1}{\#\Aut(x)}\sim_{B\to \infty}A\cdot B^{a(c)}\log(B)^{b(c)-1}.
\end{align*}
\item ($F$ is a global field) One has that $$\#\{x\in BG(F)|\text{$x$ is secure, }H(x)\leq B\}\asymp_{B\to \infty} B^{a(c)}\log(B)^{b(c)}.$$
\end{enumerate}
\end{conj}
\begin{rem}
\normalfont
Note that as the cardinalities of the automorphism groups~$\Aut(x)$ for $x\in BG(F)$ are bounded by~$\#G(\oF)$, Conjecture \ref{malletale} for the function field case is equivalent to the version where a point $x\in BG(F)$ is counted with the multiplicity $\#\Aut(x)$.
\end{rem}
\begin{rem}
\normalfont  If $\sigma\in Z^1(\Gamma_F, G(\oF))$ and $y\in B(\sG)(F)$ is secure, then $\lambda_{\sigma}(y)\in BG(F)$ is secure.  Indeed, let $\widetilde {y}\in Z^1(\Gamma_F, G(\oF))$ be a lift of~$y$. Then $\prescript{}{\widetilde y}(\sG)=\prescript{}{\Lambda_{\sigma}(y)}G$ and the claim follows from Proposition \ref{usefulcriterion}. It follows that Conjecture \ref{malletale} is valid for the finite group scheme~$G$ with the respect to the height $H$ if and only if it is valid for the finite group scheme~$\sG$ with the respect to the height $H\circ \lambda_{\sigma}$. 
\end{rem}
 %\begin{rem}
%\normalfont We remark that $\#\Aut(x)=\#\prescript{}{\widetilde x}G(F),$ where~$\widetilde x\in Z^1(\Gamma_F, G(\oF))$ is a lift of~$x$. We will not use this fact.
%\end{rem}
%\begin{conj}[''Strong Malle conjecture for \' etale group schemes"]
%Let $H$ be a height function having~$c$ for its type. There exists $A>0$ such that 
%\begin{multline*}
%\#\{x\in BG(F)|\textnormal{$x$ is strictly primitive and }H(x)\leq B\}\\\sim_{B\to \infty}A\cdot B^{a(c)}\log(B)^{b(c)-1}.
%\end{multline*}
%\end{conj}
%\begin{lem}
%Let~$G'$ be another \'etale tame $F$-group scheme and let $c':G'_*\to~\RR_{\geq 0}$ be a counting function. Suppose that Conjecture \ref{malletale} is verified for a height~$H$ on $H^1(\Gamma_F, G(\oF))$ having~$c$ for its type and a height~$H'$ on $H^1(\Gamma_F, G'(\oF))$ having~$c'$ for its type. Then the conjecture is verified for the height~$H\otimes H'$.
%\end{lem}
%\begin{proof}
%
%\end{proof}
%\begin{lem}
%Suppose that $G$ is a constant group scheme. One has that if $x\in BG(F)$ is secure, then~$x$ is a field.
%\end{lem}
%\begin{proof}
%Suppose that~$x$ is not a field and let~$\widetilde x\in \Hom(\Gamma_F, G)$ be a lift of~$x$. Then $R:=\widetilde x(\Gamma_F)$ is a proper subgroup of~$G.$  The element~$x$ is in the image of the canonical map $H^1(\Gamma_F, R)\to H^1(\Gamma_F,G)$ of the class of the map $\widetilde x:\Gamma_F\to R$, a contradiction. NOT CONTRADICTION.
%\end{proof}
%\begin{lem}
%Suppose that~$G=\mathfrak S_n$. 
%\end{lem}
%\subsection{$\mathfrak A_4$-torsors}\label{atorsors} As an application of the developed theory, we prove that the number of secure torsors is at least as predicted for the (constant) alternating group~$\Akk$ of order~$12$. The proof uses that Conjecture~\ref{malletale} is valid for~$G$ commutative, which is proven in Section~\ref{sectioncommutativecase}.

%Let~$F$ be a global field, the characteristic of which is not~$2$ or~$3.$  %
%A choice of $6$-th root of~$1$, induces an identification $\Hom(\muu_6,\Akk)\xrightarrow{\sim}\Akk$ and an identification of the pointed set~$(\mathfrak A_4)_*$ with the pointed set of $F$-conjugacy classes in~$\Akk$. The fibers of the canonical map $\Akk\to (\Akk)_*$ are the following sets: the singleton given by the trivial element, the set given by all products of two $2$-cycles in~$\Akk$ and the set given by $3$-cycles in~$\Akk$.

%We let~$H=H_{\Delta}$ be the norm of the discriminant of an $\mathfrak A_4$-torsor over~$F$. It is at almost all places determined by the formula $$(H_{\Delta})_v(x)=\exp\big(\log(q_v)\cdot c(\psi_v^{\mathfrak A_4}(i_v(x)))\big),$$where $i_v:B\mathfrak A_4(F_v)\to B\Akk(F_v)$ is the canonical map and the map $c=c_{\Delta}:(\Akk)_*\to \RR_{\geq 0}$ is the discriminant counting function: $$c([g])=12\cdot \big(1-\frac{1}{\ord(g)}\big),$$ where $\ord(g)$ denotes the order of an element~$g\in\Akk,$ the $F$-conjugacy class of which is~$[g]$. %We now apply the above theory and the result for~$G$ commutative (which will be proven in Section~\ref{sectioncommutativecase}) to give a lower bound for the number of $\mathfrak A_4$-torsors. 
%Let~$V_4$ be the Klein four subgroup of~$\mathfrak A_4$.
%\subsubsection{}In this paragraph, we prove the following result.
%\begin{thm}\label{aktorsor}
%There exists $C>0$ such that $$\#\{x\in B\mathfrak A_4(F)|x\text{ is secure and }H(x)\leq B\}>C\cdot B^{\frac{1}{6}}.$$
%\end{thm}
%\begin{proof}
%The group~$\Akk$ is solvable and by \v Safarevi\v c theorem \cite[Theorem 9.6.1]{CohomologyNF}, there exists a Galois extension~$\sigma'$ of~$F$, having~$\Akk$ for its Galois group. Let $\sigma\in\Hom(\Gamma_F,\Akk)$ be a lift, automatically surjective, of~$\sigma'$. The non-trivial elements of~$V_4$ are forming one conjugacy class and the group~$V_4$ is normal in~$\Akk$. Now, by using the surjectivity of~$\sigma$, we deduce that~$\sak$ has an order~$4$ closed subgroup~$\iota:R\hookrightarrow \sak$, such that the non-trivial elements of~$R(\oF)$ form a single $\Gamma_F$-orbit. For  $y\in (R-\{0\})(\oF),$ we have that the image of its class in~$R_*(\oF)$ for the map $R_*(\oF)\to (\Akk)_*$ is equal to the class of $y\in R(\oF)\subset \Akk$ for the map $\Akk\to(\Akk)_*$. In particular $c(y)=2,$ and hence $a(\iota^*c)=2$.
%We have that $$R_*=\Hom(\ZZ/2\ZZ, V_4)=R.$$
%Hence, $(R_*-\{0\})(\oF)$ is a single $\Gamma_F$-orbit and thus~$b(\iota^*c)=1$. It follows from Theorem~\ref{thmofequidistribution} and Theorem~\ref{satfunf} that for some $C'>0$ one has that $$\#\big\{x\in BR(F)\big|\big(H\circ (BR(F)\to B\Akk(F))\big)(x)\leq B \big\}>CB^{\frac{1}{6}}$$
%Now, Lemma~\ref{finitefibers} gives that $$\#\{x\in \Imm(BR(F)\to B\Akk (F))|H(x)\leq B\}>\frac{C'}{12}B^{\frac{1}{6}}.$$
%To finish the proof, we prove that the elements in the image of the map $BR(F)\to B\Akk(F)$ are secure. Let $r\in BR(F)$ and let $\widetilde r\in Z^1(\Gamma_F, R(\oF))$ be a lift of~$r$. We denote by~$r'$ and~$\widetilde r'$ the images of~$r$ and~$\widetilde r$ in~$BR(F)$ and $Z^1(\Gamma_F, \sak(\oF))$, respectively (so that the class of~$\widetilde r'$ in $B\sak (F)$ is~$r$). We verify that $\lambda_{\sigma}(r')$ is secure. By Proposition \ref{usefulcriterion}, it suffices to verify that~$\prescript{}{\Lambda_{\sigma}(\widetilde r')}(\Akk)$ contains no subgroup~$\kappa:T\hookrightarrow \prescript{}{\Lambda_{\sigma}(\widetilde r')}(\Akk)=\prescript{}{\widetilde r'}(\sak)$, for which $$(a(\kappa^*c), b(\kappa_*c))>(a(c), b(c)).$$ 
%For such a subgroup~$T$, the fact that $a(\kappa^*c)\geq a(c)=2$ implies that $T(\oF)=V_4$. Note that as~$R$ is commutative, we have that $\prescript{}{\widetilde r}R=R$ and we obtain that the closed immersion $R=\prescript{}{\widetilde r}R\hookrightarrow\prescript{}{\widetilde r'}(\sak)$ identifies with $\kappa:T\hookrightarrow \prescript{}{\widetilde r'}(\sak)$. By Lemma \ref{twistsandstar}, the canonical map $(\prescript{}{\widetilde r}R)_{*}\to (\prescript{}{\widetilde r'}(\sak))_{*}$ identifies with the map $\iota_{*}: R_{*}\to (\Akk)_*$. We deduce in particular that $b(\kappa^*c)=b(\iota^*c)=1$, a contradiction. The elements in the image of $BR(F)\to B\Akk(F)$ are thus secure. The statement now follows.
%\end{proof}
%\begin{lem}\label{nonconnak}
%There are $O(B^{\frac{1}{8}}\log(B))$ secure non-connected torsors of height at most~$B$.
%\end{lem}
%\begin{proof}
%Let~$x$ be a non-connected $\mathfrak A_4$-torsor and let $\widetilde x\in\Hom(\Gamma_F, \mathfrak A_4)$ be a lift of~$x$. We have that~$\widetilde x(\Gamma_F)$ is a proper subgroup ~$\iota:A\hookrightarrow \Akk$. Note that~$x$ is in the image of $BA(F)\to B\mathfrak A_4(F)$.
%
%Suppose that $A\leq V_4$. We note that because the non-trivial elements of~$V_4$ are of order~$2$, the Galois action on $(V_4)_*$ is trivial. Hence, the inclusion $\iota:V_4\hookrightarrow\Akk$ satisfies that~$a(\iota^*c)=6$ and that~$b(\iota^*c)=3$. The canonical map $BV_4(F)\to B\mathfrak A_4(F)$ is thus breaking thin and~$x$ is not secure.
%
%Suppose that~$A$ is of order~$3$. We have that $a(\iota^*c)=8$ and that $b(\iota^*(c))\leq 2.$ It follows from Theorem \ref{thmofequidistribution} and Theorem \ref{satfunf} that $$\#\big\{y\in BA(F)\big|\big(H\circ (BA(F)\to B\Akk(F))\big)(y)\leq B\big\}= O(B^{\frac{1}{8}}\log(B)).$$
%It follows that 
%\begin{multline*}\#\{z\in B\Akk(F)|z\text{ is secure and non-connected, }H(z)\leq B\}\\=O(B^{\frac{1}{8}}\log(B)). 
%\end{multline*}
%\end{proof}
%By combining Theorem \ref{aktorsor}, Lemma \ref{nonconnak} and the observation from Paragraph \ref{malleortorsors}, we obtain Theorem~\ref{akdegtw}. 
%\subsubsection{} We generalize the result of \cite{baily} on the number of quartic $\Akk$-extensions over~$\QQ$ to any global field the characteristic of which is not~$2$ or~$3$. We denote by $\iota:\Akk\hookrightarrow \mathfrak S_4$ the canonical inclusion. The resulting counting function~$c_{\iota}$ takes the value~$2$ at every non-trivial element of~$(\Akk)_*$. Let~$H_{\iota}$ be a height having $c_{\iota}$ for its type. The following result follows from Theorem \ref{aktorsor}. %We count $\Akk$-torsors for the height~$H_{\iota}$.
%\begin{cor}
%Let~$F$ be a global field of characteristic not equal to~$2$ or~$3$. For some $C>0$, we have for every~$B>1$ that
%$$\#\{x\in B\Akk(F)| x\text{ is connected, } H_{\iota}(x)<B\}\geq C B^{\frac{1}{2}}.$$
%\end{cor}
%\begin{proof}
 %
%There exists $C'>0$ such that for every $x\in B\Akk(F)$ one has that $$C'\cdot H_{\Delta}(x)\geq H_{\iota}(x)^{3}.$$
%Theorem \ref{aktorsor} gives that there exists~$C>0$, such that
%\begin{align*}
%\#\{x\in B\Akk(F)|x\text{ is connected}, H_{\iota}(x)\leq B\}\hskip-6cm&\\&\geq \#\{x\in B\Akk(F)|x\text{ is connected}, H_{\Delta}(x)\leq (C')^{-1}(B)^{{3}}\}\\  &\geq C B^{\frac{1}{2}}.
%\end{align*}
%\end{proof}
%By observations made in Paragraph~\ref{malleortorsors}, the last corollary implies Theorem~\ref{akdegf}.%\begin{proof}
%Let $A\leq G$ be commutative and normal and let $K\leq G$ be semicommutative such that $\langle A,K\rangle.$ One has a homomorphism $\phi:K\to\Autt(A)$ given by the conjugation by the elements of~$K$. The canonical map $$A\rtimes_{\phi}K\to\langle A, K\rangle=G\hspace{1cm}(a,k)\mapsto ak$$ is surjective and we denote by~$D$ its kernel. The inclusion $A\hookrightarrow G$ is the map $A\to A\rtimes_{\phi} K\to G,$ where the first map is the inclusion $a\mapsto (a,1)$. It follows that $D\cap (A\times\{1\})=\{1\}$. We have a map $A_*\to G_*$ induced by the inclusion
  %\end{proof}
%\begin{proof}
%Let~$\Theta$ be a connected $K$-torsor and let $\sigma_K\in Z^1(\Gamma_F, K(\oF))$ be a lift of~$\Theta$. We denote by~$\sigma\in Z^1(\Gamma_F, (A\rtimes_{\phi}K)(\oF))$ the image of~$\sigma$ for the map $Z^1(\Gamma_F, K(\oF))\to Z^1(\Gamma_F, (A\rtimes_{\phi}K)(\oF))$ induced by the canonical map $$K\to A\rtimes_{\phi}K\hspace{1cm}k\mapsto (1,k).$$
%We have a map $$A\rtimes_{\phi}K\to\langle A, K\rangle=G\hspace{1cm}(a,k)\mapsto ak$$ and we denote by~$\sigma_G$ the image of~$\sigma$ for the induced map $Z^1(\Gamma_F, (A\rtimes_{\phi}K)(\oF))\to Z^1(\Gamma_F, G(\oF))$. Note that the composite of the above maps $$K\to A\rtimes_{\phi}K\to G $$ is the inclusion $K\hookrightarrow G$. Hence, one has that~$\sigma_G$ is precisely the image of~$\sigma_K$ for the map induced by the inclusion. By Lemma \ref{doslemtwis} one has that $\prescript{}{\sigma_G}A=\prescript{}{\sigma}A$. Moreover, the morphism $\prescript{}{\sigma_G}A\to\prescript{}{\sigma_G}G$ is the morphism $\prescript{}{\sigma}A\to\prescript{}{\sigma}(A\rtimes_{\phi}K)\to\sG=\prescript{}{\sigma_G}G$.  
%It follows that the map $$f:B(\prescript{}{\sigma_G}A)(F)\to B(\prescript{}{\sigma_G}G)(F)=B(\sG)(F)\xrightarrow{\lambda_{\sigma}}BG(F)$$
%coincides with the map
%$$B(\prescript{}{\sigma}A)(F)\to B(\prescript{}{\sigma}(A\rtimes_\phi K))(F)\to B(\prescript{}{\sigma}G)(F)\xrightarrow{\lambda_{\sigma}}BG(F).$$
%Part (5) of Lemma \ref{listproptwist} now gives that the maps coincide with the map  $$B(\prescript{}{\sigma}A)(F)\to B(\prescript{}{\sigma}(A\rtimes_\phi K))(F)\xrightarrow{\lambda_{\sigma}}B(A\rtimes_{\phi}K)(F)\to BG(F).$$
%It follows from Lemma~\ref{}, and the fact that $ $, the last map is given by $X\mapsto (X\times_F\Theta)/N$.
%By Proposition~\ref{semidirsat}, one has that $$\#\{x\in B(A\rtimes_{\phi}K)(F)| x\text{ is connected, }H(x)\leq B \}\gg B^{\alpha}\log(B)^{\beta}.$$  The last map is given by $x\mapsto x/N$, hence, it takes connected torsors to connected torsors. The statement follows.
%\end{proof}
%\subsubsection{}We present some examples. Suppose that the characteristic of~$F$ is not~$2$ or~$3$. 
%
%Let~$A$ be a finite \' etale group such that $A(\oF)=\Akk$ is the alternating group of order~$12$. We have a continuous homomorphism $\sigma_{A}:\Gamma_F\to\Aut(\Akk)=\mathfrak \mathfrak S_4$ which defines the $\Gamma_F$-action on~$A(\oF)$. Note that as every automorphism of~$\Akk$ preserves the normal subgroup~$\mathfrak V_4\cong (\ZZ/2\ZZ)\times (\ZZ/2\ZZ)$ of order~$4,$ we deduce that~$A$ has a closed subgroup~$V$ of order~$4$. Let us denote by~$K$ the quotient finite \' etale group scheme~$A/V$.
%
%Note that if~$A$ is semicommutative, then it is isomorphic to the semidirect product $V_4\rtimes_{\phi}K$ for some $\phi:N\to \Autt(V_4).$ It follows that~$\Akk$ has a closed subgroup of order~$3$. Hence, the action of~$\Gamma_F$ on $\{1,2,3,4\}$ induced by~$\sigma_A$ must fix an element. (Indeed without loss of generality suppose that the order~$3$ subgroup is given by $\Gamma_F$-invariant subgroup generated by $(1,2,3)$. Suppose on the contrary that there exists $\gamma\in\Gamma_F$ such that $\sigma_A(\gamma)(4)=i$, with $i\neq 4$. Then $\sigma_A(\gamma) (1,2,3)\sigma_A(\gamma)^{-1}$ is a $3$-cycle fixing~$i$, hence not fixing~$4$, a contradiction with the assumption that the subgroup generated by $(1,2,3)$ is $\Gamma_F$-invariant). In other words, the map~$\sigma_A$ factorizes as $\sigma_A:\Gamma_F\to\mathfrak S_3\subset\mathfrak \mathfrak S_4$ for some of the four subgroups which are symmetric groups of order~$6$. %factorize through a subgroup of order at most~$3$ in~$\mathfrak \mathfrak S_4$.  
%Conversely, suppose there exists a factorization $\sigma:\Gamma_F\to\mathfrak S_3\subset\mathfrak \mathfrak S_4$.  
%Then~$A$ has a closed subgroup~$Z$ of order~$3$. Moreover, the restriction of the quotient map $q:A\to K$ to~$Z$ is an isomorphism. We define a homomorphism $$\phi:K\to \Autt(V)\hspace{1cm}t\mapsto (n\mapsto q^{-1}(t)(n) q^{-1}(t^{-1})).$$ One has a homomorphism $$V\rtimes_{\phi}K\xrightarrow{(n,h)\mapsto n\cdot q^{-1}(h)}\mathfrak A_4,$$ which is easily verified to be injective and surjective, hence is an isomorphism. Hence,~$A$ is semicommutative.
\section{Proof for~$G$ commutative}\label{sectioncommutativecase}
Let~$F$ be a global field and let~$G$ be a non-trivial commutative finite \' etale tame $F$-group scheme. Let~$e$ be the exponent of~$G(\oF)$. The goal of this section is to prove Conjecture \ref{malletale} on the number of $G$-torsors of bounded height.
\subsection{Adelic space of $BG$}
We recall useful properties of the ``adelic space" of the stack~$BG,$ that is of the restricted product of local cohomologies with the respect to the unramified cohomologies. 
\subsubsection{}For every $v\in M_F$, we endow the finite groups $BG(F_v):=H^1(\Gamma_v, G(\oF))$ with the discrete topology. Let~$\Sigma_G$ be the set of finite places of~$F$ such that~$G(\oF)$ is unramified at~$v$. For every~$v\in M_F^0-\Sigma_G$, we will write~$BG(\Ov)$ for the subgroup $H^1(\Gamma_v^{\un}, G(\oF))$ (a more precise notation would be $B\mathcal G(\Ov)$, where~$\mathcal G$ is an \' etale model of~$G$ over~$\Ov$, but we think that our choice will not make confusions).  % is finite (\cite[Chapter I, Corollary 2.3 and Corollary 2.13]{ArithmeticDuality}). We endow the finite groups $BG(F_v)$ with the discrete topology. The generic fiber $BG\to B\mathcal G$ is separated (because it is the base change of separated morphism $\Spec(F)\to S$), thus by the valuative criterion of separatedness \cite[\href{https://stacks.math.columbia.edu/tag/0CLU}{Lemma 0CLU}]{stacks-project}, for $\vMFz-S$, we have that the canonical homomorphism $B\mathcal G(\Ov)\to BG(F_v)$ is injective. We identify $B\mathcal G(\Ov)$ with its image in $BG(F_v)$.
We set $$BG(\AAA):=\sideset{}{'}\prod_{v\in M_F}BG(F_v)[BG(\Ov)] ,$$ where the restricted product is with the respect to the open and compact subgroups $BG(\Ov)\subset BG(F_v)$ for $v\in M_F-\Sigma_G$ (see e.g. \cite[Chapter VIII, \S 6]{CohomologyNF}). 
%As any two commutative finite \' etale models of $G$ are isomorphic over an open of $S$, we deduce that the definition of $BG(\AAA)$ does not depend on the choice of the model $\GG$. Moreover, it follows from identifications with Galois cohomologies CHANGE(\ref{equalcohomf}) and (\ref{unram}), that the space $BG(\AAA)$ identifies with the space $P^1_{M_F}(F,G)$ from \cite[Page 55]{ArithmeticDuality} and, by \cite[Proposition 3.4 (b)]{Cesnaviciuss}, with the space $H^1_{\fppf}(\AAA^{\in M_F}, G_{\AAA})$ from \cite{Cesnaviciuss}.
As for every~$\vMF$, the groups $BG(F_v)$ are Hausdorff and compact, one sees immediately that the topological group~$BG(\AAA)$ is Hausdorff and locally compact. By Lemma~\ref{klasikatreba}, the diagonal map $BG(F)\to\prod_{\vMF}BG(F_v)$ factorizes through the map $i:BG(F)\to BG(\AAA).$
%\begin{prop}\label{icocompact}
%\begin{enumerate}
%\item \cite[Proposition 3.8, Proposition 3.9]{Cesnaviciuss} The abelian group $BG(\AAA)$ is Hausdorff and locally compact.
%\item \cite[Lemma 4.8]{ArithmeticDuality} The image of $i$ lies in $BG(\AAA)$. 
%\item \cite[Proposition 4.12]{Cesnaviciuss} The map $i:BG(F)\to BG(\AAA)$ is of discrete, closed and cocompact image.
%\end{enumerate}
%\end{prop}
For $r=1,2$, let us denote by $\Sh^r(G)$ the kernel of the diagonal map $$H^r(\Gamma_F, G(\oF))\to \prod_{v\in M_F}H^r(\Gamma_v,G(\oF)).$$ The finiteness of $\Sh^r(G)$ is proven in \cite[Chapter I, Theorem 4.10]{ArithmeticDuality}.

The following result is a restatement of \cite[Theorem 8.3.20 (ii)]{CohomologyNF} (it is a stronger result than finiteness of $\Sh^1(G)$).
\begin{lem}\label{finlocun}
Let $\Sigma\supset\Sigma_G\cup M_F^{\infty}$ be a finite set of places of~$F$. The group $$i(BG(F))\cap\bigg(\prod_{v\in \Sigma}BG(F_v)\times\prod_{v\in M_F^0-\Sigma}B\GG(\Ov)\bigg)$$ is finite.
\end{lem}
%When $r=1$, we may write $i$ for $i^r$.
\subsubsection{} In this paragraph, we choose normalization of Haar measures on $i(BG(F))$ and $BG(\AAA)$ and calculate the volume of $BG(\AAA)/i(BG(F))$. The volume, but with different normalizations, has been calculated in \cite[Chapter I, Theorem 5.6]{ArithmeticDuality}.  %Our natural choice of measures gives the Tamagawa number which is compatible with Weil's conjecture on Tamagawa numbers of algebraic groups. 
If $X$ is a discrete space, let us denote by $\coun_X$ the counting measure on~$X$.
\begin{mydef}
For $v\in M_F$, let us denote by $\mu_v$ the Haar measure on $BG(F_v)$ given by $$\mu_v:=\frac{\coun_{BG(F_v)}}{\# G(F_v)}.$$
\end{mydef}
\begin{rem}
\normalfont
In other words, the measure~$\mu_v$ is the counting measure on $BG(F_v)$ weighted by the inverse of the size of the automorphism group of an element in $BG(F_v)$. 
\end{rem}
For every $v\in M_F^0-\Sigma_G$, by using \cite[Lemma 10.14]{Hararicohom}, we have that $$\mu_v(BG(\Ov))=\frac{\# BG(\Ov)}{\# G(F_v)}=\frac{\# H^1(\Gamma_v^{\un}, G)}{\# H^0(\Gamma_v^{\un}, G)}=\frac{\# H^1(\widehat{\ZZ}, G)}{\# H^0(\widehat{\ZZ}, G)}=1.$$
%\begin{lem}\label{gfv}
%\begin{enumerate}
%%\item For every $v\in M_F^0$, one has that $$|H^1(\Gamma_{\kappa_v}, G(\overline{F_v})^{I_v})|=|G(F_v)|.$$
%\item Let $\Sigma'$ be the set given by the union of $\Sigma$ and the finite set of places $v$ for which $G(\overline{F_v})$ is not ramified at $v$ (i.e. for which $G(\overline{F_v})^{I_v}=G(\overline{F_v})$). Let $v\in M_F^0-\Sigma'$. One has that $|B\GG(\Ov)|=|G(F_v)|$ and that the measure $\mu_v$ is normalized by $$\mu_v(B\GG(\Ov))=1.$$
%\end{enumerate}
%\end{lem}
%\begin{proof}
%\begin{enumerate}
%\item Let us denote by $A$ the $\Gamma_{\kappa_v}$-module $G(\overline{F_v})^{I_v}$. We note that\begin{align*}
%|G(F_v)|&=|G(\overline{F_v})^{\Gamma_v}|\\
%&=|(G(\overline{F_v})^{I_v})^{\Gamma_{\kappa_v}}|\\
%&=|H^0(\Gamma_{\kappa_v}, G(\overline{F_v})^{I_v})|\\
%&=|H^0(\Gamma_{\kappa_v}, A)|\\
%&=|A^{\Gamma_{\kappa_v}}|.
%\end{align*}
%Denote by $\Frob_v\in\Gamma_{\kappa_v}$ the Frobenius element, i.e. the element $(x\mapsto x^{q_v})_{x\in\kappa_v^{s}}$. For an element $g\in\Gamma_{\kappa_v}$, there exists a finite extension~$L_g$ such that $g$ acts on $A$ via its image in $\Gal(L_g/\kappa_v)$. In particular, $g$ acts on $A$ via a finite power of $\Frob_v$. We deduce that, $A^{\Gamma_{\kappa_v}}=A^{\sigma}.$ It follows that one has an exact sequence $$0\to A^{\Gal(\overline{\kappa_v}/\kappa_v)}\to A\xrightarrow{\sigma-\Id_{A}}A\to A/(\sigma -\Id_A)A\to 0.$$ We deduce that
%\begin{align*}
%|A^{\Gamma_{\kappa_v}}|=|A/(\sigma-\Id_A)A|=|H^{1}(\Gamma_{\kappa_v},A)||.
%\end{align*}
%The claim is proven.
%\item For every $v\in M_F^0-\Sigma',$ by \ref{unram} and by (1) that \begin{align*}|B\GG(\Ov)|&=|H^1_{\un}(\Gamma_v,G(\overline {F_v}))|\\
%&=|H^1(\Gamma_{\kappa_v},G(\overline{F_v})|)\\&=|G(F_v)|.
%\end{align*} Now, it follows that$$\mu_v(B\GG(\Ov))=\frac{|B\GG(\Ov)|}{|G(F_v)|}=1.$$ The claim is proven.
%\end{enumerate}
%\end{proof}
\begin{mydef}
Let us denote by~$\mu$ the unique Haar measure on the locally compact group $BG(\AAA)$ such that the following condition is satisfied: for every finite set of places~$\Sigma'$ containing the set~$\Sigma_G$ and the set of infinite places~$M_F^\infty$, one has for every choice of open subsets  $U_v\subset BG(F_v)$ for $v\in \Sigma'$, that $$\mu\bigg(\prod_{v\in \Sigma '}U_v\times\prod_{v\in M_F-\Sigma'}B\GG(\Ov)\bigg)=\prod_{v\in\Sigma'}\mu_v(U_v).$$ The existence of a such measure is guaranteed by \cite[Proposition 5.5]{Ramakrishnan}. %Let $\mu$ be the Haar measure on the locally compact abelian group $BG(\AAA)$ given by  applied $$\mu:=\bigotimes_{\vMF}\mu_v.$$
\end{mydef}
%By \cite[Proposition 4.4.3.1]{darda2021rational} that $\mu$ is a Haar measure on $BG(\AAA)$.
Endow $i(BG(F))$ with the measure $\frac{\#\Sh^1(G)}{\#G(F)}\cdot \coun_{i(BG(F))}$.  %We let $$\tau_{BG}:=\bigg(\mu/\bigg(\frac{\#\Sh^1(G)\cdot\coun_{i(BG(F))}}{\#G(F)}\bigg)\bigg)\big(BG(\AAA)/(i(BG(F)))\big).$$
This is the pushforward measure of the measure $\frac{\coun_{BG(F)}}{\# G(F)}$, i.e. the counting measure on $BG(F)$, weighted by the inverse of the size of the automorphism group. In the next lemma, we calculate the Tamagawa number of the stack~$BG$ (that is the volume for the quotient measure \cite[Chapter VII, \S 2, \no 7, Proposition 10]{Integrationd} of the quotient group $BG(\AAA)/i(BG(F))$ for certain choices of Haar measures). 
\begin{lem}\label{tamagawabg}
One has that:
\begin{align*}
\tau_{BG}&:=\bigg(\mu/\bigg(\frac{\#\Sh^1(G)\cdot\coun_{i(BG(F))}}{\#G(F)}\bigg)\bigg)\big(BG(\AAA)/(i(BG(F)))\big)\\
&=\frac{\# G^*(F)}{\#\Sh^2(G)}.
\end{align*}
\end{lem}
\begin{proof}
%We will compare the measures from \cite{ArithmeticDuality} with our measures. In \cite[Chapter I, Theorem 5.6]{ArithmeticDuality}, the discrete set $i(BG(F))$ is endowed with the counting measure $\coun_{i(BG(F))}$, while we have endowed it with $\frac{|\Sh^1(G)|\coun_{i(BG(F))}}{|G(F)|}.$ 
%It follows from the inflation restriction exact sequence that for every $v\in M_F^0$ one has that $$H^1(\Gamma_{\kappa_v},G(\overline{F})^{I_v})=H^1_{\un}(\Gamma_v, G(\overline{F_v})).$$
The Haar measure~$\mu'$ from \cite[Page 71]{ArithmeticDuality} on the locally compact abelian group $BG(\AAA)$ is normalized by $$\mu'\bigg(\prod_{v\in M_F^0}H^1(\Gamma_{\kappa_v}, G(\overline{F})^{I_v})\times\prod_{v\in M_F^\infty}\{1\}\bigg)=1.$$
By using that the fact that for every $\vMFz$ one has $$\# H^1(\Gamma_{\kappa_v}, G(F_v)^{I_v})=\# H^0(\Gamma_{\kappa_v}, G(F_v)^{I_v})=\# G(\oF)^{\Gamma_v}=\# G(F_v),$$ we obtain that
\begin{align*}
\mu\bigg(\prod_{v\in M_F^0}H^1(\Gamma_v^{\un}, G(\overline{F}))\times\prod_{v\in M_F^\infty}\{1\}\bigg)=\prod_{\vMFi}\#G(F_v)^{-1},
\end{align*}
and hence that $\mu':\mu=\prod_{\vMFi}\#G(F_v)^{-1}$. %By Lemma \ref{gfv}, for every $v$ in $M_F^0$, one has that $$|H^1_{\un}(\Gamma_v,G(\overline{F}))|=|G(F_v)|.$$ Furthermore, by (\ref{unram}), for every $\vMFz$, one has that $$H^1_{\un}(\Gamma_v, G(\overline{F})^{I_v})=B\GG(\mathcal O_v).$$
%We will compare the Haar measures $\mu$ and $\mu'$ on the group
%\begin{equation*}
%B\GG(\mathcal O_{\Sigma})=\prod_{v\in M_F^0-\Sigma'}B\GG(\Ov)\times\prod_{v\in M_F^\infty\cup\Sigma}\{1\}.
%\end{equation*} The measure $\mu_v$ is normalized by $$\mu(B\GG(\OO_{\Sigma}))=\bigg( \prod_{v\in\Sigma\cup M_F^{\infty}}|G(F_v)|\bigg)^{-1}.$$ The index $$\bigg(\bigg(\prod_{v\in M_F^0}H^1_{\un}(\Gamma_v,G(\overline{F}))\times\prod_{\vMFi}\{1\}\bigg):BG(\mathcal O_{\Sigma})\bigg)$$ is equal to $$\prod_{v\in\Sigma}|H^1_{\un}(\Gamma_v,G(\overline{F}))|=\prod_{v\in\Sigma}|G(F_v)|.$$
%We deduce that 
%\begin{equation*}
%\mu'(B\GG(\OO_{\Sigma}))=\prod_{v\in\Sigma}|G(F_v)|^{-1}.
%\end{equation*}
%It follows that the factor of proportionality $\mu:\mu'$ of Haar measures on $BG(\AAA)$ is equal to $\prod_{v\in M_F^{\infty}}|G(F_v)|^{-1}.$
By \cite[Chapter I, Theorem 5.6]{ArithmeticDuality} and the fact that $\#\Sh^1(G^*)=\#\Sh^2(G)$ from \cite[Chapter I, Theorem 4.20 a)]{ArithmeticDuality}, we have that
\begin{align*}
\tau_{BG}&= \frac{\# G(F)\cdot (\mu'/\coun_{i(BG(F))})\big(BG(\AAA)/(i(BG(F)))\big)}{\big(\prod_{v\in M_F^{\infty}}\# G(F_v)\big)\cdot \#\Sh^1(G)}\\
&=\frac{\# G(F)\cdot \frac{\#\Sh^1(G)\cdot \# G^*(F)}{\#\Sh^1(G^*)\cdot \#G(F)}\cdot\prod_{\vMFi}\# G(F_v)}{\big(\prod_{v\in M_F^{\infty}}\# G(F_v)\big)\cdot \#\Sh^1(G)}\\
&=\frac{\# G^*(F)}{\#\Sh^2(G)}.
\end{align*}
%\begin{multline*}(\mu'/\coun_{i(BG(F))})\bigg(BG(\AAA)/(i(BG(F)))\bigg)\\=\frac{|\Sh^1(G)|\cdot |G^*(F)|}{|\Sh^2(G)|\cdot |G(F)|}\prod_{v\in M_F^{\infty}}|G(F_v)|.\end{multline*}
%We deduce that
%\begin{align*}
%\tau_{BG}&=\bigg(\mu/\bigg(\frac{|\Sh^1(G)|\cdot\coun_{i(BG(F))}}{|G(F)|}\bigg)\bigg)\big(BG(\AAA)/(i(BG(F)))\big)\\
%&=\frac{|\Sh^1(G)|\cdot |G^*(F)|}{|\Sh^2(G)|\cdot |G(F)|}\prod_{v\in M_F^{\infty}}|G(F_v)|\cdot \frac{|G(F)|}{|\Sh^1(G)|\cdot\prod_{v\in M_F^{\infty}}|G(F_v)|}\\
%&=\frac{|G^*(F)|}{|\Sh^2(G)|},
%\end{align*} 
%as claimed.
\end{proof}
%\begin{rem}
%\normalfont
%The value of the Tamagawa number is consistent with the generalized Weil conjecture on the Tamagawa number algebraic groups. Recall that it predicts that the Tamagawa number of the algebraic group $W$ is given by $\frac{\#\Pic(W)}{\#\Sh^1(F,W)}$. In our situation, the order of the group $G^*(F)$ is precisely the cardinality of the Picard group of the stack $BG$, while $\Sh^2$ appears because of the cohomology shift. 
%\end{rem}
\subsection{Characters}Let $G^*$ be the Cartier dual of~$G$. That is, it is the \' etale $F$-group scheme corresponding to~$\Gamma_F$-module $\Hom(G(\oF), \muu_e),$ where the action is given by $\gamma\cdot f=\gamma\circ f\circ\gamma^{-1}$. 

\subsubsection{}Let~$A$ be a $\Gamma_v$-module. For every~$i\geq 1$, by \cite[Definition 1.4.3]{brauergrothendieck}, one can define residue map: $$r:H^i(\Gamma_v, A)\to H^{i-1}(\Gamma_v^{\un}, A_*),$$  and its kernel is $H^i(\Gamma_v^{\un}, A)$. Note that~$\psi_v^{A}$ is nothing other than $r:H^1(\Gamma_v, A)\to H^0(\Gamma_v^{\un}, A_*)$. For~$\vMF$, we set $$ \inv_v:\Br(F_v)=H^2(\Gamma_v, \overline{\Fv}^{\times})\to \QQ/\ZZ$$ for the invariant map at the place~$v$. For finite~$v$ such that $v(e)=0$, by \cite[Theorem 1.4.14]{brauergrothendieck}, the map $$H^2(\Gamma_v,\muu_e)\to H^2(\Gamma_v^{\un}, \ZZ/e\ZZ)=\bigg(\frac{1}{e}\ZZ\bigg)/\ZZ\subset\QQ/\ZZ\hspace{1cm}x\mapsto -r(x)$$ identifies with the restriction $\inv_v|H^2(\Gamma_v,\muu_e)\to\QQ/\ZZ$. %For~$v$ real, one has that $H^2(\Gamma_v, A)=\ZZ/2\ZZ$ and we set~$\inv_v$ to be the inclusion $\ZZ/2\ZZ\hookrightarrow\QQ/\ZZ.$ For complex~$v$, we set $\inv_v:=0$.   % From now on, we suppose that $\Sigma$ is big enough, so that $G^*$ admits a finite model $\GG^*$ over $S-\Sigma$.
\subsubsection{}For $v\in M_F$, we have Tate pairing 
\begin{equation*}
H^1(\Gamma_v, G(\oF))\times H^1(\Gamma_v, G^*(\oF))\to H^2(\Gamma_v,\overline{\Fv}^{\times})\xrightarrow{\inv_v}\QQ/\ZZ, \end{equation*}\begin{equation*}(x,\chi)\mapsto \inv_v(x\cup \chi),
\end{equation*}
where $\cup$ stands for the cup product. %The pairing is perfect
%
%which is by Tate-Schatz duality \cite[Chapter I, Corollary 2.3 and Corollary 2.13]{ArithmeticDuality} perfect. 
By \cite[Corollary 2.3]{ArithmeticDuality}, the pairing is perfect, that is, it identifies $H^1(\Fv, G^*)$ canonically with the topological dual $H^1(\Gamma_v, G(\oF))^*$.  Whenever $G$ is unramified at~$v$, the subgroups $H^1(\Gamma_v^{\un}, G(\oF))$ and $H^1(\Gamma_v^{\un}, G^*(\oF))$ are orthogonal complements for the pairing. %Moreover, by \cite[III, Corollary 1.4 and Corollary 7.2]{ArithmeticDuality} if $v\in M_F^0-\Sigma$, the subgroups $B\GG(\Ov)\subset BG(F_v)$ and $B\GG^*(\Ov)\subset BG^*(\Fv)$ are orthogonal complements for this pairing.
%\begin{prop}\cite[Proposition 4.10]{Cesnaviciuss}\label{globalpairing}
%The pairing $$BG(\AAA)\times BG^*(\AAA)\to\QQ/\ZZ,$$ $$((x_v)_v,(\chi_v)_v)\mapsto \sum_{v\in M_F} \chi_v(x_v),$$ is continuous if $\QQ/\ZZ$ is endowed with the discrete topology. Moreover, the orthogonal complement of $i(BG(F))$ is $i(BG^*(F))$, where, by abuse of notation, we also write $i$ for the diagonal map $i:BG^*(F)\to BG^*(\AAA)$.
%\end{prop}
%\begin{lem}
%Let $\chi\in B\GG^*(\Ov)=G^*(\overline F)/((\Id_{G^*(\overline F)}-\Frob_v)G^*(\overline F))$ and let $x\in B\GG(\Fv)$. Let $\widetilde \chi\in G^*(\overline F)$ be a lift of $\chi$. One has that $$(x,\chi)_{\Ta}=(\psi_v^G(x),\widetilde\chi)_{G},$$where $(\cdot,\cdot)_G$ is the canonical pairing $$G_*(\overline F)\times G^*(\overline F)\to\mu_e\hspace{1cm}(\phi,\theta)\mapsto\theta\circ\phi.$$
%\end{lem}
%\begin{proof}
%\end{proof}
\begin{lem}\label{tatepairing}
Let $v\in M_F^0$ be a place such that $G(\overline F)$ is unramified at $v$. Let $x\in H^1(\Gamma_v, G(\overline F))$ and let $\chi_v\in H^1(\Gamma_v^{\un}, G^*(\overline F))\subset H^1(\Gamma_v, G(\oF)).$ Denote by $\widetilde\chi_v:\Gamma_v^{\un}\to G^*(\oF)$ a crossed homomorphism lifting $\chi$. One has that \begin{multline*}\inv_v(x\cup\chi)=-(\widetilde\chi(\Frob_v)\circ\psi^G_v(x))\\(\in\Hom(\muu_e,\muu_e)=\ZZ/e\ZZ=\bigg(\frac{1}{e}\ZZ\bigg)/\ZZ\subset\QQ/\ZZ).$$ %=(\psi^G_v(x), \widetilde\chi)_G.
\end{multline*}
\end{lem}
\begin{proof}
%For an unramified~$\Gamma_v$-commutative module~$A$, we set $$r:H^1(\Gamma_v, A)\to H^1(\Gamma_v^{\un}, A(-1))$$ for its Serre residue. %Let $r:H^2(\Gamma_v, \muu_e)\to H^1(\kappa_v , \ZZ/e\ZZ)$ be the Serre residue. 
By \cite[Proposition 6.6]{cohomologicalinvariants}, one has that $$\inv_v(x\cup\chi)=-r(x\cup \chi)=-(r(x)\cup \chi),$$
where the cup product on the right hand side $$H^0(\Gamma_v^{\un}, G_*(\oF))\times H^1(\Gamma_v^{\un}, G^*(\oF)) \to H^1(\Gamma_v^{\un}, \ZZ/e\ZZ)=\ZZ/e\ZZ$$ is defined by the pairing % $r(x)\cup\chi\in H^1(\Gamma_v^{\un}, \ZZ/e\ZZ)\subset H^1(\Gamma_v, \ZZ/e\ZZ)$ is defined by the pairing 
$$G_*(\oF)\times G^*(\oF)\to \Hom(\muu_e,\muu_e)=\ZZ/e\ZZ\hspace{1cm}(f,g)\mapsto g\circ f.$$ We have that $r(x)=\psi^G_v(x)\in G_*(\Fv)\subset G_*(\oF)$. For the identification $$Z^1(\Gamma_v^{\un}, G^*(\oF))\xrightarrow{\sim} G^*(\oF)\hspace{1cm} y\mapsto y(\Frob_v),$$ by \cite[Chapter XIII, Proposition 1]{corpslocaux}, one has that~$\chi\in H^1(\Gamma_v^{\un}, G^*(\oF))$ is the image of~$\widetilde\chi (\Frob_v)$ under the quotient map. We deduce that 
%\[\begin{tikzcd}By \cite[Chapter XIII, Proposition 1]{corpslocaux}, the group $Z^1(\Gamma_v^{\un}, G^*(\oF)$ identifies with a subgroup of~$G^*(\oF)$ via the the homomorphism $$Z^1(\Gamma_v^{\un}, G^*(\oF))\to G^*(\oF)\hspace{1cm} y\mapsto y(\Frob_v),$$ so that~$\chi\in H^1(\Gamma_v^{\un}, G^*(\oF))$ is the image of~$\widetilde\chi (\Frob_v)$ under the quotient map.
%	{Z^1(\Gamma_v^{\un}, G^*(\oF))} & {G^*(\oF)} \\
%	{H^1(\Gamma_v^{\un}, G^*(\oF))} & {G^*(\oF)/((\Id_{G^*(\oF)}-\Frob_v)G^*(\oF)),}
%	\arrow["{y\mapsto y(\Frob_v)}", from=1-1, to=1-2]
%	\arrow[from=2-1, to=2-2]
%	\arrow[from=1-1, to=2-1]
%	\arrow[from=1-2, to=2-2]
%\end{tikzcd}\]
%and its horizontal maps are isomorphisms. We obtain that 
$$r(x)\cup\chi=(\widetilde\chi (\Frob_v))\circ\psi^G_v(x).$$ The statement follows.
%
%
%Hence, the class $\chi$ is equal to the class. The statement is proven.
%quotient map $Z^1(\Gamma_v^{\un}, G(\oF))\to H^1(\Gamma_v^{\un}, G(\oF))$ factorizes as $$Z^1(\Gamma_v^{\un}, G(\oF))\xrightarrow{\widetilde\chi\mapsto \widetilde\chi (\Frob_v)} G^*(\oF)\to G^*(\oF)/(\Id_{G^*(\oF)}-\Frob_v)G^*(\oF)=H^1(\Gamma_v^{\un}, G^*(\oF)).$$
 %class~$\chi\in H^1(\Gamma_v^{\un}, G^*(\oF))$ is the image of $\widetilde\chi(\Frob_v)$ for the map.
\end{proof}
Using the injective homomorphism $$(\QQ/\ZZ)\to S^1\hspace{1cm}q+\ZZ\mapsto \exp(2i\pi q),$$ we may identify $\QQ/\ZZ$ with a subgroup of~$S^1$. % which is injective and identifies $\QQ/\ZZ$ with a subgroup of~$S^1$.
For $x\in H^1(\Gamma_v, G(\oF))$ and $\chi\in H^1(\Gamma_v, G^*(\oF))$, we write~$$\chi(x):=-\inv_v(x\cup\chi)\in S^1.$$
%\begin{proof}
%Let $\widetilde x:\Gamma_v\to G(\oF)$ be a crossed homomorphism %(because~$\Gamma_v$ acts via~$\Gamma_v^{\un}$), 
%lifting~$x$. The composite map $\Gamma_v\to \Gamma_v^{\un}\xrightarrow{\widetilde\chi} G^*(\oF)$ is a crossed homomorphism that we denote by~$\widetilde\chi'$.
%
%The image of $(\chi, x)$ in $H^2(\Gamma_v, \muu_n)$ is equal to the class of the $2$-cocycle $$\Gamma_v^2\xrightarrow{\widetilde\chi'\otimes\widetilde x} G(\oF)\otimes G^*(\oF)\xrightarrow{(f\otimes g)\mapsto g(f)}\muu_e,\hspace{1cm}(a,b)\mapsto \widetilde\chi'(b)(\widetilde x(a)).$$
%The image of $(\chi, x)$ in $H^2(\Gamma_v, \Gm)$ is the class of the $2$-cocycle $\Gamma_v^2\to \Fv^{\times}$ given by $(a,b)\mapsto \widetilde\chi'(b)(\widetilde x(a))$.
%
%
%The map $\inv:H^2(\Fv,\muu_e)\to \ZZ/e\ZZ$ maps the class of $\widetilde\chi'(b)(\widetilde x(a))$ to . 
%\end{proof}
%\begin{proof}
%%The Tate pairing $$H^1(F_v, G)\times H^1(F_v, G^*)\to  \ZZ/(|G(\oF)|\ZZ)$$induces a canonical isomorphism $H^1(F_v, G)^*\xrightarrow{\sim} H^1(F_v, G^*)$. One hence has a commutative diagram   
%\[\begin{tikzcd}
%	{H^1(F_v,G)} & {G_*(F_v)} & 1 \\
%	{H^1(F_v,G^*)^*} & {H^1_{\un}(F_v, G^*)^*} & 1
%	\arrow["{=}", from=1-1, to=2-1]
%	\arrow[from=2-1, to=2-2]
%	\arrow[from=1-1, to=1-2]
%	\arrow[from=1-2, to=1-3]
%	\arrow[from=2-2, to=2-3]
%	\arrow["{=}", from=1-2, to=2-2]
%\end{tikzcd}\]
%where the map $H^1(F_v,G)\xrightarrow{\sim} H^1(F_v, G^*)^*$ is the canonical isomorphism. Let $x\in G_*(F_v)=\Hom(\mu_e, G)^{\Frob_v}$ and let $\widetilde x\in H^1(F_v,G)$ be a lift of $x$. The image of $\widetilde x\in H^1(F_v,G)=H^1(F_v, G^*)^*$ is the restriction $\widetilde x|_{H^1_{\un}(F_v,G^*)}$. On the other side, one has that \begin{align*}(H^1_{\un}(F_v, G^*))^*&=\big(G^*(\oF)/((\Id_{G^*(\oF)}-\Frob_v)G^*(\oF))\big)^*\\
%&=G_*(\oF)^{\Frob_v}.
%\end{align*}
%UNFINISHED.
%\end{proof}
\subsubsection{}One has a global pairing $$BG(\AAA)\times BG^*(\AAA)\to S^1,\hspace{1cm}((x_v)_v,(\chi_v)_v)\mapsto \prod_{v}\chi_v(x_v).$$ By \cite[Page 56]{ArithmeticDuality}, the pairing is perfect. By \cite[Chapter I, Theorem 4.10]{ArithmeticDuality}, the image of $i:BG(F)\to BG(\AAA)$ is discrete, closed and cocompact. Moreover, the same theorem gives that the subgroups $i(BG(F))\subset BG(\AAA)$ and $i(BG^*(F))\subset BG^*(\AAA)$ are orthogonal complements.
\begin{lem}\label{smallemch}
Let $\Sigma '\supset M_F^\infty\cup \Sigma_G$ be a finite set of places of $F$. For $v\in\Sigma'$, let $T_v\subset BG(F_v)$ be a subgroup. One has that $$\bigg(\prod_{v\in \Sigma'}T_v\times\prod_{v\in M_F^0-\Sigma'}B\GG(\Ov)\bigg)^{\perp}=\prod_{v\in\Sigma'}T_v^{\perp}\times\prod_{M_F^0-\Sigma'}B\GG^*(\Ov).$$
\end{lem}
\begin{proof}
Let $\chi \in \bigg(\prod_{v\in \Sigma'}BG(F_v)\times\prod_{v\in M_F^0-\Sigma'}B\GG(\Ov)\bigg)^{\perp}$. Let $w$ be a place of $F.$ Suppose that $w\in \Sigma'$. For every $x\in T_w,$ one has that $\chi ((x)_w,(1)_{v\neq w})=\chi_w(x)=1$, hence, $\chi_w\in T_w^\perp$. Suppose that $w\in M_F^0-\Sigma'$. For every $x\in B\GG(\OO_w),$ one has that $\chi((x)_w,(1)_{v\neq w})=1$, hence $\chi_w\in B\GG^*(\OO_w)$. The claim follows.
\end{proof}
%\begin{proof}
%Hermite-Minkowski theorem (a proof for $F$ number field can be found in \cite[]{} and for $F$ the function field of a curve over a finite field in \cite[]{}) gives that in a fixed algebraic closure of $F$, there are only finitely many extensions of $F$ which are unramified outside $\Sigma$. The group scheme $G$ is assumed to be \' etale, hence, any $G$-torsor is a product of finite field extensions. We deduce that the set $$\{x\in BG(F)|\forall v\in M_F-\Sigma: (i(x))_v\in B\GG(\Ov)\}$$ is finite. It follows that the set $$i(BG(F))\cap \bigg(\prod_{v\in \Sigma\cup M_F^{\infty}}BG(F_v)\times \prod_{v\in M_F^0-\Sigma}B\GG(\Ov)\bigg)$$ is finite.
%\end{proof}
%\begin{rem}
%\normalfont
%In Lemma $\ref{finlocun}$, one can replace the set $\Sigma$ with a sufficiently large finite set $\Sigma'\subset M_F^0$ containing it, so that the intersection reduces to the singleton $\{1\}$. 
%Now, Proposition \ref{globalpairing}, the fact that the orthogonal group of the group generated by closed subgroups is the intersection of the orthogonals of the closed subgroups \cite[Chapter II, \S 1, \no 7, Corollary 2]{TSpectrale} and the fact that the subgroup generated by a discrete subgroup and a compact subgroup is closed \cite[Chapter III, \S 4, \no 1, Corollary 1 of Proposition 1]{TopologieGj}, it follows that
%$$i(BG(F))\bigg(\prod_{v\in \Sigma'}BG(F_v)\times \prod_{v\in M_F-\Sigma'}BG(\Ov)\bigg)=BG(\AAA).$$
%This statement can be reinterpreted as finiteness of class number (see e.g. \cite{MR2904198}).
%\end{rem}
%Using Lemma \ref{finlocun} and Proposition \ref{globalpairing}, we obtain the following corollary.
%\begin{cor}
%There exists a finite set of places $\Sigma'$ which contains $\Sigma\cup M_F^\infty$, such that $$i(BG(F))\bigg(\prod_{v\in \Sigma'}BG(F_v)\times \prod_{v\in M_F-\Sigma'}BG(\Ov)\bigg)=BG(\AAA).$$
%\end{cor}
%\begin{proof}
%We use duality from \cite{Cesnaviciuss}. It follows from \cite[Proposition 4.10]{Cesnaviciuss} and from the fact that the orthogonal group of the group generated by closed subgroups is the intersection of the orthogonals of the closed subgroups \cite[Chapter II, \S 1, \no 7, Corollary 2]{TSpectrale}, that for finite set of places $\Sigma'\supset\Sigma\cup M_F^\infty$ one has that \begin{align*}\bigg((i(BG(F))\bigg(\prod_{v\in \Sigma\cup M_F^\infty}BG(F_v)\times\prod_{\vMFz-\Sigma}BG(\Ov)\bigg)\bigg)^{\perp}\hskip-8cm&\\&=(i(BG(F)))^{\perp}\bigcap\bigg(\prod_{v\in \Sigma\cup M_F^{\infty}}BG(F_v)\times\prod_{\vMFz-\Sigma}BG(\Ov)\bigg)^{\perp}\\&=i(BG^*(F))\bigcap\bigg(\prod_{v\in \Sigma\cup M_F^\infty}\{1\}\times\prod_{\vMFz-\Sigma}BG^*(\Ov)\bigg).\end{align*}
%It follows from Lemma \ref{finlocun} applied to the group $G^*$ that the last intersection is finite. Let $x_1\doots x_\ell$ be all elements of the intersection. For each $x_i\neq 1$, let us choose a place $v_i$ for which $(x_i)_{v_i}\neq 1$. Now let us set $\Sigma':=\{v_1\doots v_{\ell}\}\cup \Sigma\cup M_F^\infty$. It follows that $$i(BG^*(F))\bigcap\bigg(\prod_{v\in \Sigma'}\{1\}\times\prod_{\vMFz-\Sigma'}BG^*(\Ov)\bigg)=\{1\}.$$Hence, $$i(BG(F))\bigg(\prod_{v\in \Sigma'}BG(F_v)\times \prod_{\vMFz-\Sigma'}BG(\Ov)\bigg)=BG(\AAA),$$as claimed.
%\end{proof}
\subsection{Representation associated to a character} In this subsection, we are going to ``twist" the representation given by a $\Gamma_F$-invariant subset of $G_*(\oF)$ by a character $\chi\in H^1(\Gamma_F, G^*(\oF))$. We are then going to study the analytic behaviour of corresponding $L$-functions. The results will be used in Subsection \ref{Fouriertransformsa} to understand the analytic behaviour of Fourier transforms of heights. 

\subsubsection{}\label{definofreprho} In this paragraph, we define our representations.

Let us denote by $\CC\cdot G_*(\overline F)$ the $\CC$-vector space $$ \CC\cdot G_*(\overline F):=\{(z_{\phi}\cdot \phi)_{\phi}|\phi\in G_*(\overline F), z_{\phi}\in \CC\}.$$ For $\Gamma_F$-invariant subset $W\subset G_*(\oF),$ we write $\CC\cdot W$ for the subspace of $\CC\cdot G_*(\oF)$ generated by the elements of~$W.$  %Denote by $Z^1(F,G^*(\overline F))$ the set of continuous crossed homomorphisms $\Gamma_F\to G^*(\overline F)$ and by $B^1_{\cont}(F, G^*(\oF))$ its subset given by the crossed homomorphisms of the form $$\Gamma_F\to G^*(\oF)\hspace{1cm} \gamma\mapsto (\gamma\cdot \theta )(\theta^{-1}) $$ for some $\theta\in G^*(\oF)$.  %We will define a map $$Z^1(\Gamma_F, G^*(\overline F))\to \Hom(\Gamma_F,\GL(\CC\cdot G^*(\overline F))).$$
For $\eta\in G^*(\oF)$ and $\phi\in G_*(\oF)$ we may write $\eta(\phi)$ for $$\eta(\phi)=\eta\circ\phi\in\Hom(\muu_e,\muu_e)=\ZZ/e\ZZ=\bigg(\frac{1}{e}\ZZ\bigg)/\ZZ\subset\QQ/\ZZ\subset S^1.$$ Let $\widetilde \chi\in Z^1(\Gamma_F, G^*(\overline F)).$ For a $\Gamma_F$-invariant subset $W\subset G^*(\oF)$, we define a map $\rho^W_{\widetilde\chi}: \Gamma_F\to\End(\CC\cdot W)$ by  
\begin{align*}
(\rho^W_{\widetilde\chi}(\gamma))_{\phi, \phi'}=\begin{cases} (\widetilde\chi(\gamma))(\phi)&\quad \text{if } \phi=\gamma\cdot\phi',\\
0&\quad\text{otherwise.}
\end{cases}
\end{align*}
Note that every $\rho^W_{\widetilde\chi}(\gamma)$ is a generalized permutation matrix (i.e. it can be obtained from a permutation matrix by replacing ones by non-zero numbers) having $e$-th roots of unity for its non-zero entries. In particular, one has that $\rho^W_{\widetilde\chi}(\Gamma_F)\subset \GL(\CC\cdot W)$. 
\begin{lem}\label{rhorep} Let $\widetilde\chi\in Z^1(\Gamma_F, G^*(\overline F))$ and let $W\subset \CC\cdot G_*(\oF)$ be a $\Gamma_F$-invariant subset.
\begin{enumerate}
\item  For $\gamma,\gamma'\in\Gamma_F$, one has that $\rho^W_{\widetilde\chi}(\gamma\gamma')=\rho^W_{\widetilde\chi}(\gamma)\rho^W_{\widetilde\chi}(\gamma')$.
\item The homomorphism $\rho^W_{\widetilde\chi}:\Gamma_F\to\GL(\CC\cdot W)$ is continuous. In other words, $\rho^W_{\widetilde\chi}$ is a Galois representation.
\item Let $\theta\in G^*(\oF)$ and let $\gamma\in\Gamma_F$. We denote by $D^{\theta}(\gamma)$ the diagonal matrix such that for every $\phi\in G_*(\oF)$ the entry at the place $(\phi,\phi)$ is $\theta(\phi)$. One has that $$\rho^W_{\widetilde\chi}(\gamma)\circ D^{\theta}(\gamma)=D^{\theta}(\gamma)\circ\rho^W_{(\widetilde\chi)(\gamma\cdot \theta)(\theta^{-1})}(\gamma).$$Thus, the representations $\rho_{\widetilde\chi}$ and $\rho_{(\widetilde\chi)(\gamma\cdot\theta)(\theta^{-1})}$ are isomorphic.
\end{enumerate}
\end{lem}
%\item The map $$Z^1(F, G^*(\overline F))\to \CC\hspace{1 cm}\widetilde\chi\mapsto \tr(\rho^W_{\wchi})$$ is $B^1_{\cont}(F, G^*(\overline F))$-invariant.
%\item There exists a finite Galois extension $L^{\chi}/F$ such that $\rssssho_{\widetilde\chi}$ factorizes through the canonical map $\Gamma_F\to\Gamma_{L^{\chi}/F}$.
\begin{proof}
\begin{enumerate}
\item Let $\gamma, \gamma'\in\Gamma_F$ and let $\phi, \phi'\in G_*(\overline F).$ Suppose that $\phi\neq (\gamma\cdot\gamma')\cdot\phi'$. We have that $$(\rho^W_{\widetilde\chi}(\gamma)\cdot\rho^W_{\widetilde\chi}(\gamma'))_{\phi, \phi'}=\rho^W_{\widetilde\chi}(\gamma)_{\phi',\gamma'\cdot\phi}\cdot \rho^W_{\widetilde\chi}(\gamma')_{\gamma'\cdot\phi,\phi}=0=(\rho^W_\chi(\gamma\gamma'))_{\phi, \phi'}.$$
Suppose that $\phi=\gamma\cdot(\gamma'\cdot\phi')$. By using that $\gamma^{-1}\cdot\phi=\gamma^{-1}\circ\phi\circ\gamma$ and by using that $\Gamma_F$ acts trivially on $\Hom(\mu_e,\mu_e)$, we obtain that
\begin{align*}
(\rho^W_{\wchi}(\gamma\cdot\gamma'))_{\phi, \phi'}&=\big(\widetilde \chi(\gamma\gamma')\big)(\phi)\\
&=\big(\widetilde\chi(\gamma)\cdot\big(\gamma\cdot \widetilde\chi(\gamma')\big)\big)(\phi)\\
&=\big(\widetilde\chi(\gamma)\cdot\big(\gamma\circ(\widetilde\chi(\gamma'))\circ\gamma^{-1}\big)\big)(\phi)\\
&=\big((\widetilde\chi(\gamma))(\phi)\big)\cdot \big(\gamma\circ(\widetilde\chi(\gamma'))\circ\gamma^{-1}\circ\phi\big)\\
&=\big((\widetilde\chi(\gamma))(\phi)\big)\cdot (\gamma\circ\widetilde\chi(\gamma')\circ(\gamma^{-1}\cdot\phi)\circ\gamma^{-1})\\
&=\big((\widetilde\chi(\gamma))(\phi)\big)\cdot (\widetilde\chi(\gamma') (\gamma^{-1}\cdot\phi))\\
&=(\rho^W_{\widetilde\chi}(\gamma))_{\phi,\gamma^{-1}\cdot\phi}\cdot(\rho^W_{\widetilde\chi}(\gamma'))_{\gamma^{-1}\cdot\phi, (\gamma')^{-1}\cdot\gamma^{-1}\cdot \phi}\\
&=\big(\rho^W_{\wchi}(\gamma)\cdot\rho^W_{\wchi}(\gamma')\big)_{\phi, \phi'}.
\end{align*}
%\item Let $\theta\in G^*(\overline F)$ be an element, and consider a boundary map $b_{\theta}:\gamma\mapsto (\gamma\cdot \theta)\theta^{-1}.$ For $\phi\in G_*(\overline F)$, we have that \begin{align*}
%(b_\theta(\gamma))(\phi)&=((\gamma\cdot\theta)\theta^{-1})(\phi)\\
%&=(\gamma\cdot\theta)(\phi)\cdot (\theta^{-1}(\phi)))\\
%&=(\gamma\circ\theta\circ\gamma^{-1})(\phi)\cdot (\theta^{-1})(\phi)\\
%&=(\gamma\circ\theta\circ\gamma^{-1}\circ\phi)\cdot (\theta^{-1})(\phi)\\&=(\gamma\circ\theta\circ (\gamma^{-1}\cdot\phi)\circ\gamma^{-1})\cdot \theta^{-1}(\phi)\\&=\theta (\gamma^{-1}\cdot\phi)\cdot\theta ^{-1}(\phi)\\
%&=\theta((\gamma^{-1}\cdot\phi)\phi^{-1})
%\end{align*}
%\item There exists a finite Galois extension $L^{\chi}/F$ such that $$\chi\in H^1(\Gamma_F, G^*(\overline F))=\varprojlim_{\substack{L/F\text{ finite Galois}}}H^1(\Gamma_{L/F}, G^*(\overline F)) $$ is an element of $ H^1(\Gamma_{L^\chi/F}, G^*(\overline F))$. Let us lift $\chi$ to a crossed homomorphism $\Gamma_{L^\chi/F}\to G^*(\overline F)$ and let $\widetilde \chi^L:\Gamma_F\to G^*(\overline F)$ be the crossed homomorphism which is the composite of the lift and the canonical homomorphism $\Gamma_F\to\Gamma_{L^{\chi}/F}$. For $\gamma,\gamma'\in\Gamma_F$ with the same images in $\Gamma_{L^\chi/F}$, one has that $\widetilde\chi^{L}(\gamma)=\widetilde\chi^{L}(\gamma'),$ hence, one has that  $$\rho_{\widetilde\chi}(\gamma)=\rho_{\widetilde\chi^L}(\gamma)=\rho_{\widetilde\chi^L}(\gamma')=\rho_{\widetilde\chi}(\gamma').$$
\item The kernel of the crossed homomorphism $\widetilde\chi:\Gamma_F\to G^*(\oF)$ is an open normal subgroup of finite index. %, and, hence, in Galois correspondance defines a finite Galois extension $K/F$. 
One has that $\widetilde\chi$ is $\ker (\widetilde\chi)$-invariant. %It follows that $\widetilde\chi$ factorizes through the quotient map $\Gamma_F\to(\Gamma_{F}/\ker(\wchi))$.
 There exists a normal subgroup~$N\subset\Gamma_F$ of finite index such that for $\gamma,\gamma'\in\Gamma_F$ having the same image in $\Gamma_F/N$ one has that $\gamma\cdot \phi=\gamma'\cdot \phi$ for every $\phi\in W$.  Now, it follows from the definition of $\rho^W_{\wchi}$ that for $\gamma,\gamma'\in \Gamma_F$ with the same image in $\Gamma_F/( N\cap\ker (\widetilde\chi))$, one has that $\rho^W_{\wchi}(\gamma)=\rho^W_{\wchi}(\gamma')$, i.e. that $\rho^W_{\wchi}$ factorizes through $\Gamma_F/( N\cap\ker (\widetilde\chi))$. The subgroup $N\cap\ker (\widetilde\chi))\subset\Gamma_F$ is an open normal subgroup of finite index. We deduce that $\rho^W_{\wchi}$ is continuous.
\item Let $\phi,\phi'\in~W$. By using that the canonical action of $\Gamma_F$ on $\Hom(\muu_e,\muu_e)$ is trivial and that $\rho^W_{\wchi}(\gamma)_{\phi, \phi'}=0$ unless $\phi'=\gamma^{-1}\cdot\phi$, we obtain that
\begin{align*}
(D^{\theta}(\gamma)^{-1}\circ \rho^W_{\widetilde\chi}(\gamma)\circ D^{\theta}(\gamma))_{\phi,\phi'}&=(\rho^W_{\widetilde\chi}(\gamma))_{\phi,\phi'}\cdot\frac{\theta(\phi')}{\theta (\phi)}\\
&=(\rho^W_{\widetilde\chi}(\gamma))_{\phi,\phi'}\cdot \frac{\theta\circ\gamma^{-1}\circ\phi\circ\gamma}{\theta(\phi)}\\
&=(\rho^W_{\widetilde\chi}(\gamma))_{\phi,\phi'}\cdot\frac{\gamma\circ\theta\circ\gamma^{-1}\circ\phi}{\theta(\phi)}\\
&=(\rho^W_{\widetilde\chi}(\gamma))_{\phi,\phi'}\cdot\frac{(\gamma\cdot \theta)(\phi)}{\theta(\phi)}\\
&=(\rho^W_{(\widetilde\chi)(\gamma\cdot\theta)(\theta^{-1})}(\gamma))_{\phi,\phi'}.
\end{align*}
The claim follows.
\end{enumerate}
\end{proof}
For $\chi\in H^1(\Gamma_F, G^*(\oF))$, we will denote by~$\rho^W_{\chi}$ the (isomorphism class of) representation~$\rho^W_{\widetilde\chi}$ for a lift~$\widetilde\chi$ of~$\chi$. In the rest of the subsection, the choice of a lift may be implicit. As similar matrices have same characteristic polynomials, the notations such as $\tr(\rho^W_{\chi}(\Frob_v))$ or $\tr(\rho^W_{\wchi}(\Frob_v))$ are unambiguous. 

%
%Clearly, if $v$ is a place such that $G^*(\oF)$ is unramified at~$v$, then for every $\chi\in H^1(F, G(\oF))$, the representation~$\rho^W_{\chi}$ is also unramified at~$v$. By abuse of the notation, we write $\rho^W_{\wchi}:\Gamma_{\kappa_v}\to\GL(\CC\cdot W)$ for the induced representation from the restriction $\rho^W_{\wchi}|_{\Gamma_v}$. Let $\Frob_v\in\Gamma_{\kappa(v)}$ be the Frobenius element. In the next lemma, we calculate the trace  $\rho^W_{\wchi}(\Frob_v)$.
\begin{lem}\label{traceofl} Let $\chi\in H^1(\Gamma_F, G^*(\oF))$.  Let $v\in M_F^0-\Sigma_G$ be such that the character~$\chi$ is unramified at~$v$ (i.e. such that~$\chi_v|_{H^1(\Gamma_v^{\un}, G(\oF))}=1$). One has that %the representation~$\rho_{\chi}$ is unramified at~$v$ and that %The representation $\rho_{\widetilde\chi}$ is unramified at $v$ and induces a representation $\Gamma_v^{\un}\to \GL(\CC\cdot W)$.
$$\tr(\rho^W_{\chi}(\Frob_v))=\frac{1}{\# G(F_v)}\sum_{y\in (\psi_v^G)^{-1}(W^{\Frob_v})}\chi_v(y).$$%where by $\tr(\rho^W_{\chi}(\Frob_v))$ is meant the quantity $\tr(\rho^W_{\chi}(\widetilde{\Frob_v}))$ which does not depend on the choice of a lift $\widetilde{\Frob_v}\in\Gamma_v$ of $\Frob_v\in\Gamma_v^{\un}$.   %where $\chi_v\in H^1(F_v, G^*(\oF))$ is the image of $\chi$ for the canonical morphism $H^1(F, G^*(\oF))\to H^1(F_v, G^*(\oF))$.
\end{lem}
\begin{proof}
For every $\phi\in W$, one has that $$(\rho^W_{\widetilde\chi}(\Frob_v))_{\phi, \phi}=\begin{cases} (\widetilde\chi(\Frob_v))(\phi)&\quad \text{if } \phi=\Frob_v\cdot\phi,\\
0&\quad\text{otherwise.}
\end{cases}$$
Lemma \ref{tatepairing} gives that for every $y\in H^1(F_v, G(\oF))$ such that $\psi^G_v(y)=~\phi$, one has that $\widetilde\chi_v(\Frob_v)(\phi)=\chi_v(y),$ where $\widetilde\chi_v$ is the composite morphism $\Gamma_v\to \Gamma_F\xrightarrow{\widetilde{\chi}}G^*(\oF)$.  We deduce that  \begin{align*}\tr(\rho^W_{\chi}(\Frob_v))=\tr(\rho^W_{\widetilde\chi}(\Frob_v))&=\sum _{\phi\in W^{\Frob_v}}\widetilde\chi(\Frob_v)(\phi)\\
&=\sum _{\phi\in W^{\Frob_v}}\widetilde\chi_v(\Frob_v)(\phi)\\
&=\sum_{y\in(\psi^G_v)^{-1}(\phi)}\frac{1}{\#(\psi_v^G)^{-1}(\phi)}\sum_{\hspace{0.2cm}\phi\in W^{\Frob_v}}\chi_v(y).
\end{align*}
%It follows from the equality 
%\begin{align*}
%\#H^1(\Gamma^{\un}_v, G(\oF))=\# H^1(\widehat{\ZZ}, G(\oF))=\# H^0 (\widehat{\ZZ}, G(\oF))&=\# H^0(\Gamma^{\un}_v, G(\oF))\\
%&=\# G(\oF)^{\Frob_v}\\
%&=\# G(F_v)
%\end{align*} from \cite[Lemma 10.14]{Hararicohom} and 
By Lemma \ref{funddefh} for every $\phi\in W$, one has that \begin{equation*}\#(\psi_v^G)^{-1}(\phi)=\#H^1(\Gamma^{\un}_v, G(\oF))=\# G(F_v).\end{equation*}
The claim follows.
\end{proof}
\subsubsection{}In this paragraph, we will study $L$-functions of the representations~$\rho^W_{\chi}$.  Let $W\subset G_*(\oF)$ be a $\Gamma_F$-invariant subset. For a lift $\wchi:\Gamma_F\to G^*(\oF)$ of $\chi\in H^1(\Gamma_F, G^*(\oF))$ and a place $v\in M_F^0$, we set $$L_v\big(s,\rho_{\chi}^{W}\big)=\det\big(1-q_v^{-s}\rho_{\wchi}^{W}(\Frob_v)|(\CC\cdot W)^{I_v}\big)^{-1}.$$ %
By Part (3) of Lemma \ref{tatepairing}, one has that $L_v(s, \rho^W_{\chi})$ does not depend on the choice of $\widetilde\chi$, as notation suggests.  The factors $L_v(s,\rho_{\chi}^W)$ are the local Artin $L$-factors of a Galois representation~$\rho^W_{\chi},$ and hence define non-vanishing holomorphic functions in the domain $\Re(s)\geq 1$. In the next lemma, we summarize the properties of a global $L$-function. We recall that in the function field case, we write~$q$ for the cardinality of the field of constants of~$F$.
\begin{lem}\label{lfunf}
Let~$\Sigma\subset M_F$ be finite. 
\begin{enumerate}
\item The product $$L^{\Sigma}(s,\rho_{\chi}^W):=\prod_{v\in M_F^0-\Sigma}L_v(s,\rho_{\chi}^{W})$$ converges normally on compacts in the domain $\Re(s)>1$ and the function $s\mapsto L^\Sigma(s,\rho^W_{\chi})$ is holomorphic in this domain. Moreover, $s\mapsto L^\Sigma(s,\rho^W_{\chi})$ extends to a meromorphic function on~$\CC$ with a possible pole at $s=1$ which is of order at most $\#(W/\Gamma_F)$.
\item If $\chi=1$ is the trivial character, then the pole of $s\mapsto L^{\Sigma}(s,\rho_{\chi}^W)$ at $s=1$ exists and its order is precisely $\#(W/\Gamma_F)$. 
\item If~$F$ is a number field, then $s\mapsto L^\Sigma(s,\rho^W_{\chi})$ is holomorphic in the domain $\{\Re(s)\geq 1\}-\{1\}$. If~$F$ is a function field, then $s\mapsto L^{\Sigma}(s, \rho^W_{\chi})$ is holomorhic in the domain $$\{\Re(s)>0\}-\{1+2ki\pi\log(q)^{-1}|k\in\ZZ\}.$$
\end{enumerate}
\end{lem}
\begin{proof}
\begin{enumerate}
\item As $s\mapsto L_v(s,\rho^W_{\chi})$ are non-vanishing and holomorphic in the domain $\Re(s)>0$, we can suppose that $\Sigma=\emptyset$. The representations~$\rho^W_{\chi}$ are defined over a finite rang $\ZZ$-module (e.g. over $\ZZ[\xi_e],$ where~$\xi_e$ is $e$-th primitive root of~$1$). Now, \cite[Proposition 2.49]{bourquifzh} applies, and the claims on convergence and meromorphic extension follow. Moreover, by the same proposition, the order of the pole at~$s=1$ is equal to the dimension of the subspace $(\CC\cdot W)^{\rho_{\chi}^W}$ fixed by~$\rho^W_{\chi}$. Let~$Q$ be a $\Gamma_F$-orbit in~$W$ and let $\phi\in Q$. For a vector $0\neq (w_{\phi})_{\phi}\in\CC\cdot Q$, which is fixed by~$\rho^W_{\chi}$, one has for every $\gamma\in\Gamma_F$ that $w_{\gamma\cdot\phi}=\widetilde\chi(\gamma)(\phi)\cdot w_{\phi}$. It follows that the dimension of the subspace of $\CC\cdot Q$ fixed by~$\rho^W_{\chi}$ is no more than~$1$. It follows that the subspace $(\CC\cdot W)^{\rho^W_{\chi}}$ fixed by the representation $\rho^{W}_{\chi}$ is of the dimension no more than $\# (W/\Gamma_F)$. We deduce that the order of the pole at $s=1$ is no more that $\#(W/\Gamma_F)$. 
\item Let us now suppose that~$\chi$ is the trivial character, then $\wchi':\gamma\mapsto 1\in G^*(\oF)$ is a lift of~$\chi$. In this case, every generalized permutation matrix $\rho^W_{\wchi'}(\gamma)$ fixes the vector subspace $$\{(x_Q)_{\phi\in Q}|Q\in(W/\Gamma_F), x_Q\in\CC\} \subset \CC\cdot W.$$ Hence, the subspace $(\CC\cdot W)^{\rho^W_{\chi}}$ is of the dimension $\# (W/\Gamma_F)$.  The claim on the order of the pole for the case $\chi=1$ follows.
\item In the case~$F$ is a number field, by Brauer theorem \cite{brauertheorem}, one has that $s\mapsto L(s,\rho^W_{\chi})$ is holomorhic in the domain $\{\Re(s)\geq 1\}-\{1\}$. Suppose that~$F$ is a function field. Let us decompose $\rho^W_{\chi}=\oplus_{i=1}^a\rho_i^{b_i}$ into a direct sum of the irreducible representations, where for $i\neq j$, the representations~$\rho_i$~and $\rho_j$ are not isomorphic and $\rho_1$ is the trivial representation. By using \cite[Page 196]{Milne} and \cite[Chapter VII, \S 6, Theorem 4]{basicnt}, we obtain that \begin{align*}
L(s,\rho^W_{\chi})=\prod_{i=1}^aL(s,\rho_i^{b_i})&=L(s,\rho_1)^{b_1}\cdot\prod_{i=2}^aL(s,\rho_i)^{b_i}\\&=\zeta_F(s)^{b_1}\cdot \prod_{i=2}^aL(s,\rho_i)^{b_i}\\
&=\bigg(\frac{P(q^{-s})}{(1-q^{-s})(1-q^{1-s})}\bigg)^{b_1}\prod_{i=2}^{b_i}L(s,\rho_i)^{b_i},
\end{align*}
where $\zeta_F$ is the Dedekind zeta function of~$F$ and~$P$ is a polynomial such that $P(0)=1$. For every $2\leq i\leq a$, one has by \cite[Pages 196-197]{Milne} that $s\mapsto L(s,\rho_i)$ is holomorphic on~$\CC$. It follows that in the domain $\Re(s)>0$, the set of the poles of $s\mapsto L(s,\rho^W_{\chi})$ is contained in the set $\{1+2ki\pi\log(q)^{-1}|k\in\ZZ\}$, as claimed.
\end{enumerate}
 %is fixed by every permutation matrixfor every $\gamma\in\Gamma_F$, one has that $\rho^W_{\wchi'}(\gamma)=\rho^W_{\wchi'}(1),$ hence $$\dim\big((\CC\cdot W)^{\rho^W_{\wchi'}}\big)=\dim\big((\CC\cdot W)^{\rho^W_{\wchi'}(1)}\big)=(W/\Gamma_F).$$
%If for some $\gamma\in\Gamma_F$ and some $\phi\in Q$ one has that $\widetilde \chi(\gamma)(\phi)\neq 1,$ then for $a\in\CC$ the vector $(a)_{\phi\in\QQ}$ is not fixed by $\rho^{Q}_{\widetilde\chi}(\gamma)$. It follows that for every $\Gamma_F$-orbit $Q$ in $G_*(\oF)_1$, one has at most $1$ subspace which is $\rho^Q_{\widetilde\chi}$-invariant, which is moreover of dimension~$1$. 
%\end{enumerate}
\end{proof}
% In the next two lemmas, we study analytic properties of $s\mapsto L^\Sigma(s,\rho^W_{\chi})$. 
%\begin{lem}\label{lfunnf} Suppose that $F$ is a number field. 
%The function $s\mapsto L^\Sigma(s,\rho_{\chi}^W)$ extends to a meromorphic function on $\CC$ which is holomorphic in the domain $\{\Re(s)\geq 1\}-\{1\}$. The order of the possible pole of $s\mapsto L^\Sigma(s,\rho_{\chi}^{W})$ at $s=1$ is no more than $\#(W/\Gamma_F).$ Moreover, if $\chi=1$, then $s\mapsto L^\Sigma(s,\rho_{\chi}^{W})$ has a pole at $s=1$ and it is of order precisely $\# (W/\Gamma_F).$
%The For a lift $\wchi:\Gamma_F\to G^*(\oF)$ of $\chi$, let $\rho_{\wchi}^{W}:\Gamma_F\to\GL(\CC\cdot W) $ the representation given by Lemma \ref{rhorep}. 
%The product $$L^{\Sigma}(s,\rho_{\chi}^{G_*(\oF)_1}):=\prod_{v\in M_F^0-\Sigma}L_v(s,\rho_{\wchi}^{W})$$  converges absolutely and uniformly for $\Re(s)>1$ and the function $\mapsto L^\Sigma(s,\rho_\chi^W)$ is holomorphic in this domain. Moreover, the following claims are true.
%\begin{enumerate}
%\item Suppose $F$ is a number field. The function $s\mapsto L^\Sigma(s,\rho_{\chi}^W)$ extends to a meromorphic function in the domain $\Re(s)\geq 1$ with the only possible pole at $s=1$. The order of the possible pole of $s\mapsto L^\Sigma(s,\rho_{\chi}^{W})$ at $s=1$ is no more than  $\#(W/\Gamma_F).$ Moreover, the order of the pole at $s=1$ is precisely $\# (W/\Gamma_F)$ if $\chi=1$.
%\item Suppose $F$ is a finite extension of $\FF_{q}(t)$. The function $s\mapsto L^\Sigma(s,\rho_{\chi}^W)$ is a rational function in $q^{-s}$ with a possible pole at $s=1$ of order no more than $\# (W/\Gamma_F).$ Moreover, the order of the pole at $s=1$ is precisely $\# (W/\Gamma_F)$ if $\chi=1$.
%\end{enumerate}
%\end{lem}
%\begin{proof}
%There exists a finite Galois extension $F'/F$ such that $$\chi\in H^1(\Gamma_{F}, G^*(\overline F))=\varprojlim_{\substack{K/F\text{ finite Galois}}}H^1(\Gamma_{K/F}, G^*(\overline F)) $$ is an element of $ H^1(\Gamma_{F'/F}, G^*(\overline F))$. Let us lift $\chi$ to a crossed homomorphism $\Gamma_{F'/F}\to G^*(\overline F)$ and let $\widetilde \chi^{'}:\Gamma_F\to G^*(\overline F)$ be the crossed homomorphism which is the composite of the lift and the canonical homomorphism $\Gamma_F\to\Gamma_{F'/F}$. Then $\wchi^{'}$ is a lift of $\chi$ and the representation $\rho_{\wchi^{'}}^W$ factorizes through the canonical morphism $\Gamma_F\to\Gamma_{F'/F}$.
% For $\gamma,\gamma'\in\Gamma_F$ with the same images in $\Gamma_{F'/F}$, one has that $\widetilde\chi^{F'}(\gamma)=\widetilde\chi^{F'}(\gamma'),$ hence, one has that  $$\rho_{\widetilde\chi}(\gamma)=\rho_{\widetilde\chi^{F'}}(\gamma)=\rho_{\widetilde\chi^{F'}}(\gamma')=\rho_{\widetilde\chi}(\gamma').$$
%\cite[Corollary 5]{MurtyArtin} one has that for a non-trivial permutation the $L$-function is SOMETHINGS.
%It follows from Brauer's theorem \cite{brauertheorem} that $s\mapsto L(s,\rho_{\chi}^W)$ is meromorphic on~$\CC$ and that it %\cite[Theorem 7]{zetafunlfun}
%is holomorphic and non-vanishing in the domain $\{\Re(s)\geq 1\}-\{1\}$. 
%\end{proof}
%Let us now treat the function field case. We need the following lemma. Note that for 
%\begin{lem}
% Let $v_1, v_2\in M_F-\Sigma_G$ such that~$\chi$ and~$\rho^W_{\chi}$ are unramified at~$v_1$ and~$v_2,$ with $\# \kappa_{v_1}\leq \#\kappa_{v_2}$, i.e. with $\kappa_{v_1}\subset\kappa_{v_2}$. Let $\phi\in G_*(F).$ One has that $$(\widetilde\chi(\Frob_{v_2})(\phi))=(\widetilde\chi(\Frob_{v_1})(\phi))^{[\kappa_{v_2}:\kappa_{v_1}]}.$$
%\end{lem}
%\begin{proof}
%Here for $i=1, 2$, the notation $\wchi(\Frob_{v_1})$ means $\wchi_{v_1}(\widetilde{\Frob_{v_i}}),$ where~$\widetilde{\Frob_{v_i}}$ is a lift of~$\Frob_{v_i}$. WHY IT DOESNOT DEPEND?
%Let $v_1, v_2\in M_F-\Sigma_G$ such that~$\chi$ and~$\rho^W_{\chi}$ are unramified at~$v_1$ and~$v_2,$ with $\# \kappa_{v_1}\leq \#\kappa_{v_2}$, i.e. with $\kappa_{v_1}\subset\kappa_{v_2}$. We will establish for every~$Q\in (W/\Gamma_F)$ that $$[\kappa_{v_2}:\kappa_{v_1}]\cdot \wchi(\Frob_{v_1})\big(\prod_{\phi\in Q}\phi\big)=\wchi(\Frob_{v_2})\big(\prod_{\phi\in Q}\phi\big).$$ Denote by $\overline{\Frob}_{v_i}\in\Gamma_{\kappa_{v_i}}$ the Frobenius element for $i=1,2$. By the fact that the inertia groups~$I_{v_i}$ act trivially on $G^*(\oF)$ and the fact that $\widetilde\chi_{v_i}$ is~$I_{v_i}$ invariant, we deduce factorization $$\widetilde\chi_{v_i}:\Gamma_{v_i}\to \Gamma_{\kappa_{v_i}}\xrightarrow{\overline\chi_{v_i}} G^*(\oF)$$ with $\overline\chi_{v_i}$ crossed homomorphism satisfying that $\overline\chi_{v_i}(\overline{\Frob}_v)=\widetilde\chi_{v_i}(\Frob_v)$. But, one has that $\overline{\Frob}_{v_2}=[\kappa_{v_2}:\kappa_{v_1}]\cdot \overline{\Frob}_{v_2}$ and using that $$ $$, we obtain that \begin{multline*}\widetilde\chi(\Frob_{v_2})=\widetilde\chi_{v_2}(\Frob_{v_2})=\widetilde\chi_{v_2}([\kappa_{v_2}:\kappa_{v_1}]\cdot \Frob_{v_1})\\=[\kappa_{v_2}:\kappa_{v_1}]\widetilde\chi_{v_2}(\Frob_{v_1})=[\kappa_{v_2}:\kappa_{v_1}]\overline\chi_{v_1}(\overline\Frob_{v_1})=[\kappa_{v_2}:\kappa_{v_1}]\widetilde\chi(\Frob_{v_1}).\end{multline*}
%\end{proof}
%We denote by~$\mathbb F_q$ the field of constants of~$F$ and by~$q$ its cardinality. Let us denote by $\Frob_{q}\in\Gal(\overline{\mathbb F_q}/\mathbb F_q)$ the automorphism $x\mapsto x^{q}$. We adapt calculations of \cite[Lemma 2.11]{vardrapeau} to analyse the behaviour of the $L$-functions.

%The Frobenius $\Frob_{q}$ acts as a permutation matrix on $\CC\cdot W$, thus, its eigenvalues are roots of unity.
%Note that $m_1$ is the number of $\mathcal Y\in\pi_0(\mathcal J_0BG)_{\min}$ which are fixed by the action of the Frobenius, i.e. the number of $\mathcal Y\in\pi_0(\mathcal J_0BG)_{\min}$ having an $F$-point. We have denoted this number by $b(G)$. 
%\begin{lem}\label{lfunff} %Let $\chi\in H^1(F, G^*(\oF))$
%Suppose that~$F$ is a function field.
%\begin{enumerate}
%\item The function $s\mapsto L(s,\rho^W_{\chi})$ is a rational function in $q_F^{-s}$ with a possible pole at $s=1$ of order at most $\# (W/\Gamma_F)$.
%\item Suppose that $\chi=1$. One has that $s\mapsto L(s,\rho^W_{\chi})$ has a pole at $s=1$ which is of order $\# (W/\Gamma_F)$.
%\end{enumerate}
%\end{lem}
%\begin{proof}
%\begin{enumerate}
%\item For $v\in M_F-\Sigma_G$ such that~$\chi$ and ~$\rho^W_{\chi}$ are unramified at~$v$, one has that:
%\begin{align*}
%L_v(s,\rho_{\chi}^W)\hskip-2cm&\\&=\det\big(1-q_v^{-s}\rho^W_{\chi}(\Frob_v)|(\CC\cdot W)\big)^{-1}\\
%&=\prod_{Q\in (W/\Gamma_F)}\det\big(1-q_v^{-s}\rho^Q_{\chi}(\Frob_v)|(\CC\cdot Q)\big)^{-1}\\
%&=\prod_{Q\in (W/\Gamma_F)}\big(1-q_v^{-s\cdot\# Q}\prod_{\phi\in Q}\widetilde\chi (\Frob_v)(\phi)\big)^{-1}\\
%&=\prod_{Q\in (W/\Gamma_F)}\big(1-q_v^{-s\cdot\# Q}\cdot\widetilde\chi(\Frob_v)\big(\prod_{\phi\in Q}\phi\big)\big)^{-1}\\
%&=\prod_{Q\in (W/\Gamma_F)}\bigg(1-(q_{v_1}^{-s\cdot\# Q}\cdot\widetilde\chi(\Frob_{v_1})\big(\prod_{\phi\in Q}\phi\big))^{[\kappa_v:\kappa_{v_1}]}\bigg)^{-1}\\
%&=\prod_{Q\in (W/\Gamma_F)}\big(1-(q_F^{-s\cdot\#Q }\alpha_{Q})^{[\kappa_v:\mathbb F_q]}\big)^{-1}.
%\end{align*}
%We deduce that FORMULA FOR Z IS PROBABLY FALSE, I NEED EIGENVALUES
%\begin{align*}
%L^{\Sigma}(s,\rho_\chi^W)&=\prod_{v\in M_F^0-\Sigma}L_v(s,\rho_\chi^W)\\&=\prod_{v\in M_F-\Sigma}\prod_{Q\in (W/\Gamma_F)}\big(1-(q_F^{-s\cdot \# Q}\alpha_Q)^{ [\kappa_v:\mathbb F_q]}\big)^{-1}\\
%%&=\prod_{Q\in (W/\Gamma_F)}Z(S, q_F^{-s\cdot \#Q}\alpha_Q)\times\\&\quad\quad\quad\quad\quad\times\prod_{v\in\Sigma}\prod_{Q\in (W/\Gamma_F)}\big(1-(q_F^{-s\cdot \#Q }\alpha_{Q})^{[\kappa_v:\mathbb F_q]}\big),
%\end{align*}
 %The crossed homomorphism $\widetilde\chi|{\Gamma_{v_i}}=\widetilde\chi_{v_i}$ 
%The convergence for $\Re(s)>1$ follows from the fact that the eigenvalues of Frobenius are roots of unity. 
%\item The primes $v\in M_F^0-\Sigma$ define an open subscheme $U$ of the scheme $S.$ The representation $\rho^W_{\wchi}$ is unramified at the places in $M_F^0-S$, i.e. it factors through $\pi_1(U, \overline \eta)$ for some $\eta \in U(\overline{\mathbb F}_q)$. The induced representation $\rho^W_{\wchi}:\pi_1(U,\overline\eta)\to\GL(\CC\cdot W)$ defines a locally constant constructible sheaf $E'$ on $U$ such that $E'_{\overline{\eta}}=V$ as $\pi_1(U, \overline\eta)$-modules and a sheaf $E=j_*E'$. One has that 
%\begin{align*}
%L^{\Sigma}(s,\rho^W_{\wchi})&=\prod_{v\in M_F^0-\Sigma}\det\big(1-q^{-s}\rho^W_{\wchi}(\Frob_v)\big)^{-1}\\
%&=Z(S, E, q^{-s})\cdot \prod_{v\in\Sigma}\det\big(1-q^{-s}\Frob_v^{[\kappa(v):\mathbb F_q]}|E_{\overline v}\big),
%\end{align*}where $\overline v$ is a $\overline{\mathbb F_q}$-point of $X$ mapping to $v$ and $$Z(S, E, t)=\prod_{v\in M_F^0}\det\big(1-t\cdot\Frob_v^{[\kappa(v):\mathbb F_q]}|E_{\overline v}\big).$$ Let $\Frob_q\in\Gamma_{\mathbb F_q}$ be the Frobenius element and let $\overline E$ be the pullback of~$E$ for the canonical morphism $\overline S=S\times_{\mathbb F_q}\overline {\mathbb F_q}\to S$. By \cite[Chapter VI, Theorem 13.3]{Milne}, one has that \begin{align*}Z(S, E, q^{-s})\hskip-2cm&\\&=\dfrac{\det\big(1-q^{-s}\Frob_q|H^1_{\et}(\overline S,\overline E)\big)}{\det\big(1-q^{-s}\Frob_q|H^0_{\et}(\overline S,\overline E)\big)\cdot \det\big(1-q^{-s}\Frob_q|H^2_{\et}(\overline S,\overline E)\big)}\\
%&=\dfrac{\det\big(1-q^{-s}\Frob_q|H^1_{\et}(\overline S,\overline E)\big)}{(1-q^{-s})\cdot \det\big(1-q^{-s}\Frob_q|H^2_{\et}(\overline S,\overline E)\big)}.
%\end{align*}
%\item We write $\overline{\Frob_v}$ for the linear map $\CC\cdot W\to \CC\cdot W$ given by $(x_{\phi})_{\phi}\mapsto (q_Fx_{\phi})_{\phi}$. We note that for $v\in M_F^0$, one has that $\Frob_v=\Frob_q^{[\kappa(v):\mathbb F_q]}$. Let $(\alpha_i)_{i\in I}$ be the set of the eigenvalues of $\rho^W_{\wchi}(\Frob_q)$ and let $(m_i)_{i\in I}$ be the set of their multiplicities. We have that:
%Let $S_0$ be the set of closed points of $S$. One has that 
%\begin{align*}L^\Sigma(s, \rho^W_{\wchi})&:=\prod_{v\in M_F^0-\Sigma}\det\big(1-q_v^{-s}\cdot\rho^W_{\widetilde\chi}(\Frob_v)\big)^{-1}
%\\&=\prod_{v\in M_F^0-\Sigma}\det\big(1-q^{-s[\kappa(v):\mathbb F_q]}\cdot\rho^W_{\wchi}(\Frob_{q})^{[\kappa(v):\mathbb F_{q}]}\big)^{-1}\\
%&=\prod_{v\in M_F^0-\Sigma}\det\big(1-(q^{-s}\cdot\rho^W_{\wchi}(\Frob_{q}))^{[\kappa(x):\mathbb F_{q}]}\big)^{-1}\\
%&=\prod_{v\in M_F^0-\Sigma}\prod_{i\in I}\big(1-(q^{-s}\alpha_i)^{[\kappa(v):\mathbb F_q]}\big)^{-m_i} \\
%&=\bigg(\prod_{i\in I}Z(S, q^{-s}\alpha_i)^{m_i}\bigg)\cdot \prod_{v\in\Sigma}\prod_{i\in I}\big(1-(q^{-s}\alpha_i)^{[\kappa(v):\mathbb F_q]}\big)^{m_i},
%\end{align*}
%where $Z(S,t)$ is the zeta function of $S=\Spec(F[t])$ given by $$Z(S,t)=\exp\bigg(\sum_{n>0}\frac{|S(\mathbb{F}_{q^n_F})|t^n}{n}\bigg).$$
%Recall that Weil's theorem \cite[Theorem 4]{basicnt}, one has$$Z(S,t)=\frac{\det(1-t\Frob_{q}|H^1_{\et}(S_{\overline{F_q}}))}{(1-t)(1-q_Ft)}$$where $\Frob_q:\overline{\mathbb F_q}\to\overline{\mathbb F_q}$ is given by $x\mapsto x^{q_F}$.
%Hence, one has that \begin{align*}L^\Sigma(s,\rho^W_{\chi})\hskip-2cm&\\&=\prod_{Q\in (W/\Gamma_F)}\frac{\det(1-(q_F^{-s\cdot \#Q}\alpha_Q)\Frob_{q}|H^1_{\et}(S_{\overline{F_{q}}}))}{(1-q_F^{-s\cdot \#Q}\alpha_Q)(1-q_F^{1-s\cdot\#Q}\alpha_Q)}\times\\&\quad\quad\quad\quad\quad\times\prod_{v\in\Sigma}\prod_{Q\in (W/\Gamma_F)}\big(1-(q_F^{-s\cdot\#Q}\alpha_Q)^{[\kappa(v):\mathbb F_q]}\big).
%\end{align*}
%We deduce that $s\mapsto L^\Sigma(s,\rho^W_{\chi})$ is a rational function in $q_F^{-s}$ and that it has a pole at $s=1$ if and only if~$1$ appears among $\{\alpha_Q\}_Q$. Suppose that~$1$ appears among $\{\alpha_Q\}_{Q\in (W/\Gamma_F)}$ and that $\alpha_{Q_0}=1$. The order of $s\mapsto L^\Sigma(s,\rho^W_{\chi})$ at $s=1$ is equal to~$\#Q_0$. The matrix~$\rho^W_{\wchi}(\Frob_q)$ is a permutation matrix having $e$-th roots of unity for non-zero elements. The dimension of the eigenspace for the eigenvalue~$1$ for every permutation block is no more than~$1$. It follows that the multiplicity~$m_{Q_0}$ is not greater than the number of permutation blocks, i.e. not greater than the number of $\Gamma_F$-orbits $\# (W/\Gamma_F)$. The claim follows.
%
%The function $$s\mapsto \prod_{i\in I}\bigg(\frac{\det(1- q^{-s}\alpha_i\Frob_q|H^1_{\et}(S_{\overline {F_q}})}{1-q^{-s}\alpha_i}\bigg)^{m_i}\prod_{v\in\Sigma}\prod_{i\in I}\big(1-(q^{-s}\alpha_i)^{[\kappa(v):\mathbb F_q]}\big)^{m_i}$$is holomorphic and non-vanishing in the domain $\{\Re(s)>0\}$. Moreover, $s\mapsto \frac{1}{(1-q^{-s+1}\alpha_i)^{m_i}}$ is meromorphic in the domain $\Re(s)>0$ and has a pole at $s=ss1$ if and only if $1$ appears among the eigenvalues $\{\alpha_i\}_i$. It follows that $s\mapsto L^\Sigma(s,\rho^W_{\wchi})$ has a pole at $s=1$ if and only if  Moreover, if $1$ appears among $\{\alpha_i\}_i$, wlog suppose that $1=\alpha_1,$ then the order of the pole is equal to the multiplicity $m_1$. As $\rho^W_{\wchi}(\Frob_q)$ is a quasi-permutation matrix, the multiplicity of the eigenvalue $1$ is not bigger than the number of quasi-permutation blocks, i.e. not bigger than the number of orbits $W/\Gamma_F$. 
%\item If $\chi=1$ is the trivial character, then the crossed homomorphism $\gamma\mapsto 1\in G^*(\oF)$ is a lift of $\chi$ and thus $\rho^W_{\wchi}(\Frob_q)$ is a permutation matrix having only $1$-s for non-zero elements. The dimension of the eigenspace for the eigenvalue $1$ for each permutation block is then equal to $1$, and thus the multiplicity of the eigenvalue $1$ is equal to the number of permutation blocks, i.e. the number of $\Gamma_F$-orbits~$\# (W/\Gamma_F)$.
%\end{enumerate}
%\end{proof}
\subsection{Analysis of height zeta functions}\label{Fouriertransformsa} In this subsection, we will define and analyse the height zeta functions. %study analytic properties of certain Fourier transforms.  %We will use the results in Subsection \ref{poisson}. 
Let $c:G_*(\oF)\to\QQ_{\geq 0}$ be a normalized counting function and let~$H$ be a height having~$c$ for its type. Let $\Sigma_H\supset \Sigma_G\cup M_F^{\infty}$ be the finite set of places such that for~$v\in M_F^0-\Sigma_H,$ one has that $H_v:BG(F_v)\to \RR_{>0}$ is given by$$H_v(x)=q_v^{c(\psi^G_v(x))}.$$ %For~$v\in \Sigma_H$, we choose functions $H_v:BG(F_v)\to\RR_{>0}$. We obtain a height~$H:BG(F)\to \RR_{>0}$.
The definition of height extends to $H:BG(\AAA)\to \RR_{>0}$ by $$H((x_v)_v)=\prod_{\vMF}H_v(x_v)$$so that for $x\in BG(F)$ one has that $H(x)=H(i(x)),$ where $i:BG(F)\to BG(\AAA)$ is the diagonal map. 
\subsubsection{} The following lemma will be used later.
\begin{lem}\label{codx}
Let $x\in BG(\AAA).$ We define $$\Sigma_x:=\Sigma_H\cup\{v\in M^0_F-\Sigma|i_v(x)\not\in B\GG(\Ov)\}.$$  We set $$C(x):=\prod_{v\in\Sigma_x}\frac{\max_{z\in BG(F_v)}H_v(z)}{\min_{z\in BG(F_v)}H_v(z)}.$$ For every $y\in BG(\AAA)$, one has that $$C(x)^{-1}H(y)\leq H(xy)\leq C(x)H(y).$$
\end{lem}
\begin{proof}
For $v\in M_F-\Sigma_x$, one has that $H_v(x)=1.$ For $v\in\Sigma_x$, we set $$C_v(x):=\frac{\max_{z\in BG(F_v)}H_v(z)}{\min_{z\in BG(F_v)}H_v(z)},$$ so that $C(x)=\prod_{v\in\Sigma_x}C_v(x).$ For every $y\in BG(\AAA)$, one has that \begin{align*}
H(xy)&=\prod_{v\in M_F}H_v(xy)\\
&=\prod_{v\in\Sigma_x}H_v(xy)\cdot\prod_{v\in M_F-\Sigma_x}H_v(y)\\
&\leq \prod_{v\in\Sigma_x}C_v(x)H_v(y)\cdot\prod_{v\in M_F-\Sigma_x}H_v(y)\\
&=C(x)H(y).
\end{align*}
Analogously, one verifies that $C(x)^{-1}H(y)\leq H(x)$.
\end{proof}
For $s\in\CC$ and $\chi_v\in (BG(F_v))^*$,  we define $$\wH_v(s,\chi_v):=\int_{BG(F_v)}H_v^{-s}\overline{\chi_v}\mu_v.$$ 
For every $\chi_v\in BG(F_v)^*,$ the function $s\mapsto\wH_v(s,\chi_v),$ being a sum of finitely many entire functions, is an entire function.

For a character $\chi\in BG(\AAA)^*=BG^*(\AAA)$ we denote by $$ \wH(s,\chi):=\int_{BG(\AAA)}H^{-s}\overline{\chi}\mu.$$
For a place $v,$ the set given by the elements $x\in BG(F_v)$ such that~$H_v$ is invariant for the multiplication by~$x$, forms an open subgroup of the discrete group $BG(F_v)$. We denote it by~$K_{H_v}$. Note that for $v\in M_F^0-\Sigma_H$, one has that $K_{H_v}=B\GG(\Ov).$ We denote by~$K_H$ the compact and open subgroup $$K_{H}:=\prod_{\vMF}K_{H_v}\subset BG(\AAA).$$
We denote $$\Xi_H:=i(BG(F))K_H.$$ 
The following lemma identifies the finite group of characters for which the Fourier transform does not vanish.
\begin{lem}
\label{importantchar}
\begin{enumerate}
\item The subgroup $\Xi_H^{\perp}\subset BG(\AAA)^*$ is finite.
\item The homomorphism $$BG(\AAA)=(BG(\AAA)^*)^*\to (\Xi_H^\perp)^*,\hspace{1cm} x\mapsto x|_{\Xi_H^\perp}$$ is surjective and its kernel identifies is the group $\Xi_H$. In particular, the subgroup $\Xi_H$ is open, closed and of the index $\#\Xi_H^\perp$ in $BG(\AAA)$.
\item There exists a finite set of places~$(\Sigma_G\cup M_F^{\infty})\subset\Sigma$ and a height~$H'$ for which~$H_v'=1$ if $v\in \Sigma$ and $H_v'=q_v^{c_v(\psi_v(x))}$  if $v\in M_F-\Sigma$, such that $$\Xi_{H'}=BG(\AAA).$$
\item Let $\chi \in (BG(\AAA)^*-\Xi_H^{\perp})$ be a character. One has that $\wH(s,\chi)=~0$.
\end{enumerate}
\end{lem} 
\begin{proof}
\begin{enumerate}
\item %Let $\Sigma '$ be the finite subset of $M_F$ given by the union of the set $\Sigma_H$ and the set of places $v$ for which $K_{H_v}\not\supset B\GG(\Ov).$ 
By \cite[Chapter 2, \S 1, \no 7, Corollary 2 of Theorem 4]{TSpectrale} and by Lemma \ref{smallemch}, one has that
\begin{align*}
\Xi_H^\perp&=(i(BG(F))K_H)^\perp\\&=(i(BG(F))^{\perp}\cap K_H^{\perp}\\&=i(BG^*(F))\cap\bigg(\prod_{v\in\Sigma_H}K_{H_v}^{\perp}\times\prod_{v\in M_F^{0}-\Sigma_H}B\GG^*(\Ov)\bigg)
\end{align*}
The finiteness of the later intersection follows from Lemma~\ref{finlocun}.
\item The homomorphism $BG(\AAA)\to (\Xi_H^\perp)^*$ is the dual homomorphism \cite[Chapter II, \S 1, \no 7, Page 124]{TSpectrale} of the inclusion of a finite, and hence closed, subgroup $\Xi_H^\perp\subset BG(\AAA)^*,$ hence by \cite[Chapter II, \S 1, \no 7, Theorem 4]{TSpectrale} is surjective and its kernel is the group $(\Xi_H^\perp)^{\perp}=\overline{\Xi_H}$. But, the subgroup $\Xi_H\subset BG(\AAA)$ is the subgroup generated by a discrete subgroup $i(BG(F))$ and a compact subgroup $K_H$, hence, by \cite[Chapter III, \S 4, \no 1, Corollary 1 of Proposition 1]{TopologieGj}, is closed in $BG(\AAA)$. Thus $\Xi_H=\overline{\Xi_H}=(\Xi_H^\perp)^\perp$ and, hence, $\Xi_H$ is the kernel of $BG(\AAA)\to (\Xi_H^\perp)^*$.
%. Now, by \cite[Chapter II, \S 1, \no 7, Theorem 4]{TSpectrale}, one has that is the kernel of the map $BG(\AAA)\to \Xi_H^\perp$ which is (by the same theorem) surjective. It follows that the index of $\Xi_H$ is equal to the cardinality $|(\Xi_H^\perp)^*|=|\Xi_H|$. In particular, $\Xi_H$ is open. The claim is proven. 
\item By Part (2), we need to find~$H'$ for which $\Xi_{H'}^{\perp}$ is a singleton. There exists a finite $\Sigma_H\cup\Sigma_G\cup\Sigma_{G^*}\cup M_F^{\infty}\subset \Sigma\subset M_F$ such that $$\Xi_H^{\perp}\cap \big(\prod_{v\in\Sigma}BG^*(F_v)\times\prod_{v\in M_F^0-\Sigma}BG^*(\Ov)\big)=\{1\}.$$ Now, let us set~$H'_v=1$ for $v\in\Sigma$ and $H'_v=H_v$ for $v\not\in\Sigma$, so that $K_{H}\subset K_{H'}$ and that $K_{H'}=\prod_{v\in\Sigma}BG^*(F_v)\times\prod_{v\in M_F^0-\Sigma}BG^*(\Ov).$ Thus $$\Xi_{H'}^{\perp}=(i(BG(F))K_{H'})^{\perp}=(i(BG(F))K_HK_{H'})^{\perp}\subset\Xi_H^{\perp}\cap K_{H'}=\{1\}$$and the claim follows. 
\item Let $\vMF$ be such that $\chi_v|_{K_{H_v}}\neq 1$. As $x\mapsto H_v(x)^{-s}$ is $K_{H_v}$-invariant, we deduce that $\wH_v(s,\chi_v)=0$. It follows that $\wH(s,\chi)=0$.
\end{enumerate}
\end{proof}
%Denote by $\pi_0(\mathcal{J}_0BG)_{\min}$ the subset of $\pi_0(\mathcal{J}_0BG)$ given by the components having the minimal age with the respect to $V$. Note that $\pi_0(\mathcal{J}_0BG)_{\min}$ is a $\Gal (\overline{F}/F)$ submodule of the $\Gal (\overline{F}/F)$ module $\pi_0(\mathcal{J}_0BG)$. We denote by $\ZZ\piomin$ the free abelian group generated by $\piomin$. If $F$ is a number field, we let $\Frob_v\in\Gal(\overline{F}/F)$ be a Frobenius element corresponding to the place $v$. In the function field case, we set $\Frob_v\in\Gal(\overline F/F)$ to be the element that maps $x$ to $x^{q_F}$.  For a finite place $v$, we denote by $$L_v(s,\pi_0(\mathcal{J}_0BG)_{\min}):=\det(1-q_v^{-s}\Frob_v|{(\ZZ\pi_0(\mathcal{J}_0BG)_{\min})_{\QQ}})^{-1}, $$ the corresponding local Artin $L$-function. We denote by $$L(s,\pi_0(\mathcal{J}_0BG)_{\min})):=\prod_{v\in M_F^0}L_v(s,\pi_0(\mathcal{J}_0BG)_{\min})$$ the global Artin $L$-function. We denote $$a(G):=(\min_{\mathcal{Y}\in\pi_0(\mathcal J_0BG)}\age(\mathcal{Y},V))^{-1}$$ and $$b(G):=|\{\mathcal Y\in\pi_0(\mathcal J_0BG)_{V-\min}|\mathcal Y(F)\neq \emptyset\}|.$$
%\subsection{Growth of $L$-function in the number field case}
%\begin{lem}\label{lnumbf}
%The function $s\mapsto L(s,\pijbg)$ is holomorphic in $s$ in the domain $\Re(s)>{1}$. It admits a pole of order $b(G)$ at $s=1$ and there exists $\epsilon>0$ such that $s\mapsto L(s,\pijbg)$ extends meromorphically to the domain  $\Re(s)>1-\epsilon$. Moreover, there exists $\delta >0$ such that for every compact $\mathcal K\subset ]a(G)-\epsilon,+\infty[,$ there exists $C(\mathcal K)>0$ such that $$\bigg|\frac{(s-1)^{b(G)}L(s,\pijbg)}{s^{b(G)}}\bigg|\leq C(\mathcal K)(1+|\Im(s)|)^{\delta}$$whenever $\Re(s)\in K$.
%\end{lem}
\subsubsection{} 
%\begin{lem}Let $\chi: BG(\AAA)\to S^1:=\{z\in\CC, |z|=1\}$ be a character. The Fourier transform $\wH(s,\chi)$ converges absolutely and uniformly for every $s$ in the domain $\Re(s)>1$ and defines a holomorphic function in this domain. Moreover, $\wH(s,\chi)$ extends to a meromorphic function in the domain $\Re(s)\geq 1$ with a possible pole at $s=1$ of order at most $b(G)$.
%\end{lem}
We will compare the Fourier transform at the character~$\chi$, with the $L$-function of a representation $\rho^W_{\chi}$. Recall the notation $G_*(\oF)_1=\{x\in G_*(F)| c(x)=1\}.$ % and $$G_*(F_v)_1=(G_*(\oF)_1)\cap G_*(\Fv).$$
\begin{lem}\label{ftlfun} Let $\chi\in H^1(\Gamma_F,G^*(\oF))\cap K^{\perp}\subset BG(\AAA)^*$ be a character and let~$\Sigma\subset M_F^0$ be finite. The product %s $\epsilon>0$, such that
$$\prod_{\vMFz-\Sigma}\frac{{\wH_v(s,\chi_v)}}{L_v(s,\rho_{\chi}^{G_*(\oF)_1})}$$ converges normally on the compacts in the domain $\Re(s)>~\lambda^{-1}$ to a holomorphic function, where \begin{align*}
1<\lambda:=\begin{cases} \min_{x\in (G_*(F_v)-(\{1\}\cup G_*(F)_1))}c(x), \text{if } \{1\}\cup G_*(F_v)_1\subsetneq G_*(F_v),\\
2\quad\quad\quad\quad\quad\quad\quad\quad\quad\quad\quad\quad\text{, otherwise}.
\end{cases}
\end{align*}
(Here, $G_*(\Fv)_1=(G_*(\oF)_1)^{\Gamma_v}$.)
\end{lem}
\begin{proof}
Let us set $\Sigma_{\chi}$ to be the union of~$\Sigma_H$ and the finite set of places for which~$\chi$ is ramified. By Lemma \ref{traceofl}, for $\vMFz-\Sigma_\chi$ and $s\in\{z:\Re(z)>0\}$, one has that
\begin{align*}
L_v\big(s,\rho^{G_*(\oF)_1}_{\chi}\big)^{-1}&=\det\big(1-q_v^{-s}\rho^{G_*(\oF)_1}_{\chi}(\Frob_v)\big)\\
&=1-q_v^{-s}\Tr\big(\rho^{G_*(\oF)_1}_{\chi}(\Frob_v)\big)+ O(q_v^{-2\Re(s)})\\
&=1-\bigg(\frac{q_v^{-s}}{\#G(F_v)}\sum_{y\in(\psi^G_v)^{-1}(G_*(F_v)_1)}\chi_v(y) \bigg)+O(q_v^{-2\Re(s)}).
\end{align*}
%Let $\frac{1}{2}>\eta>0$ be chosen so that for every $\mathcal{Y}\in\pi_0(\mathcal{J}_0BG)$ with $\age(\mathcal Y,V)>a(G,V),$ one has that $$a(G,V)\eta<\age(\mathcal Y,V).$$
%It follows from part (2) of Lemma \ref{gfv} that there exists a finite set of finite places $\Sigma'$ containing $\Sigma$ such that for every $v\in M_F^0-\Sigma'$ one has that $\#B\GG(\Ov)=\#G(F_v)$  and that $H_v(x)=q_v^{c(\psi_v(x))}$. 
On the other side, for $\vMF^0-\Sigma_{\chi}$ and $s\in\{z:\Re(z)>0\}$, one has that \begin{align*}
\wH_v(s,\chi_v)&=\int_{BG(F_v)}H_v^{-s}\overline{\chi_v}\mu_v\\
&=\frac{1}{\#G(F_v)}\sum_{x\in BG(F_v)} H_v(x)^{-s}\chi_v(x)\\
&=\frac{1}{\# G(F_v)}\bigg(\sum_{x\in B\GG(\Ov)}1 + q_v^{-s}\sum_{x\in (\psi^G_v)^{-1}(G_*(F_v)_1)}\chi_v(x)+\\&\quad\quad\quad\quad\quad+\sum_{x\in BG(F_v)-\big(B\GG(\Ov)\cup(\psi^G_v)^{-1}(G_*(F_v)_1) \big)}H_v(x)^{-s}\chi_v(x)\bigg)\\
&=1+\bigg(\frac{q_v^{-s}}{\#G(F_v)}\sum_{y\in(\psi_v^G)^{-1}(G_*(\Fv)_1)}\chi_v(y)\bigg)+ O(q_v^{-\Re(s)\cdot\lambda}).
\end{align*}
%one thus has that \begin{align*}
%\wH_v(s,\chi_v)&=\int_{BG(F_v)}H_v^{-s}\overline{\chi_v}\mu_v\\
%&=\frac{1}{|G(F_v)|}\sum_{x\in BG(F_v)} H_v(x)^{-s}\overline{\chi_v(x)}\\
%%&=\frac{1}{|B\GG(\Ov)|}\bigg(\sum_{x\in B\GG(\Ov)}1 + \sum_{x\in BG(F_v)_{V-\min}}H_v(x)^{-s}\overline{\chi_v(x)}+\\&\quad\quad\quad\quad\quad\quad\quad+\sum_{x\in BG(F_v)-\big(B\GG(\Ov)\cup BG(F_v)_{V-\min}\big)}H_v(x)^{-s}\overline{\chi_v(x)}\bigg)\\
%&=1+q_v^{-s}\Tr(\Frob_v|\CC\pi_0(\mathcal J_0BG))+O(q_v^{-\Re(s)(1+\eta)})
%\end{align*}
Hence, for $v\in M_F^0-\Sigma_\chi$ and $\Re(s)>0$ one has that $$L_v\big(s,\rho_{\widetilde\chi}^{G_*(\oF)_1}\big)^{-1}\wH_v(s,\chi_v)=1+O(q_v^{-\Re(s)\cdot \lambda}).$$ It follows from \cite[Chapter VII, \S 5, Corollary 1]{basicnt} that the product $$\prod_{v\in (M_F^0-\Sigma)-\Sigma_\chi}L_v(s,\rho_{\wchi}^{G_*(\oF)_1})^{-1}\wH_v(s,\chi_v),$$ %and hence the product $$\prod_{v\in M_F^0}L_v(s,\rho_{\wchi}^{G_*(\oF)_1})^{-1}\wH_v(s,1)$$
converges normally on the compacts to a holomorphic function in the domain $\Re(s)>\lambda^{-1}.$ 
Hence the product $$\prod_{v\in M_F^0-\Sigma}L_v(s,\rho_{\chi}^{G_*(\oF)_1})^{-1}\wH_v(s,\chi_v),$$ converges normally on the compacts to a holomorphic function in the domain $\Re(s)>\lambda^{-1}.$ 
\end{proof}

\begin{cor}\label{trivialcharacter}
Let $\chi\in H^1(\Gamma_F, G^*(\oF))\cap K_H^\perp\subset BG^*(\AAA)$ be a character and let $\lambda>1$ be as in Lemma \ref{ftlfun}.  Let $\Sigma\subset M_F$ be finite.
\begin{enumerate}
\item The product defining $\wH_{M_F-\Sigma}(s, \chi):=\prod_{v\in M_F-\Sigma}\wH_v(s,\chi_v)$ converges normally on the compacts in the domain $\Re(s)>1$. The function $s\mapsto\wH_{M_F-\Sigma}(s,\chi)$ extends to a meromorphic function in the domain $\Re(s)> \lambda^{-1}$ with a possible pole at $s=1$ of order at most $b(c)=(G_*(\oF)_1/\Gamma_F)$. In the case~$F$ is a number field, one also has that $s\mapsto\wH_{M_F-\Sigma}(s,\chi)$ is holomorphic in the domain $\{\Re(s)\geq 1\}-\{1\}$. In the case~$F$ is a function field, one also has that $s\mapsto\wH_{M_F-\Sigma}(s,\chi)$ is holomorphic in the domain $\{\Re(s)\geq 1-\epsilon\}-\{1+2ki\pi\log(q)^{-1}|k\in\ZZ\}$ for some $\epsilon>0$.
\item If $\chi=1$ is the trivial character, then the meromorphic function $s\mapsto\wH(s,1)=\wH_{M_F}(s,1)$ admits a pole at $s=1$ of order $b(c)$. Moreover, one has that $$\lim_{s\to 1^+}(s-1)^{b(c)}\wH(s,1)=\lim_{s\to 1^+}(s-1)^{b(c)}\int_{BG(\AAA)}H^{-s}\mu>0.$$
\end{enumerate}
\end{cor}
\begin{proof}
\begin{enumerate}
%\item This claim follows by combining Lemma \ref{ftlfun}, Parts (1) and (3) of Lemma \ref{lfunf} and the fact that the functions  $s\mapsto\wH_v(s,\chi_v)$ are entire.
\item The claims follow by combining Lemma \ref{ftlfun}, Lemma \ref{lfunf} and the fact that the functions  $s\mapsto\wH_v(s,\chi_v)$ are non vanishing for $\Re(s)\geq 0$ if $\chi=1$.
\item From Lemma \ref{ftlfun}, Part (2) of Lemma \ref{lfunf} and the fact the functions  $s\mapsto\wH_v(s,1)$ are non vanishing for $\Re(s)\geq 0$, 
% one has that $s\mapsto L^\Sigma(s,\rho^{G_*(\oF)_1}_{1})$ has a pole at $s=1$ of order precisely $b(c).$ 
we deduce that $s\mapsto\wH(s,1)$ has a pole at~$s=1$ of order~$b(c)$. The principal value at the pole is equal to \begin{align*}\lim_{s\to 1^+}(s-1)^{b(c)}\wH(s,1)=\lim_{s\to 1^+}(s-1)^{b(c)}\int_{BG(\AAA)}H^{-s}\mu.
\end{align*} and is strictly bigger than zero, as it is nonzero and for $s>1$ one has that $(s-1)^{b(c)}\wH(s,1)>0.$ %because $\wH(s,1)$ is defined by a converging Euler product.
%\item For $v\in M_F-\Sigma_\chi$, let us set $$\beta_v(s,\chi):=L_v\big(s,\rho_{\chi}^{G_*(\oF)_1}\big)^{-1}\wH_v(s,\chi_v) $$ and for $v\in \Sigma_\chi$, let us set $$\beta_v(s, \chi_v):=\wH_v(s,\chi_v).$$ It follows from Lemma \ref{ftlfun} together with the fact that all functions $s\mapsto\wH_v(s,\chi_v)$ are entire and non-vanishing, that the product $$\beta_{M_F-\Sigma_\chi}(s,\chi):=\prod_{v\in M_F-\Sigma'}\beta_v(s,\chi_v)$$ converges absolutely and normally on the compacts in the domain $\Re(s)> \lambda^{-1}$ to a holomorphic function. The function, furthermore, does not vanish in the domain of convergence. By Lemma \ref{lfunf}, the product $$\prod_{v\in M_F^0-\Sigma'}L_v(s,\rho_{\chi}^{G_*(\oF)_1})$$ converges for $\Re(s)>1$ and $s\mapsto \prod_{v\in M_F-\Sigma'}L(s,\rho_{\chi}^{G_*(\oF)_1})$ extends to a meromorphic function on $\CC$, with a possible pole at $s=1$ of order no more than $b(c)$. Moreover, if $F$ is a number field, the function is holomorphic in $\{\Re(s)\geq 1\}-\{1\}.$ 
%We deduce that 
%\begin{align*}
%s\mapsto\wH_{M_F-\Sigma'}(s,\chi)&=\prod_{v\in M_F-\Sigma'}\wH(s,\chi_v)\\&=\prod_{v\in M_F^0-\Sigma'}\big(L_v(s,\rho_{\chi}^{G_*(\oF)_1})\cdot\beta_v(s,\chi_v)\big)\prod_{\vMFi-\Sigma'}\beta_v(s,\chi)\\
%&=\bigg(\prod_{v\in M_F-\Sigma'}L(s,\rho_{\chi}^{G_*(\oF)_1})\bigg)\cdot \beta_{M_F-\Sigma'}(s,\chi)
%\end{align*}
%converges for $\Re(s)>1$ and extends to a meromorphic function in the domain~$\Re(s)> \lambda^{-1}$, with a possible pole at $s=1$ which is of order no more than $b(c)$. Moreover, if $F$ is a number field, then the function $s\mapsto \wH_{M_F-\Sigma'}(s,\chi)$ is holomorphic in the domain $\{\Re(s)\geq 1\}-\{1\}$. 
%
\end{enumerate}
\end{proof}
The following lemma will be used later.
\begin{lem}\label{intfinind}
Let~$V, W\subset BG(\AAA)$ be open subgroups of finite index. Let~$y_1\in V$. One has that $$\lim_{s\to 1^+}(s-1)^{b(c)}\int_{(y_1W)\cap V}H^{-s}\mu>0.$$ 
\end{lem}
\begin{proof}
Let $y_1\doots y_r$ be a set of representatives of classes of~$W$ in~$BG(\AAA)$. For~$s>1$, by Lemma~\ref{codx}, one has that \begin{align*}\int_{BG(\AAA)}H^{-s}\mu&=\sum_{i=1}^k\int_{y_iW}H^{-s}\mu\\
&=\sum_{i=1}^r\int_{y_1W}H(y_iy_1^{-1}t)^{-s}d\mu(t)\\
&\leq\bigg( \sum_{i=1}^rC_i^{-s}\bigg)\cdot\int_{y_1W}H^{-s}\mu
\end{align*}
for certain $C_i>0$. Now, Part~(2) of Corollary~\ref{trivialcharacter} gives that
\begin{align*}
0&<\lim_{s\to 1^+}(s-1)^{b(c)}\int_{BG(\AAA)}H^{-s}\mu\\
&\leq\lim_{s\to 1^+}(s-1)^{b(c)}\bigg(\sum_{i=1}^rC_i^{-s}\bigg)\cdot\int_{y_1W}H^{-s}\mu\\
&=\bigg(\sum_{i=1}^rC_i^{-1}\bigg)\lim_{s\to 1^+}(s-1)^{b(c)}\int_{y_1W}H^{-s}\mu,
\end{align*}
and hence $$\lim_{s\to 1^+}(s-1)^{b(c)}\int_{y_1W}H^{-s}\mu>0.$$
Note that it follows (by setting $W=V$ and $y_1=1$) that $$\lim_{s\to 1^+}(s-1)^{b(c)}\int_{V}H^{-s}\mu>0.$$ %Let now $x_1\doots x_k$ be a set of representatives of classes of~$V$ in~$BG(\AAA)$. 
Let $y_1', y_2'\doots y_k'\in V$ be a set of representatives of classes of $V\cap W$ in~$V$ such that $y_1'=y_1$. For~$s>1$, by Lemma~\ref{codx} we have that:
\begin{align*}
\int_{V}H^{-s}\mu&=\sum_{i=1}^k\int_{ (y_i')(W\cap V)}(H^{-s})\mu\\
&=\sum_{i=1}^k\int_{(y_1')(W\cap V)}(H((y_i')(y_1')^{-1}t)^{-s})d\mu(t)\\
&\leq \bigg(\sum_{i=1}^k(C'_i)^{-s}\bigg)\int_{(y_1)(W\cap V)}H^{-s}\mu\\
\end{align*}
for some~$C'_i>0$. We obtain that
\begin{align*}
0&<\lim_{s\to 1^+}(s-1)^{b(c)}\int_{V}H^{-s}\mu\\
&\leq\lim_{s\to 1^+}(s-1)^{b(c)}\bigg(\sum_{i=1}^k(C_i')^{-s}\bigg)\cdot\int_{(y_1)(W\cap V)}H^{-s}\mu\\
&=\bigg(\sum_{i=1}^k(C_i')^{-1}\bigg)\lim_{s\to 1^+}(s-1)^{b(c)}\int_{(y_1)(W\cap V)}H^{-s}\mu.
\end{align*}
Finally, by remarking that $(y_1W)\cap V\supset (y_1)(W\cap V)$, we deduce that $$\lim_{s\to 1^+}(s-1)^{b(c)}\int_{(y_1W)\cap V}H^{-s}\mu\geq \lim_{s\to 1^+}(s-1)^{b(c)}\int_{(y_1)(W\cap V)}H^{-s}\mu>0,$$as claimed.
\end{proof}
%\begin{cor}Suppose $F$ is a number field. Let $\chi\in BG(\AAA)^*$ be a character. There exists  $\epsilon>0$ such that the meromorphic function $s\mapsto (s-1)^{b(c)}\wH(s,\chi)$ is holomorphic in the domain $\Re(s)>1-\epsilon$. Moreover, there exists $\delta>0$ such that for every compact $\mathcal K\subset ]1-\epsilon,+\infty[$ one has that there exists $C(\mathcal K)>0$ such that $$\bigg|\frac{(s-1)^{b(c)}\wH(s,\chi)}{s^{b(c)}}\bigg|\leq C(\mathcal K)(1+|\Im(s)|)^{\delta}$$whenever $\Re(s)\in\mathcal K.$
%\end{cor}
%\begin{proof}
%%By 
%\end{proof}
%\begin{cor}\label{fourtranschi}
%Let $\chi\in BG(\AAA)^*.$ One has that the integral defining $\wH(s,\chi)$ converges absolutely and uniformly in the domain $\Re(s)>1$. The function $s\mapsto\wH(s,\chi)$ is holomorphic in this domain. There exists $\epsilon>0$ such that $s\mapsto\wH(s,\chi)$ extends to a meromorphic function in the domain $\Re(s)>1-\epsilon$ with a pole at $s=1$ of order at most $b(c)$.
%\end{cor}
%\begin{proof}
%By Lemma \ref{lnumbf} and by Lemma \ref{lfunf}, in the domain $\Re(s)>~1$, the product $\prod_{\vMFz}L_v(s,\pijbg)$ converges absolutely and uniformly to the holomorphic function $s\mapsto L(s,\pijbg)$. By Lemma \ref{closehl}, there exists $\epsilon'>0$ such that the product $$\prod_{v\in M_F^0}\frac{\wH_v(s,\chi_v)}{L_v(s,\pijbg)}$$ converges absolutely and uniformly to a holomorphic function in the domain $\Re(s)>1-\epsilon'$. This function does not vanish at zero, as it is defined by a converging product. For $v\in M_F^\infty$, the function $s\mapsto\wH_v(s,\chi_v)$ is holomorphic and non vanishing at $s=1$. It follows that $$\wH(s,\chi)=L(s,\chi)\cdot\prod_{v\in M_F^0}\frac{\wH_v(s,\chi_v)}{L_v(s,\pijbg)}\cdot\prod_{\vMFi}\wH_{v}(s,\chi_v)$$ is holomorphic in the domain $\Re(s)>1$, that it extends to a meromorphic function in the domain $\Re(s)>1-\epsilon$ for some $\epsilon>0$ and that it admits a pole of order $b(c)$ at $s=1$.
%\end{proof}
%If $F$ is a number field, we will need the following estimate in order to use Tauberian results.
%\begin{cor}Suppose $F$ is a number field. Let $\chi\in BG(\AAA)^*$. There exists $1>\epsilon >0$, such that the function $s\mapsto (s-1)^{b(c)}\wH(s,\chi)$ is holomorphic in the domain $\Re(s)>1-\epsilon$. Moreover, there exists $\delta>0$ such that for every compact $\mathcal K$ in the domain $]1-\epsilon,\infty[$, there exists $C(\mathcal K)>0$ such that $$\bigg|\frac{(s-1)^{b(c)}Z(s)}{s^{b(c)}}\bigg|\leq C(\mathcal K)(1+|\Im(s)|)^{\delta}$$ if provided $\Re(s)\in\mathcal K$. 
%\end{cor}
%\begin{proof}
%It follows from Lemma \ref{lnumbf} that $L(s,\chi)$ satisfies that 
%\end{proof}
\subsubsection{}In this paragraph we apply the Poisson formula for the height zeta function.
\begin{mydef}
For $s\in\CC$, by $Z(s)$ we denote the formal (non-empty) sum $$Z(s):=\frac{1}{\#G(F)}\sum_{x\in BG(F)}H(x)^{-s}.$$
\end{mydef}
The following fact follows from Proposition \ref{estimatenumber} and \cite[Chapter~2, Theorem 2.1]{widder}.
\begin{prop}\label{convz}
There exists $\eta>0$ such that the series defining~$Z$ converges absolutely and uniformly for $\Re(s)>\eta$. The function $s\mapsto Z(s)$ is holomorphic in the domain $\Re(s)>\eta$.
\end{prop}
The following conditions we need to verify in order to apply the Poisson formula holds.
\begin{lem}\label{condofpoiss}
Let $s\in\CC$ be such that $\Re(s)>\eta$. Let $A>0$ and let $g:BG(\AAA)\to\RR_{\geq 0}$ be a bounded continuous function. For every $x\in BG(\AAA)$, one has that the series $$Z^g(x,s):=\frac{\#\Sh^1(G)}{\# G(F)}\sum_{y\in i(BG(F))}g(xy)^{-s}H(xy)^{-s}$$ converges absolutely. The function $$Z^g(-,s):BG(\AAA)\to\CC,\hspace{1cm}x\mapsto Z^g(x,s)$$ is continuous.
\end{lem}
\begin{proof}
For $x\in \BGA$, we define $$\Sigma_x:=\Sigma\cup\{v\in M^0_F-\Sigma|i_v(x)\not\in B\GG(\Ov)\}$$ and $$U_{x}:=\prod_{v\in\Sigma_x}\{x_v\}\times\prod_{v\in M_F^0-\Sigma_x}B\GG(\Ov).$$ For every $x\in BG(\AAA)$, we are going to prove that the series $Z^g(x,s)$ converges absolutely and uniformly on $U_x$. For every $x'\in U_x$ and every $y\in i(BG(F))$, it follows from Lemma \ref{codx} that there exists $C(x)>0$ such that $$\big|H(x'y)^{-s}\big|=H(x'y)^{-\Re(s)}\leq C(x)^{-\Re(s)}H(y)^{-\Re(s)}.$$ Let $A>0$ be such that for every~$x\in BG(\AAA)$ one has $g(x)\leq A$. We have that 
\begin{align*}
\frac{\# \Sh^1(G)}{\# G(F)}\sum_{y\in i(BG(F))}\bigg|g(xy)^{s}H(xy)^{-s}\bigg|\hskip-3,7cm&\\&\leq \frac{\#\Sh^1(G)}{\# G(F)}\sum_{y\in i(BG(F))}A^{\Re(s)}C(x)^{-\Re(s)}H(y)^{-\Re(s)}\\
&=\frac{A^{\Re(s)}\cdot C(x)^{-\Re(s)}}{\# G(F)}{\sum_{y\in BG(F)}H(y)^{-\Re(s)}}\\
&=A^{\Re(s)}\cdot C(x)^{-\Re(s)}{Z(\Re(s))}.
\end{align*}
Now using Proposition \ref{convz}, we deduce that the series $Z^g(x,s)$ converges absolutely and uniformly in the domain $x'\in U_x$. It follows that $x'\mapsto Z^g(x',s)$ is continuous on $U_x$. We deduce that $x\mapsto Z^g(x,s)$ is continuous on $BG(\AAA)$, as claimed. 
\end{proof}
If $W$ is a locally compact abelian group endowed with a Haar measure~$dw$, we will denote by $dw^*$ the Haar measure on the character group~$W^*$ which is the dual of $dw$ \cite[Chapter II, \S 1, \no 3, Definition~4]{TSpectrale}. It is characterized by the following property: it is the unique Haar measure on~$W^*$ such that for every $f\in L^1(W, dw)$ for which the Fourier transform $\widehat f:\chi\mapsto\widehat f(\chi)=\int_{W}f\overline{\chi}dw$ satisfies $\widehat f \in L^1(W^*, dw^*)$, the \textit{Fourier inversion formula} \cite[Chapter II, \S 1, \no 4, Proposition 4]{TSpectrale} is valid: \begin{equation}\label{fourierinversion}f(w)=\int_{W^*}\overline{\chi(w)}\widehat{f}(\chi) dw^*(\chi)
\end{equation} for every $w\in W$.

In the next proposition, we apply the Poisson formula \cite[Chapter II, \S1, \no 8, Proposition 8]{TSpectrale}.
\begin{prop}\label{poisson}
In the domain $\Re(s)>\max(\eta,1)$,  one has that $$Z(s)=\frac{1}{\tau_{BG}}\sum_{\chi \in \Xi_H^{\perp}}\int_{BG(\AAA)}H^{-s}\overline{\chi}\mu.$$ 
\end{prop}
\begin{proof}
We will denote by $d\chi$ the dual Haar measure on the group $(BG(\AAA)/i(BG(F)))^*$ of the Haar measure $\mu/\coun_{i(BG(F))}$ on the group $BG(\AAA)/i(BG(F)).$ The space $(BG(\AAA)/i(BG(F))$ is compact and by Lemma \ref{tamagawabg}, one has that $\tau_{BG}= \frac{\# G^*(F)}{\# \Sh^2(G)},$ where 
\begin{align*}
\tau_{BG}=\bigg(\mu/\bigg(\frac{\# \Sh^1(G)\coun_{i(BG(F))}}{\# G(F)}\bigg)\bigg)\big(BG(\AAA)/(i(BG(F)))\big).
\end{align*}
Thus $$(\mu/\coun_{i(BG(F))})(BG(\AAA)/i(BG(F)))=\frac{\# G^*(F)\cdot \#\Sh^1(G)}{\# \Sh^2(G)\cdot \# G(F)}.$$
It follows that the dual measure on the discrete group $$i(BG^*(F))=i(BG(F))^{\perp}=(BG(\AAA)/i(BG(F)))^*$$ is given by$$d\chi=\frac{\# \Sh^2(G)\cdot \# G(F)}{\# G^*(F)\cdot \#\Sh^1(G)}\coun_{i(BG^*(F))}.$$
%We apply the Poisson formula .
The necessary conditions of \cite[Chapter II, \S1, \no 8, Proposition 8]{TSpectrale} have been verified  in Lemma~\ref{condofpoiss} (we set $g=1$) and applying it gives:
\begin{align*}
Z(s)&=\frac{1}{\# G(F)}\sum_{x\in BG(F)}H(x)^{-s}\\
&=\frac{\#\Sh^1(G)}{\# G(F)}\sum_{x\in i(BG(F))}H(x)^{-s}\\
&=\frac{\# \Sh^1(G)\cdot\# \Sh^2(G)\cdot \# G(F)}{\# G(F)\cdot\# G^*(F)\cdot \# \Sh^1(G)}\sum_{\chi\in i(BG^*(F))}\wH(s,\chi)\\
&=\frac{1}{\tau_{BG}}\sum_{\chi\in (i(BG^*(F))}\wH(s,\chi),
\end{align*}
whenever every sum converges. It follows from Proposition \ref{convz} that the sums $\sum_{x\in BG(F)}H(x)^{-s}$ and $\sum_{x\in i(BG(F))}H(x)^{-s}$ converge in the domain $\Re(s)>~\eta$. By Corollary \ref{trivialcharacter}, for every $\chi\in BG(\AAA)^*$, the integral defining the Fourier transform $\wH(s,\chi)$ converges in the domain $\Re(s)>1$. Moreover, by Lemma \ref{importantchar}, one has for every $\Re(s)>1$ that $\wH(s,\chi)=0$ unless $\chi$ is an element of the finite set $(\Xi_H)^{\perp}.$ Hence, the sum $\sum_{\chi\in (i(BG^*(F))}\wH(s,\chi)$ reduces to the finite sum $\sum_{\chi\in\Xi_H^{\perp}}\wH(s,\chi).$ We deduce that $$Z(s)=\frac{1}{\tau_{BG}}\sum_{\chi\in \Xi_H^\perp}\wH(s,\chi)$$ for $\Re(s)>\max(1,\eta)$ as claimed.
\end{proof}
\begin{thm}\label{residueofz}
Let $\lambda>1$ be as in Lemma \ref{ftlfun}. The function $s\mapsto~Z(s)$ extends to a meromorphic function in the domain $\Re(s)>\lambda^{-1}.$ The meromorphic function $Z$ admits a pole of order $b(c)$ at $s=1$. One has that \begin{equation*}\lim_{s\to 1}(s-1)^{b(c)}Z(s)=\frac{\tau_H }{\tau_{BG}},\end{equation*}
where $$\tau_H:=\#\Xi_H^\perp\cdot\lim_{s\to 1^+}(s-1)^{b(c)}\int_{\#\Xi_H^\perp}H^{-s}\mu>0.$$
Moreover, if~$F$ is a number field, the function~$Z$ is holomorphic in the domain $\{\Re(s)\geq 1\}-\{1\}$ and if~$F$ is a function field, then the function~$Z$ is holomorphic in the domain $$\{\Re(s)\geq 1-\epsilon\}-\{1+2ki\pi\log(q)^{-1}|k\in\ZZ\}$$ for some $\epsilon>0$.
\end{thm}
\begin{proof}
For every $\chi\in \Xi_H^{\perp}$, by Part (1) of Corollary \ref{trivialcharacter}, one has that $s\mapsto\wH(s,\chi)$ extends to a meromorphic function in the domain $\Re(s)>\lambda^{-1}$ with a possible pole at $s=1$ of order at most $b(c)$. It follows from this fact and Proposition \ref{poisson} that the function $$s\mapsto Z(s)=\sum _{\chi\in \Xi_H^\perp}\wH(s,\chi)$$ extends to a meromorphic function in the domain $\Re(s)>\lambda^{-1}$ with a possible pole at $s=1$ of order at most $b(c)$.

Now, we are going to prove that the function $Z$ has a pole at $s=1$ which is of order $b(c)$. We will apply Poisson formula to the function $$BG(\AAA)^*\to\CC,\hspace{1cm}\chi\mapsto \frac{1}{\tau_{BG}}\wH(s,\chi)$$ when $\Re(s)>\max(\eta,1)$, for the closed inclusion of locally compact abelian groups $$\Xi_H^\perp\subset BG(\AAA)^*,$$ where the finite group $\Xi_H^\perp$ is endowed with the counting measure and the group $BG(\AAA)^*$ is endowed with the measure $\mu^*$. We first calculate the dual measure of the measure $(\mu/\coun_{\Xi_H^\perp})^*$ on $$(BG(\AAA)^*/\Xi_H^\perp)^*=\Xi_H.$$ By Lemma \ref{importantchar}, the group $\Xi_H$ identifies with the kernel of the surjective homomorphism $$BG(\AAA)=(BG(\AAA)^*)^*\to (\Xi_H^{\perp})^*,\hspace{1cm}x\mapsto x|_{\Xi_H^\perp}$$ and is an open and closed subgroup of the index $\# (\Xi_H^\perp)^*=\#\Xi_H^\perp$ in $BG(\AAA)$. The dual measure of the measure $\coun_{\Xi_H^\perp}$ on the dual group $(\Xi^\perp_H)^*$ is given by $\frac{1}{|\Xi_H^\perp|}\cdot~\coun_{(\Xi^\perp_H)^*},$ thus we have an equality of the Haar measures on $(BG(\AAA)^*/\Xi_H^\perp)^*=\Xi_H$: $$(\mu^*/\coun_{\Xi_H^\perp})^*=\#\Xi_H^\perp\cdot \mu|_{\Xi_H}.$$ 
Note that for $x\in (BG(\AAA)^*)^*=BG(\AAA),$ by the Fourier inversion formula (\ref{fourierinversion}), one has for $\Re(s)>\eta$ that
\begin{align*}
\int_{BG(\AAA)^*}\wH(s,\chi)\overline{x(\chi)}d\mu^*(\chi)\hskip-1cm&\\&=\int_{BG(\AAA)^*}\bigg(\int_{BG(\AAA)}H(y)^{-s}\overline{\chi(y)}d\mu(y) \bigg)\overline{x(\chi)}d\mu^*(\chi)\\
&= H(x)^{-s}.
\end{align*}
Now, for $\Re(s)>\max(\eta,1),$ the Poisson formula and Proposition \ref{poisson} give $$Z(s)=\sum_{\chi\in\Xi_H^\perp}\frac{1}{\tau_{BG}}\wH(s,\chi)=\frac{\#\Xi_H^\perp}{\tau_{BG}}\cdot\int_{\Xi_H}H^{-s}\mu.$$
%The equality is valid, as an equality of meromorphic functions in the domain $\Re(s)>\lambda^{-1}.$ Moreover, the functions have a possible pole at $s=1$ which is of order at most~$b(c)$. To prove that the pole of $Z$ exists and that it is of the order $b(c)$, it suffices to prove that the limit 
%$$\lim_{s\to 1^+}(s-1)^{b(c)}\int_{\Xi_H}H^{-s}\mu$$exists and is strictly bigger than $0$. Part (2) of Corollary \ref{trivialcharacter} gives that $s\mapsto\int_{BG(\AAA)}H^{-s}\mu $ is a meromorphic function in the domain $\Re(s)>\lambda^{-1}$  with a pole at $s=1$ of order $b(c)$. 
%Let $x_1\doots x_{|\Xi_H^{\perp}|}\in BG(\AAA)$ be a set of representatives modulo $\Xi_{H}$. Using Lemma \ref{codx}, we obtain for $s>1$ that \begin{align*}
%\int_{BG(\AAA)}H^{-s}\mu&=\sum_{i=1\doots |\Xi_H^\perp|}\int_{\Xi_H^\perp}H^{-s}(x_iy)^{-s}d\mu(y)\\
%&\leq \sum_{i=1\doots |\Xi_H^\perp|}C(x_i)^{-s}\int_{\Xi_H}H^{-s}\mu
%\end{align*} for certain $C(x_i)>0$.
 Corollary \ref{trivialcharacter} gives $$0<\lim_{s\to 1^+} (s-1)^{b(c)}\int_{\Xi_H}H^{-s}\mu.$$ %It follows that for $s>1,$ one has that
%\begin{align*}
%0&<\lim_{s\to 1^+} (s-1)^{b(c)}\int_{BG(\AAA)}H^{-s}\mu\\
%&\leq\lim_{s\to 1^+}(s-1)^{b(c)}\sum_{i=1\doots {|\Xi^\perp_H|}}C(x_i)^{-s}\int_{\Xi_H}H^{-s}\mu\\
%&\leq \bigg(\sum_{i=1\doots {|\Xi^\perp_H|}}C(x_i)^{-1}\bigg)\lim_{s\to 1^+}(s-1)^{b(c)}\int_{\Xi_H}H^{-s}\mu,
%\end{align*} and
%Lemma~\ref{intfinind} gives that $\lim_{s\to 1^+}(s-1)^{b(c)}\int_{\Xi_H}H^{-s}\mu>0.$ 
and we deduce that the function $s\mapsto Z(s)$ admits a pole at $s=1$ of order exactly $b(c).$ The principal value at this pole is $$\lim_{s\to 1^+}(s-1)^{b(c)}Z(s)=\frac{ \#\Xi_H^\perp}{\tau_{BG}}\lim_{s\to 1^+}(s-1)^{b(c)}\int_{\Xi_H}H^{-s}\mu.$$ 
If~$F$ is a number field, then it follows from Part (1) of Corollary \ref{trivialcharacter} that $s\mapsto Z(s)=\sum_{\chi\in\Xi_H^{\perp}}\wH(s,\chi)$ is holomorphic in the domain $\{\Re(s)\geq 1\}-\{1\}$. If~$F$ is a function field, then it follows from Part (1) of Corollary \ref{trivialcharacter} that $s\mapsto Z(s)=\sum_{\chi\in\Xi_H^{\perp}}\wH(s,\chi)$ is holomorphic in the domain $\{\Re(s)\geq 1\}-\{1+2ki\pi\log(q)^{-1}|k\in\ZZ\}$.  The proof of the theorem is completed.
\end{proof}
%When $F$ is a number field, applying a Tauberian result \cite[Theorem 3.3.2]{aniso} gives us following theorem.
%\begin{lem}\label{verticalz} Suppose $F$ is a number field. There exists $\epsilon >0$ such that the function $s\mapsto (s-1)^{b(c)}Z(s)$ is holomorphic in the domain $\Re(s)>1-\epsilon$. Moreover, there exists $\delta>0$ such that for every compact $\mathcal K\subset ]1-\epsilon, +\infty[$ one has that $$\bigg|\frac{(s-1)^{b(c)}Z(s)}{s^{b(c)}}\bigg|\leq C(\mathcal K)(1+|\Im(s)|)^\delta.$$
%\end{lem}
%\begin{thm}\label{main}
%Suppose that $F$ is a number field. One has that 
%\begin{multline*}
%N_{H}(B)=\# \{x\in BG(F)| H(x)\leq B\}\\\sim_{B\to\infty}\frac{\#\Sh^2(G)}{\# G^*(F)}\frac{\tau_H}{(b(c)-1)!}\cdot B\log(B)^{b(c)-1},
%\end{multline*}
%where $$\tau_H= \lim_{s\to 1^+}(s-1)^{b(c)}\cdot\#\Xi_H^\perp\cdot\int_{\Xi_H}H^{-s}\mu.$$
%\end{thm}
%\begin{proof}
%The conditions of \cite[Theorem A.1]{Igusa} have been verified in Theorem \ref{residueofz} and in Lemma \ref{verticalz}. The principal value has been calculated in Theorem \ref{residueofz}. The result is now a direct application of \cite[Theorem~A.1]{Igusa}.
%\end{proof}
%\subsubsection{}In this paragraph, we suppose that~$F$ is a function field and we establish asymptotic bounds for the number of torsors of bounded height.
%\begin{lem}
%\begin{enumerate}
%\item There exists a rational function~$f$ such that~$f$ has a pole at~$s=1$ and no other poles in the domain $\Re(s)\geq 1$ and a bounded holomorphic function~$\beta$ defined in the domain $\Re(s)\geq 1$, such that~$\wH(s,1)=f(q^{-s})\beta(q^{-s}).$ 
%\item The meromorphic function $s\mapsto\wH(s,1)$ has poles at points $1+2ik\pi$ for~$k\in\ZZ$ or order~$b(c)$ and no other poles in the domain~$\Re(s)>1-\epsilon$ for some~$\epsilon>0$
%\end{enumerate}
%\end{lem}
%\begin{proof}
%We have that $$H^{\Sigma}(s,1)=\prod_{v\in M_F^0-\Sigma}\wH(s,1) $$
%For $v\not\in\Sigma$, one has that \begin{align*}
%\wH_v(s,1)&=\int_{BG(F_v)}H^{-s}_v\mu_v\\
%&=\frac{1}{\# G(\Fv)}\sum_{x\in BG(F_v)}H(x)^{-s}\\
%&=1+ \# (G_*)_1(\Fv)\cdot q_v^{-s} +\sum_{r\in c(G_*(\oF))-\{0,1\}}{\# (G_*)_r(\Fv)} \cdot q_v^{-rs}.
%\end{align*}
%\end{proof}
%We need the following Tauberian result.
%\begin{lem}
%Let~$\lambda, a >0$. Consider the Dirichlet series $$D(s):=\sum_{n=0}^{\infty}a_n\exp(-sn\lambda).$$We suppose that~$D(s)$ converges absolutely to a holomorphic (necessary $2i\pi\lambda^{-1}$-periodic) function in the domain $\Re(s)>a,$ that it extends to a meromorphic $2i\pi\lambda^{-1}$-periodic function in the domain $\Re(s)> a-\epsilon,$ for some $\epsilon>0$, with the poles at the points $\{a+2ik\pi\lambda^{-1}|k\in\ZZ\}$ which are all of order~$b>0$, and no other poles in the domain $\Re(s)>a-\epsilon$. One has that
%\begin{equation*}
%\sum_{n=0}^Ma_n\sim_{M\to\infty} \bigg(\frac{\exp(ab\lambda)\big(\lim_{s\to 1}(s-1)^r\cdot D(s)\big)}{(b-1)!\cdot (1-\exp(-a\lambda))}\bigg) \exp(\lambda a M)\big(M a\big).
%\end{equation*}
%\end{lem}
%\begin{proof}Consider the variable $t=\exp(-s\lambda)$. The power series $$W(t):=\sum_{n=0}^{\infty}a_nt^{n}$$ converges absolutely to a holomorphic function whenever $|t|<R:=\exp(-a\lambda)$. Moreover, it extends to a meromorphic function in the domain $|t|<R+\epsilon$, with the only pole at the point~$R$, which is of order~$b$ and one has that $$\lim_{t\to R}(t-R)^{b}W(t)= $$ Now, by Part (c) of \cite[Theorem A.5]{lagemanartinschreier}, one has that STOP HERE THE PRINCIPAL VALUE IS  A BIT COMPLICATED, must take into account all Laurent coefficients
%\begin{align*}\sum_{n=0}^ma_m&\sim_{m\to\infty}C\cdot R^{-m}\log(R^{-m})^{b-1}\\
%&\sim_{m\to\infty}C\cdot \exp(ma\lambda)(ma\lambda)^{b-1}
%\end{align*} for some 
%\end{proof}
%\begin{thm}\label{mainfunction}
%One has that $$\#\{x\in BG(F)|H(x)\leq B\}\asymp_{B\to\infty}B\log(B)^{b(c)-1}. $$
%\end{thm}
\subsection{Equidistribution}\label{secequid}  We keep notations from previous subsection. Theorem \ref{residueofz} allows modifications of heights at finitely many places. By using \cite[Proposition 2.10]{Igusa}, this will give us that $G$-torsors are equidistributed in a subspace of the infinite product $\prod_{v\in M_F}BG(F_v)$ for a certain measure related to the height $H$.
%We start off by defining the measure. Let us denote by $\mu_{\infty}$ the counting measure on $BG(F)$ multiplied by the factor $\frac{1}{\# G(F))}$. Let $K$ be an open and compact such that $i(BG(F))\cap K=\{(1)_v\}$. One has that $$i(BG(F))\times K\to \Xi_K:=i(BG(F))K\hspace{1cm} (x,k)\mapsto xk$$ is an isomorphism which takes the measure $\mu_{\infty}\times K$ to $\mu|_{\Xi_K}$. EXPLAIN?
 By abuse of notation, if $f:\prod_{\vMF}BG(F_v)\to \CC$ is a function, we may write also~$f$ for the composite map $f\circ i:BG(F)\to\CC$ and the composite map $BG(\AAA)\to \prod_{\vMF}BG(F_v)\xrightarrow{f}\CC$. We note that the inclusion $BG(\AAA)\to \prod_{\vMF}BG(F_v)$ is continuous. The functions of form~$\otimes_v f_v$ with~$f_v\equiv 1$ for almost all~$v$ will be called {\it elementary} functions.
\subsubsection{}In the following lemma we will define a measure on the product~$\prod_{\vMF}BG(F_v)$ for which rational points equidistribute. More practical formula will be provided in Proposition \ref{genfz}. 
\begin{lem}\label{existmeas}
For every continuous $f:\prod_{\vMF}BG(F_v)\to~\CC$, the series $$\frac{1}{\# G(F)}\sum_{x\in BG(F)}f(x)H(x)^{-s}$$converges absolutely to a holomorphic function in the domain $\Re(s)>1$ and the limit $$\omega_H(f):=\frac{\tau_{BG} }{\# G(F)}\lim_{s\to 1^+}(s-1)^{b(c)}\sum_{x\in BG(F)} f(x)H(x)^{-s}$$exists. Moreover, the map $f\mapsto\omega_H(f)$ defines a Radon measure on on the space $\prod_{\vMF}BG(F_v)$.
\end{lem}
\begin{proof}It follows from Theorem \ref{residueofz} and \cite[Section 2.3, Proposition~7]{coursdarithmetique} that the abscissa of absolute convergence of the Dirichlet series $Z(s)$ is $\Re(s)=1$.  Now, by boundedness of $f$, the first claim follows. Let us prove the existence of the limit.
\begin{enumerate}[label=(\alph*)]
\item We first suppose that $f=\otimes_v f_v$ is elementary, with $f_v$ strictly positive functions. %:BG(F_v)\to~\RR_{> 0}$ such that $f_v\equiv 1$ for almost all~$v$. 
One has that 
\begin{align*}
%\lim_{s\to 1^+}\frac{(s-1)^{b(c)}}{\# G(F)}\sum_{x\in BG(F)} f(x)H(x)^{-s}\hskip-3cm&\\\#\Sh^1(G)
\frac{\tau_{BG} }{\# G(F)}\lim_{s\to 1^+}{(s-1)^{b(c)}}\sum_{x\in BG(F)} \big(f(x)^{-1}\cdot H(x)\big)^{-s}%\\
%&=\lim_{s\to 1^+}\frac{(s-1)^{b(c)}}{\# G(F)}\sum_{x\in BG(F)} \big(H((f_v\cdot c_v)_v)(x)\big)^{-s},
\end{align*}exists by Theorem \ref{residueofz}. By the compactness of $\prod_{\vMF}BG(F_v)$, one has that $f^{s}\to f$ uniformly when $s\to 1^+$. Thus, one has that $$\frac{\tau_{BG} }{\# G(F)}\lim_{s\to 1^+}{(s-1)^{b(c)}}\sum_{x\in BG(F)} \big(f(x)^{s}-f(x)\big)\cdot H(x)^{-s} =0.$$By subtracting the later limit from the former, we deduce that $\omega_H(f)$ exists. 
\item Suppose that $f$ is a finite sum of functions $\otimes_v f_v$ as in (a). Then $\omega_H(f)$ exists. Moreover, if $\lambda\in\RR$ is scalar, then $\omega_H(\lambda\cdot(\otimes_v f_v))=\lambda\cdot \omega_H(\otimes_v f_v)$ exists.
\item Consider the family of functions which are of form $g(f_1\doots f_n)$, where $n\geq 1$ is an integer, $g\in\RR[X_1\doots X_n]$ is a polynomial and $f_i$ are as in (a). The functions from (a) are stable for multiplications, hence, this family is given by the linear combinations of elementary functions. 
\item The family of elementary functions contains constants and separates the points of $\prod_{\vMF}BG(F_v)$. By Stone-Weierstrass theorem \cite[Chapter X, \S 4, \no 2, Theorem 3]{TopologieGd} and by (c), one has that any continuous $f:\prod_{\vMF}BG(F_v)\to\RR$ is a uniform limit of functions which are linear combinations of functions of form $\otimes_v f_v$. Thus, the limit $\omega_H(f)$ exists. 
\item For a continuous complex valued function $f=\Re(f)+i\cdot\Im(f)$, one has that $\omega_H(f)$ exists because $\omega_H(\Re(f))$ and $\omega_H(\Im(f))$ exist. \end{enumerate}
Let us prove that~$\omega_H$ is a measure. If~$f$ is a non-negative function, then it is a limit of finite sums of non-negative functions~$\otimes_v f_v$, hence~$\omega_H(f)$ is non-negative. By \cite[Chapter III, \S 1, \no 6, Theorem 1]{Integrationj}, one has that~$\omega_H$ is a Radon measure. 
\end{proof}
%\begin{multline*}
%\lim_{s\to 1^+}(s-1)^{b(c)}\frac{1}{\# G(F)}\sum_{x\in BG(F)} f(i(x))H(x)^{-s}\\=\lim_{s\to 1^+}(s-1)^{b(c)}|\Xi_H^{\perp}|\int_{\Xi_H}f\cdot H(x)^{-s}\mu
%\end{multline*}
%exists and is strictly positive. Now, by \cite[Chapter III, \S 1, \no 6, Theorem~1]{Integrationj}, the map $$f\mapsto \lim_{s\to 1^+}\frac{(s-1)^{b(c)}}{\# G(F)} \sum_{x\in BG(F)} f(x)H(x)^{-s}$$ defines a measure $\omega_H$ on $\prod_{\vMF}BG(F_v)$. 
%\begin{mydef}
%We define a functional $\omega_H:\mathcal C^0_{K_H}(\prod_{\vMF}BG(F_v), \CC)\to\CC$ by
%\begin{align*}
%\omega_H:f&\mapsto \lim _{s\to 1^+}(s-1)^{b(c)}\sum _{\chi\in \Xi_H^{\perp}}\int_{BG(\AAA)}H^{-s}f\chi\mu\\
%%&=\lim_{s\to 1^+}\int_{\Xi_H}H^{-s}f\mu.
%\end{align*}
%where $\mathcal C^0(-,-)$ stands for the continuous functions which are $K_H$-invariant.
%\end{mydef}
%\begin{lem}
%One has that $$\overline{i(BG(F))}=j(\Xi_H).$$
%\end{lem}
%\begin{proof}
%\end{proof}
\begin{lem}\label{supportofmeasureo}
One has that $$\supp(\omega_H)=\overline{i(BG(F))}$$ and that $$ \tau_H=\omega_H(\overline{i(BG(F))}).$$
\end{lem}
\begin{proof}
We denote by~$\mathbf 1_X$ the characteristic function of~$X$. If we denote by $U=\prod_{\vMF}BG(F_v)-i(\overline{BG(F)})$, then \begin{align*}\omega_H(U)&=\omega_H(\mathbf 1_U)\\&=\frac{\tau_{BG}}{\# G(F)}\lim_{s\to 1^+}{(s-1)^{b(c)}}\sum_{x\in BG(F)}\mathbf{1}_U(x)H(x)^{-s}\\&=0,
\end{align*}
thus we obtain that $$\supp(\omega_H)\subset\overline{i(BG(F))}.$$ 
Let us prove the reverse inclusion. Let $(y_v)_v\in i(BG(F)),$ let~$\Sigma_y\supset\Sigma_H$ be a finite set of places of~$F$ containing all places~$v$ for which $y_v\not\in BG(\Ov)$ and let $U_y\subset\prod_{\vMF}BG(F_v)$ be an open neighbourhood of~$(y_v)_v$ of the form $$\prod_{v\in\Sigma_y}\{y_v\}\times \prod_{v\in M_F-\Sigma_x}BG(F_v) .$$   We will prove that $\omega_H(U_y)>0$.  Let us write $$K=\prod_{v\in\Sigma_y}\{1\}\times\prod_{v\in M_F-\Sigma_y}BG(\Ov)$$ and $\Xi_K=i(BG(F))K\subset BG(\AAA)$, one has that the index~$(BG(\AAA):\Xi_K)$ is finite.  We denote by $j:BG(\AAA)\to \prod_{\vMF}BG(F_v)$ the inclusion, which is a continuous homomorphism. One has that 
\begin{align*}j^{-1}(U_y)&=j^{-1}\bigg(\prod_{v\in\Sigma_y}\{y_v\}\times \prod_{M_F-\Sigma_y}BG(F_v)\bigg)\\&=\yyy j^{-1}\bigg(\prod_{v\in\Sigma_y}\{1\}\times\prod_{v\in M_F-\Sigma_y}BG(F_v)\bigg)\\
&=\yyy V,
\end{align*}
where $\yyy=((y_v)_{v\in \Sigma_y},(1)_{v\in M_F-\Sigma_y})$ and $$V=j^{-1}\bigg(\prod_{v\in\Sigma_y}\{1\}\times\prod_{v\in M_F-\Sigma_y}BG(F_v)\bigg).$$ Note that $\yyy\in\Xi_K$ and that the subgroup $V\subset BG(\AAA)$ is of finite index. %It follows from Part (1) of \ref{importantchar}  
Now, Lemma \ref{intfinind}, the fact that $\sum_{\chi\in\Xi_K^{\perp}}\chi=\#\Xi_K^{\perp}\cdot\mathbf{1}_{\Xi_K}$ and Proposition~\ref{thmofequidistribution} give that
\begin{align*}
\omega_H(U_y)&=\omega_H(\mathbf{1}_{U_y})\\
&=\tau_{BG}\lim_{s\to 1^+}(s-1)^{b(c)}Z^{\mathbf{1}_{U_y}}(s)\\
&=\lim_{s\to 1^+}(s-1)^{b(c)}\sum_{\chi\in\Xi_K^{\perp}}\int_{BG(\AAA)}(\mathbf 1_{j^{-1}(U_y)}\cdot H^{-s})\chi\mu\\
&=\lim_{s\to 1^+}(s-1)^{b(c)}\#\Xi_K^{\perp}\cdot \int_{BG(\AAA)}\mathbf 1_{\Xi_K}\mathbf 1_{j^{-1}(U_y)}H^{-s}\mu\\
&=\#\Xi_K^{\perp}\cdot\lim_{s\to 1^+}(s-1)^{b(c)}\int_{\Xi_K\cap j^{-1}(U_y)}H^{-s}\mu\\
&>0.
\end{align*}
It follows that $\supp(\omega_H)\supset i(BG(F))$ and hence that $$\supp(\omega_H)=\overline{i(BG(F))}.$$
Finally, we have that \begin{align*}\omega_H(\overline{i(BG(F))})&=\omega_H\big(\prod_{\vMF}BG(F_v)\big)\\&=\omega_H(\mathbf{1}_{\prod_v BG(F_v)})\\&= \tau_{BG}\lim_{s\to 1^+}\frac{(s-1)^{b(c)}}{\# G(F)}\sum_{x\in BG(F)}H(x)^{-s}\\&=\tau_{BG}\cdot\lim_{s\to 1^+}(s-1)\cdot Z(s)\\&=\tau_H.
\end{align*} 
%\item One has that $$(\otimes g_v)\cdot H=H((c_vg_v)_v).$$ Let $f:\prod_{v\in M_F}BG(F_v)\to \RR_{\geq 0}$ be a continuous function. Let us first suppose that for every $v$, one has that $g_v\leq 1$.  We have that
%\begin{align*}
%\omega_{H(g_vc_v)_v)}(f)&=\lim_{s\to 1^+}\frac{(s-1)^{b(c)}}{\# G(F)}\sum_{x\in BG(F)}f(x)H(x)^{-s}\prod_v g_v(x)^{-s}\\&\geq \lim_{s\to 1^+}\frac{(s-1)^{b(c)}}{\# G(F)}\sum_{x\in BG(F)}f(x)H(x)^{-s}\prod_v g_v(x)^{-1}\\
%&=\omega_H((\otimes g_v)^{-1}\cdot f).
%\end{align*}
%If $\delta>1$, by taking the limit over $s\in]1,\delta[, $ we deduce that \begin{align*}\omega_{H(g_vc_v)_v)}(f)&\leq \lim_{s\to 1^+}\frac{(s-1)^{b(c)}}{\# G(F)}\sum_{x\in BG(F)}f(x)H(x)^{-s}\prod_v g_v(x)^{-\delta}\\
%&=\omega_H((\otimes g_v)^{-\delta}\cdot f).
%\end{align*}
%By the compactness of $\prod_{\vMF}BG(F_v)$, one has $\lim_{\delta\to 1^+}(\otimes g_v)^{-\delta}=(\otimes g_v)^{-1}$ and by the fact that $\omega_H$ is a measure we have that $$\lim_{\delta\to 1^+}(\omega_H((\otimes g_v)^{-\delta}\cdot f))=\omega_H((\otimes g_v)^{-1}\cdot f),$$ We deduce that \begin{equation}\label{jedn}\omega_{H((c_vg_v)_v)}(f)=\omega_{H}((\otimes g_v)^{-1}\cdot f).
%\end{equation}Suppose that for every $v$ one has that $g_v\geq 1$. Then the reasoning as above is valid, but the sign $\geq$ should be replaced by $\leq$ and vice versa. We deduce that for such $g_v$, one has that (\ref{jedn}) is valid. For a general $(g_v)_v$, we set $g_v'=\min (g_v,1)$ and $g_v''=\max (g_v,1)$ so that $g_v=g_v''\cdot g_v'$. Then it follows from above that \begin{align*}\begin{split}\label{dvaa}\omega_{H((c_vg_v)_v)}(f)= \omega_{H((c_vg_v')_v)}((\otimes g_v'')^{-1}\cdot f)&=\omega_H((\otimes g_v'')^{-1}\cdot(\otimes g_v')^{-1}\cdot f)\\&=\omega_H((\otimes g_v)^{-1}\cdot f)
%\end{split}
%\end{align*}
%for every $f:\prod_{\vMF}BG(F_v)\to\RR_{\geq 0}$. Every continuous function $f:\prod_{\vMF}BG(F_v)\to\RR$ writes as a difference of two positive valued functions and it follows that (\ref{dvaa}) is valid for every continuous real valued function~$f$. Finally by treating separately the real and imaginary part of~$f$, we deduce that (\ref{dvaa}) is valid for every continuous complex function~$f$. The statement follows. 
%\end{enumerate}
\end{proof}
\begin{rem}
\normalfont Let~$F$ be a number field. Harari proves in \cite[Page 522]{haraprox} that {\it Brauer-Manin obstruction is the only obstruction to the weak approximation for~$G$}. Let us explain what he means by this. We assume that~$G$ is embedded as a closed subgroup scheme of $\GL_n$ for certain $n\geq 1$ and let us set $X=\GL_n/G$ to be the quotient variety. Then the closure of the image of $X(F)$ for the diagonal map to $\prod_{\vMF}X(F_v)$ is the Brauer-Manin set $\big(\prod_{\vMF}X(F_v)\big)^{\Br}$, i.e. the subset of $\prod_{\vMF}X(F_v)$ where the Brauer-Manin obstruction vanishes. Moreover, he explains that one could have chosen other embedding in $\GL_n$ or in $\SL_n$ and the same property would be true. Demeio told us that one can define Brauer-Manin set $\big(\prod_{\vMF}BG(F_v)\big)^{\Br}$ as the image of $\big(\prod_{\vMF}X(F_v)\big)^{\Br}$ for the induced map $\prod_{\vMF}X(\Fv)\to\prod_{\vMF}BG(F_v)$ (although there is a better definition in terms of certain pairings).  Then it is not hard to verify that the fact that Brauer-Manin obstruction is the only obstruction to the weak approximation for~$G$ translates into  $$\overline{i(BG(F))}=\bigg(\prod_{\vMF}BG(F_v)\bigg)^{\Br}.$$ 
We may compare this with Colliot-Th\' el\`ene's conjecture \cite[Conjecture 14.1.2]{brauergrothendieck} that the Brauer-Manin obstruction is the only obstruction to the weak approximation for rationally connected varieties.
\end{rem}
% It is a conjecture of Colliot-Th\'el\`ene \cite[Conjecture14.1.2]{brauergrothendieck} that the analogue of the set $\overline{i(BG(F))}$ for rationally connected varieties coincides with the {\it Brauer-Manin set} of a variety. We explain in this remark that the corresponding statement is true in our situation. Let us choose an embedding $G\hookrightarrow \GL_n$ and set $X=\GL_n/G$ be the quotient variety. One has a $1$-equivalence of stacks $X/\GL_n\xrightarrow{\sim} BG.$ For $\vMF$, we endow the sets $X(F_v)$ with usual topologies. By \cite[Section 2]{Cesnavicius}, the maps $X(F_v)\to BG(F_v)$ are continuous (where $BG(F_v)$ are discrete). We denote by $$\bigg(\prod_{\vMF}X(F_v)\bigg)^{\Br}\subset\prod_{\vMF}X(F_v)$$ the Brauer-Manin set of~$X$ and by $$\bigg(\prod_{\vMF}BG(F_v)\bigg)^{\Br}\subset\prod_{\vMF}BG(F_v)$$the image of $\big(\prod_{\vMF}X(F_v)\big)^{\Br}$ for the canonical map $$\prod_{\vMF}X(F_v)\to\prod_{\vMF}BG(F_v).$$  By the fact \cite{kuc} that~$X$ is stably $F$-birational to a torus, it follows from \cite{kuc} that the {\it Brauer-Manin obstruction is the only obstruction to the weak approximation of}~$X$, that is $\overline{X(F)}=\big(\prod_{\vMF}X(F_v)\big)^{\Br}$, where~$X(F)$ is embedded diagonally in $\prod_{\vMF}X(F_v)$. 
%For every~$v$, the diagram 
%\[\begin{tikzcd}
%	{X(F)} & {X(F_v)} \\
%	{BG(F)} & {BG(F_v)}
%	\arrow[from=1-1, to=1-2]
%	\arrow[from=1-1, to=2-1]
%	\arrow[from=2-1, to=2-2]
%	\arrow[from=1-2, to=2-2]
%\end{tikzcd}\]
%is commutative, and we deduce that $i(BG(F))$ is the image of $X(F)$ for the composite map $X(F)\to \prod_{\vMF}X(F_v)\to\prod_{\vMF}BG(F_v)$. Finally, the map $\prod_{\vMF}X(F_v)\to\prod_{\vMF}BG(F_v)$ is continuous and we deduce that $$\overline{i(BG(F))}=\bigg(\prod_{\vMF}BG(F_v)\bigg)^{\Br}.$$ 
%\end{rem}
%THIS MAY BE NEEDEDThe following is valid:
%\begin{prop}
%One has that the index $$\bigg(\big(\prod_{\vMF}BG(F_v)\big):\overline{i(BG(F))}\bigg)$$ is finite.
%\end{prop}
%\begin{proof}
%Let $F(G^*)$ be a finite Galois extension of $F$, such that $(G^*)_{F(G^*)}$ is a constant group. Let $\Sigma'$ be the finite set of the finite places $v$ such that $\gcd(\# G, \#\kappa(v))\neq 1$ or such that the Galois group $\Gamma_{F(G^*)_{v'}/F_v}$ is not cyclic for any $v'$ extending $v$. 
%By \cite[Theorem 9.2.3 (vii)]{CohomologyNF}, for every finite subset $T\subset M_F-\Sigma'$, one has that $$BG(F)\to\prod_{v\in T}BG(F_v)$$is surjective. It follows that $$i^{M_F-\Sigma'}:BG(F)\to\prod_{v\in M_F-\Sigma'}BG(F_v)$$ is of dense image. %Let us denote by $BG(F)_{\Sigma'}$ the kernel of $$i^{\Sigma'}:BG(F)\to\prod_{v\in\Sigma'}BG(F_v).$$ 
% Let $x\in\prod_{\vMF-\Sigma'}BG(F_v)$. For an open neighbourhood~$U_{n}$ of~$x$, there exists an element $b_n\in BG(F)$ such that $i^{M_F-\Sigma'}(b_n)\in~U_n$. Moreover, the kernel of the homomorphism $i^{\Sigma'}:BG(F)\to\prod_{v\in \Sigma'}BG(F_v)$ is of the finite index in $BG(F)$, hence, there exists a subsequence $U_{\phi(n)}$ of~$U_n$ such that every two $b_{\phi(n)}$ are congruent to certain $b\in BG(F)$ modulo $\ker(i^{\Sigma'})$. We deduce that $i^{M_F-\Sigma'}(b^{-1})\cdot x$ is a limit point of $i^{M_F-\Sigma'}(a_n)$, where $a_n:=b_{\phi(n)}^{-1}x\in\ker (i^{\Sigma'}).$ It follows that $\overline{i^{M_F-\Sigma'}(\ker(i^{\Sigma'}))}$ is of finite index in $\prod_{\vMF-\Sigma'}BG(F_v).$ We deduce that \begin{align*}\bigg(\big(\prod_{\vMF}BG(F_v)\big):\overline{i(BG(F))}\bigg)\hskip-4cm&\\&\leq \bigg(\big(\prod_{\vMF}BG(F_v)\big):\big(\prod_{v\in\Sigma'}\{1\} \times \overline {i^{M_F-\Sigma'}(BG(F))}\big)\bigg)\\
%&\leq \bigg(\big(\prod_{\vMF}BG(F_v)\big):\big(\prod_{v\in\Sigma'}\{1\} \times \overline {i^{M_F-\Sigma'}(\ker(i^{\Sigma'}))}\big)\bigg)\\
%&\leq \big(\prod_{v}\#BG(F_v)\big)\cdot \big(\big(\prod_{v\in M_F-\Sigma'}BG(F_v)\big):\overline{i^{M_F-\Sigma'}(\ker(i^{\Sigma'}))}\big)
% \end{align*} is finite.
%\end{proof}
\subsubsection{}We modify the Dirichlet series by multiplying by the value of~$f$. We calculate the limit when the new series is multiplied by $(s-1)^{b(c)}$ and $s\to 1$.
\begin{prop}\label{genfz}
Let $f:\prod_{\vMF}BG(F_v)\to \RR_{\geq 0}$ be a continuous function. We set $$Z^f(s):=\frac{1}{\# G(F)}\sum_{x\in BG(F)}f(x)^sH(x)^{-s}.$$ 
\begin{enumerate}
\item For every~$\epsilon>0$, the series $Z^f(s)$ converges absolutely and uniformly to a holomorphic function in the domain $\Re(s)>1+\epsilon$. One has that $$\lim_{s\to 1^+}(s-1)^{b(c)}Z^f(s)=\frac{\omega_H(f)}{\tau_{BG}}.$$
\item Suppose that $f=\sum_{i=1}^n\lambda_if_i$, with $\lambda_i\in\RR$ and~$f_i$ elementary. Let $\Sigma$ be the set of all places containing $\Sigma_G\cup M_F^\infty$ such that for every $v\in M_F-\Sigma$ one has that $(f_i)_v\equiv 1$ and that~$H_v$ is~$BG(\Ov)$-invariant. One has that \begin{align*}
Z^f(s)&=\frac{1}{\tau_{BG}}\sum_{\chi\in \Xi^{\perp}_K}\int_{BG(\AAA)}f^sH^{-s}\chi\mu\\&=\frac{1}{\tau_{BG}}\sum_{\chi\in \Xi_K^{\perp}}\bigg(\int_{\prod_{v\in \Sigma}BG(F_v)}f_{\Sigma}^sH_{\Sigma}^{-s}\chi_{\Sigma}\mu_{\Sigma}\bigg) \cdot \prod_{v\in M_F-\Sigma}\wH_v(s,\chi_v),\end{align*}
where~$K=\prod_{\vMFz-\Sigma}BG(\Ov),$ and $\Xi_K^{\perp}$ denotes the set of characters in $BG^*(F)$ which vanish on~$K$, while the subscript~$\Sigma$ denotes the restriction of a function or a measure to $\prod_{v\in\Sigma}BG(F_v)$. If~$f\neq 0$, the function $s\mapsto Z^f(s)$ extends to a meromorphic function in the domain $\Re(s)\geq 1$ and admits a pole at $s=1$ of order~$b(c)$. Moreover, if~$F$ is a number field, then $s\mapsto Z^f(s)$ is holomorphic in the domain $\{\Re(s)\geq 1\}-\{1\}$ and if~$F$ is a function field, then $s\mapsto Z^f(s)$ is holomorphic in the domain $\{\Re(s)\geq 1-\epsilon\}-\{1+2ki\pi\log(q)^{-1}|k\in\ZZ\}$ for some~$\epsilon>0$.
\end{enumerate}
\end{prop}
\begin{proof}
\begin{enumerate}
\item   It follows from Theorem \ref{residueofz} and \cite[Section 2.3, Proposition~7]{coursdarithmetique} that the Dirichlet series $Z(s)$ converges absolutely and uniformly in the domain $\Re(s)\geq 1+\epsilon$. The claim on the convergence and holomorphicity of $Z^f(s)$ in the domain $\Re(s)>1+\epsilon$ now follows by the fact that~$f$ is bounded. By the same theorem, one has $\lim_{s\to 1^+}(s-1)^{b(c)}Z(s)=\tau_H$ and by the compactness of $\prod_{\vMF}BG(F_v)$, one has that $f^s\to f$ uniformly. We deduce that $$\lim_{s\to 1^+}\frac{(s-1)^{b(c)}}{\# G(F)}\sum_{x\in BG(F)}(f(x)^s-f(x))H(x)^{-s}=0.$$ %$\varepsilon>0$, one has that Suppose first that $f>0$.
%\begin{align*}
%0\leq&\lim_{s\to 1^+}\frac{(s-1)^{b(c)}}{\# G(F)}\sum_{x\in BG(F)}(f(x)^s-f(x))H(x)^{-s}\\
%&\leq\lim_{s\to 1^+}\frac{(s-1)^{b(c)}}{\# G(F)}\sum_{x\in BG(F)}\varepsilon^sH(x)^{-s}\\
%&\leq \varepsilon \tau_H
%\end{align*}
Thus
\begin{align*}\lim_{s\to 1^+}(s-1)^{b(c)}Z^f(s)&=\lim_{s\to 1^+}\frac{(s-1)^{b(c)}}{\# G(F)}\sum_{x\in BG(F)}f(x)H(x)^{-s}\\&\quad+\lim_{s\to 1^+}\frac{(s-1)^{b(c)}}{\# G(F)}\sum_{x\in BG(F)}(f(x)^s-f(x))H(x)^{-s}\\
&=\frac{\omega_H(f)}{\tau_{BG}}.
\end{align*}
%HERE USEFUL MAYBE Let us now treat the case when there may be some~$x$ such that $f(x)=0$. We write $f=f_1-f_2$, with $f_1\geq f_2>0$. When $s\to 1^+$, one has that $f(x)^s\to f(x)$ and that $\lim_{s\to 1}(f_1(x)^s-f_2(x)^s)\to f_1(x)-f_2(x)=f(x)$ uniformly in~$x$, hence $f_1(x)^s-f_2(x)^s\to f(x)^s$ uniformly in~$x$. We deduce that 
%\begin{align*}\lim_{s\to 1^+}(s-1)^{b(c)}Z^f(s)&=\lim_{s\to 1^+}\frac{(s-1)^{b(c)}}{\# G(F)}\sum_{x\in BG(F)}f(x)^sH(x)^{-s}\\
%&=\frac{(s-1)^{b(c)}}{\# G(F)}\sum_{x\in BG(F)}(f_1(x)^s-f_2(x)^s)H(x)^{-s}%\\&=\omega_H(f_1)-\omega_H(f_2)\\
%&=\omega(f).
%\end{align*}
\item %There exists a finite set of places $\Sigma',$ containing $\Sigma\cup M_F^{\infty}$, such that for every $i$ and every $v\in M_F-\Sigma',$ one has that $(f_i)_v\equiv 1$ and that $H_v$ is $B\GG(\Ov)$-invariant. Let $K=\prod_{v\in\Sigma'}B\GG(\Ov)$. For a function $g:\prod_{\vMF}BG(F_v)\to\CC$, let us denote by~$g_{\Sigma'}$ the restriction $g_{\Sigma'}=g|_{\prod_{v\in\Sigma'}BG(F_v)}$. %The function $(f^{-1}\cdot H)^{-s}$ is $K$-invariant, hence, formally one has that $$Z^f(s)=\sum_{\chi\in \Xi_K^{\perp}}\widehat{H(f)}(s,\chi),$$where $\widehat{H(f)}(s,\chi)$ is the Fourier transform of the function $(f^{-1}\cdot H)^{-s}$ and $\Xi_K^{\perp}$ is the subset of $BG^*(F)$ given by $\chi$ for which $\chi|_K\neq 1$. Note that $$\widehat{H(f)}(s,\chi)= $$
%Let $\Sigma'$ be a finite set of places of $F$ containing $\Sigma\cup M_F^{\infty}$ such that for every $v\in M_F-\Sigma'$, one has that $f_v\equiv 1$ and that $H_v:x\mapsto q^{\psi_v(c(x))}$. We set $K=\prod_{v\in M_F-\Sigma'}B\GG(\Ov)$ and $\Xi_K=BG^*(F)\cap K^{\perp}$.
In Lemma~\ref{condofpoiss}, we have verified the conditions of Poisson formula, and applying it gives:
\begin{align*}
Z^f(s)\hskip-2cm&\\&=\frac{1}{\tau_{BG}}\sum_{\chi\in \Xi^{\perp}_K}\int_{BG(\AAA)}f^sH^{-s}\chi\mu\\
&=\frac{1}{\tau_{BG}}\sum_{\chi\in \Xi^{\perp}_K}\bigg(\int_{\prod_{v\in \Sigma'}BG(F_v)}f_{\Sigma'}^sH_{\Sigma'}^{-s}\chi_{\Sigma'}\mu_{\Sigma'}\bigg) \cdot \prod_{v\in M_F-\Sigma'}\wH_v(s,\chi_v),
\end{align*}
whenever every quantity converges. It is evident that $$s\mapsto \int_{\prod_{v\in \Sigma'}BG(F_v)}f_{\Sigma'}^sH_{\Sigma'}^{-s}\chi_{\Sigma'}\mu_{\Sigma'}$$ is an entire function. By Corollary \ref{trivialcharacter}, one has that $\prod_{v\in M_F-\Sigma'}\wH_v(s,\chi_v)$ converges absolutely to a holomorphic function in the domain $\Re(s)>1$ and extends to a meromorphic function in the domain $\Re(s)\geq 1$ with a possible pole at $s=1$ of order no more than~$b(c)$.  We deduce that the same is valid for $s\mapsto Z^f(s)$. As $f\neq 0$, it follows from (1) that the pole exists and that is of order $b(c)$. Moreover, if $F$ is a number field, the function $s\mapsto\prod_{v\in M_F-\Sigma'}\wH_v(s,\chi_v)$ is holomorphic in the domain $\{\Re(s)\geq 1\}-\{1\}$. If~$F$ is a function field, the function $s\mapsto\prod_{v\in M_F-\Sigma'}\wH_v(s,\chi_v)$ is holomorphic in the domain $\{\Re(s)\geq 1-\epsilon\}-\{1+2ki\pi\log(q)^{-1}|k\in\ZZ\}$ for some~$\epsilon>0$. The two claims on holomorphicity of~$Z^f$ follow. The statement is proven.
\end{enumerate}
\end{proof}
%If $U\subset\prod_{\vMF}BG(F_v)$ is an open subset, for $s\in\CC$, we define series $$Z^U(s):=\frac{1}{|G(F)|}\sum_{\substack{x\in BG(F)\\i(x)\in U}}H(x)^{-s}.$$
%It follows from Proposition \ref{estimatenumber}, that there exists $\eta>0$ such that the series converges absolutely and uniformly for $\Re(s)>\eta$. The function $s\mapsto Z^U(s)$ is holomorphic in this domain.
\begin{rem}
\normalfont
By combining (1) and (2) of Proposition \ref{genfz}, we obtain a formula which can be used to compute $\omega_H(f)$ when $f=\sum \lambda_if_i,$ with~$\lambda_i\in\RR$ and~$f_i$ are elementary. By linearity, it is possible to write down the formula for~$\lambda_i\in\CC$. 
\end{rem}
\begin{rem}
\normalfont We remark that, contrary to the case of varieties,  the measure~$\omega_H$ is not always an Euler product of local measures, but only a finite sum of Euler products. This may happen even when $\overline{i(BG(F))}=\prod_{\vMF}BG(F_v)$ (e.g. when $G=\muu_m$). 
\end{rem}
\begin{thm}\label{satfunf}
Suppose that~$F$ is a function field. Let $$f:\prod_{\vMF}BG(F_v)\to \RR_{\geq 0}$$ be a continuous function such that~$\omega_H(f)> 0.$ %and which writes $f=\sum_{i=1}^n\lambda_if_i$ with $\lambda_i\in\RR$ and~$f_i$ elementary (i.e. as in Part (2) of Proposition \ref{genfz}). %Consider an elementary open~$U\subset\prod_{\vMF}BG(F_v)$ such that $\omega_H(U)>0$. 
One has that $$\#\{x\in BG(F)|H(x)\leq f(x)B\}\asymp_{B\to\infty} B\log(B)^{b(c)-1}.$$
\end{thm}
\begin{proof}
We split the proof in several parts. %By the fact that~$f$ is bounded it suffices to prove that the number of points is at least $B\log(B)^{b(c)-1}$.  
We start by making several simplifications on~$H$. 
\begin{enumerate}
\item We can suppose that~$c$ takes rational values. Indeed, for $r\in c(G_*(\oF))-\{0,1\}$, let us choose $r', r''\in\QQ_{>0}$ such that $1<r'<r<r''$ and such that for every $r, t\in c(G_*(\oF))-\{0,1\}$ one has that if $r<t$ then $r'<t'$ and $r''<t''$. We define a counting function $c':G_*(\oF)\to\QQ_{\geq 0}$ by $c'(x):= c(x)$ if $c(x)\in\{0,1\}$, otherwise $c'(x)=r'$, where $r=c(x)$. Analogously, one defines~$c''$. One can define a height~$H'$ by~$H'_v=H_v$ for $v\in\Sigma_{H}$ and $H'_v=q_v^{c'(\psi_v(x))}$. Analogously, one defines~$H''$. We have that that $H'(x)\leq H(x)\leq H''(x)$ for every $x\in BG(F)$. Thus if the statement is valid for the heights~$H'$ and~$H''$, it is also valid for the height~$H$. 
\item We prove that we can suppose that $f=\mathbf 1_U$ is the characteristic function of an open set~$U$ and that it is elementary (we call such open subsets elementary). The elementary open sets form a basis for the topology of $\prod_{\vMF}BG(F_v)$, hence the open set $\{x\in\prod_{\vMF}BG(F)|f(x)>0\}$ contains an elementary open~$V$. By finiteness of the spaces $BG(F_v)$, the elementary open sets are also closed. Hence, the set~$V$ is compact and it follows that there exists $\epsilon>0$ such that for every $x\in\prod_{\vMF}BG(F_v)$ one has that $f(x)\geq \epsilon\cdot\mathbf 1_V$. On the other side the function~$f$ is bounded and let~$C>0$ be such that for every $x\in\prod_{\vMF}BG(F_v)$, one has that $f(x)\leq C\cdot\mathbf 1_{\prod_{\vMF} BG(F_v)}$. %Note that if the claim is valid for 
%\item We prove that we can suppose that~$n=1$ and that~$f=f_1=\mathbf 1_U$ is the characteristic function of an open subset~$U$ of $\prod_{\vMF}BG(F_v)$, which is a finite union of elementary open sets (i.e. their characteristic functions are elementary). Let us set \begin{align*}U:&=\big\{x\in\prod_{\vMF}BG(F_v)|f(x)=0\big\}\\&=\big\{x\in \prod_{\vMF}BG(F_v)|\sum_{i=1}^nf_i(x)=0\big\}.\end{align*} There exists a finite set of places~$\Sigma'\subset M_F$ such that every $v\in M_F-\Sigma'$ one has that $f_{iv}=1$ for all~$i$, where $f_{iv}:BG(F_v)\to\RR_{\geq 0}$ are the functions defining~$f_i$.  We have that 
%\begin{align*}
%U=\big\{x\in \prod_{v\in\Sigma'}BG(F_v)\big|f_{\Sigma'}(x)=0\big\}\times\prod_{v\in M_F-\Sigma'}BG(F_v),
%\end{align*}
%where subscript has the same meaning as in Proposition~\ref{genfz}. 
%In particular, the closed set~$U$, by finiteness of every~$BG(F_v)$, is also open. The set $\prod_{\vMF}BG(F_v)-U$ is thus compact. Now, the continuous function~$f$ is strictly positive on the set $\prod_{\vMF}BG(F_v)-U$ and there are two constants $C_1, C_2>0$ such that $$f\big(\prod_{\vMF}BG(F_v)-U\big)\subset ]C_1, C_2[.$$ We have that 
\begin{align*}
\{x\in BG(F)|H(x)\leq \epsilon\mathbf 1_{V}B\}&\subset\{x\in BG(F)| H(x)\leq f(x)B\}\\&\subset \{x\in BG(F)|  H(x)\leq  CB\}.
\end{align*}
It follows thus to prove the statement for the characteristic function~$f=\mathbf 1_U$ of elementary open sets~$U$.
\item %Let $U=\cup_{j=1}^k\prod_{v\in M_F}U^j_{v}$ be a union of elementary open sets. 
Changing~$H$ at finitely many places will not change the validity of the statement, hence we can suppose that there exists a finite set of places $\Sigma_G\subset\Sigma$ such that for every $v\in\Sigma$ one has that~$H_v=1$. % and that $U^j_v=BG(F_v)$ for every~$1\leq j\leq k$. 
This simplification provides that % by Part (3) of Lemma \ref{importantchar}, we can suppose that there exists a finite set of places~$\Sigma_G\subset\Sigma$ such that the following conditions are satisfied: for every~$v\in\Sigma$, one has that~$H_v=1$, for every~$v\in M_F-\Sigma$ one has that $H_v=q_v^{c_v(x)}$ and $\Xi_H=BG(\AAA)$. The formula from  Proposition \ref{poisson} simplifies to $Z(s)=\frac{1}{\tau_{BG}}\wH(s,1).$ We also note that the set~
$$H(BG(F))\subset \{q^{\frac mr}|m\in\ZZ_{\geq 0}\}$$ for some~$r\in\ZZ_{>0}$.  
%\item We will use this part in Part (5). Let~$\rho$ be a Galois representation. We recall that in the domain $\Re(s)>0$, one has that $s\mapsto L(s,\rho)$ has poles in the set $\{1+2ki\pi\log(q)^{-1}|k\in\ZZ\}$. 
 %We verify that in the domain $\Re(s)>0$ one has that $s\mapsto L(s,\rho)$ has poles in the set $\{1+2ki\pi\log(q)^{-1}|k\in\ZZ\}$. Indeed, for every orbit~$W$ of~$\Gamma_F$ for its action on $G_*(\oF)_1$, we decompose the representation~$\rho^W_{1}$ into a direct sum of irreducible representations. For every~$W$, exactly one irreducible representation is the trivial one, namely the one corresponding to the subspace $\{(t)_{\phi}|\phi\in W\}$ of $\CC\cdot W$. For every non-trivial irreducible subrepresentation~$\tau$ of~$\rho^W_1$, we have by \cite[Pages 196-197]{Milne} that the corresponding~$L$-function $L(-,\tau)$ is holomorphic on~$\CC$. %Denote by $\mathcal Z(\mathcal C, T)$ the zeta function of~$\mathcal C$. 
%Now,  by using \cite[Page 196]{Milne} and \cite[Chapter VII, \S 6, Theorem 4]{basicnt}, we obtain that
%\begin{align*}
%L(s,\rho^{G_*(\oF)_1} )\hskip-2cm&\\&=\prod_{W\in (G_*(\oF)_1)/\Gamma_F}L(s,\rho^W_1)\\&=\prod_{W\in (G_*(\oF)_1)/\Gamma_F}\zeta_F(s)\prod_{1\neq \tau\subset \rho^W_1\text{ is irreducible}}L(s,\tau)\\
%&=\prod_{W\in (G_*(\oF)_1)/\Gamma_F}\frac{P(q^{-s})}{(1-q^{-s})(1-q^{1-s})}\prod_{1\neq \tau\subset \rho^W_1\text{ is irreducible}}L(s,\tau),
%\end{align*}
%where~$\zeta_F$ is the Dedekind zeta function of~$F$ and~$P$ is a polynomial such that $P(0)=1$. It follows that for some~$\epsilon>0$ the poles with $\Re(s)\geq 1-\epsilon$ are given by $s=1+2ki\pi\log(q)^{-1}$ for $k\in\ZZ.$ %and are all of order $b(c)$ by Lemma \ref{lfunf}. 
%We write
%\begin{multline*}Z^f(s)\\=\frac{\# \Sh^2(G)}{\# G^*(F)}\sum_{\chi\in \Xi_K^{\perp}}\bigg(\int_{\prod_{v\in \Sigma}BG(F_v)}f_{\Sigma}^sH_{\Sigma}^{-s}\chi_{\Sigma}\mu_{\Sigma}\bigg) \cdot \prod_{v\in M_F-\Sigma}\wH_v(s,\chi_v),\end{multline*}
% with the notation as in Part (2) of Proposition \ref{genfz}. The functions $$s\mapsto\int_{\prod_{v\in \Sigma}BG(F_v)}f_{\Sigma}^sH_{\Sigma}^{-s}\chi_{\Sigma}\mu_{\Sigma} $$ are entire. By combining Lemma \ref{} with the Part (3) of Lemma \ref{}, we obtain that the functions $$ $$ have the set of poles in the set $\{1+\}$

%Indeed, this follows from the fact that \begin{equation*}
%Z(s)=\frac{\wH(s,1)}{\tau_{BG}}=\frac{L(s,\rho_1^{G_*(\oF)})}{\tau_{BG}}\prod_{v\in M_F^0}L_v(s,\rho_1^{G_*(\oF)_1})^{-1}\wH_v(s,1),
%\end{equation*}
%Part (3) and Lemma \ref{ftlfun}.
%\begin{align*}
%Z(s)&=\frac{\wH(s,1)}{\tau_{BG}}\\&=\frac{\prod_{v\in \Sigma}\wH(s,1)}{\tau_{BG}}\prod_{v\in M_F-\Sigma}L_v(s,\rho_1)\cdot \prod_{v\in M_F-\Sigma}L_v(s,\rho_1)^{-1}\wH_v(s,1)\\
%&=
%\end{align*}
%\item We set $Q=\exp(r\cdot(\log_q\circ H)),$ so that for $B>0$, one has $H(x)\leq B$ if and only if $\log(Q(x))\leq r\log_q(B).$ We deduce the equivalences
%\begin{align*}
%\#\{x\in BG(F)|H(x)\leq B\}\asymp B\log(B)^{b(c)}\hskip-6cm&\\&\iff\#\{x\in BG(F)|\log(Q(x))\leq r\log_q(B)\}\asymp B\log(B)^{b(c)}\\
%&\iff \#\{x\in BG(F)|\log(Q(x))\leq rB\}\asymp q^B\log(q^B)^{b(c)}\\
%\#\big\{x\in BG(F)|\log(Q(x))\leq\frac{\log(B)}{\log(q)}\big\}\asymp B\log(B)^{b(c)}\\
%&\iff \#\{x\in BG(F)|Q(x)\leq \exp(rB)\}\asymp q^BB^{b(c)}\\
%&\iff \#\{x\in BG(F)|Q(x)\leq \exp(B)\}\asymp q^{B/r}B^{b(c)}\\
%&\iff \#\{x\in BG(F)|Q(x)\leq B\}\asymp q^{\log(B)/r}\log(B)^{b(c)}\\
%&\iff \#\{x\in BG(F)|Q(x)\leq B\}\asymp B^{\log(q)/r}\log(B)^{b(c)}.
%\end{align*}
%\item For $n\in\ZZ_{\geq 0}$, we denote by~$a_n$ the number of $x\in BG(F)$ having the height equal to~$q^{n/r}$, counted with multiplicity $\# G(F)^{-1}$. We are going to prove that $\sum_{n=1}^Ba_n\sim_{B\to\infty}B\log(B)^{b(c)}.$ 
%We have that $$\frac{\wH(s,1)}{\tau_{BG}}=Z(s)=\sum_{n\geq 0}a_n q^{-(n/r)s}=\sum_{n\geq 0}a_n\exp(-sn\log(q)/r),$$or equivalently, $$\frac{\wH(rs\log(q),1)}{\tau_{BG}}=Z(rs\log(q))=\sum_{n\geq 0}a_n\exp(-sn)$$in the domain $\Re(s)>\frac{1}{r\log(q)}-\varepsilon$ for some~$\varepsilon>0$.
%\begin{align*}
%\wH(s,1)&=A\prod_{v\in M_F-\Sigma}\bigg(\sum_{}1+q_v^{-s}\bigg)\\
%&=A\prod_{v\in M_F-\Sigma}\bigg(\exp(-s[\kappa_v:])\bigg)
%\end{align*} one has that
%\item We claim that for some~$\varepsilon>0$, in the domain $\Re(s)>1-\varepsilon$ one can write $\wH(s,1)=f(q^{-s})\beta(q^{-s})$, where~$f$ is rational with a pole at~$q^{-1}$ of order~$b(c)$ and $\beta$ holomorphic and bounded in the strip $0<\Re(s)<1$.NO OTHER POLES IN THE VERTICAL STRIP. For~$\vMFz$, we set $$L_v'(t)=$$ so that $L_v'(q^{-s})=L_v(s,\rho_1)$ are the factors.  By Lemma \cite[Example 13.6 a)]{Milne}, the product $f(t)=\prod_{\vMFz}L_v'(t)$ is a rational function of~$t$. Moreover, we have seen that~$s\mapsto L(s,\rho_1)=f(q^{-s})$ has a pole at~$s=1$ of order~$b(c)$, hence~$f$ has a pole at~$q^{-1}$ of order~$b(c)$. Let us set $\beta_v()$ so that $$\beta_v(q^{-s})=L_v()^{-1}.$$ By \ref{}The product defining~$\beta(q^{-s})$ converges absolutely for every $s\geq 1-\varepsilon$ for some~$\varepsilon>0$ and defines a holomorphic function in the domain $\Re(s)>1-\varepsilon$. It follows that $\beta(t):=\prod_v \beta_v(t)$ converges absolutely for $t<1+\varepsilon'$ for some $\varepsilon'>0$ and defines a holomorphic and bounded function in the domain $0<\Re(s)<1$. The claim is proven.
%\item One has that 
%\begin{align*}\sum_{n\geq 0}a_n\exp(-ns)&=\frac{\wH(rs\log(q),1)}{\tau_{BG}}\\&=f(q^{-rs\log(q)})\beta(q^{-rs\log(q)})\\
%&=f(\exp(-s\log(q)^2r))\beta(\exp(-s\log(q)^2r))
%\end{align*}
\item Let us complete the proof. For $n\in\ZZ_{\geq 0}$, we denote by~$a_n$ the number of $x\in BG(F)$ with $i(x)\in U$ having the height equal to~$q^{n/r}$, counted with multiplicity $\# G(F)^{-1}$. Let us introduce the variable $t=~q^{-s/r}$. We define power series $Q(t)=\sum _{n=0}^{\infty}a_nt^{n}$ so that $$Q(t)=Q(q^{-s/r})=Z^f(s).$$
By Proposition \ref{genfz}, the series defining $Z^f(s)$ converges absolutely for $\Re(s)> 1$ to a holomorphic function, the function $s\mapsto Z^f(s)$ extends to a meromorphic function in the domain $\Re(s)>1-\epsilon$ for some~$\epsilon>0.$ Moreover, as $\omega_H(\mathbf 1_U)=\omega_H(f)>0$, the function~$Z^f$ has a pole at~$s=1$ of order~$b(c)$ and no poles outside of the set $\{1+2ik\pi\log(q)^{-1}|k\in\ZZ\}$. %Moreover, the function~$s\mapsto Z^f(s)$ is $2i\pi \log(q)^{-1}r$-periodic by the assumptions on~$H$.
Hence, the series~$Q(t)$ converges absolutely for $|t|<q^{-1/r}$ to a holomorphic function and it extends to a meromorphic function of the disk $|t|<q^{-1/r}+\epsilon$ for some~$\epsilon$ with the only pole at~$t=q^{-1/r}$ of order $b(c)$. By Part~(c) of \cite[Theorem A.5]{lagemanartinschreier}, for some $C>0$, one has that $$\sum_{n=0}^Ma_n\sim_{M\to\infty} C\cdot q^{M/r}\log(q^{M/r})^{b(c)-1}.$$
Finally, $B\asymp_{B\to\infty} q^{\lfloor\log_{q}(rB)\rfloor/r}$, and hence 
\begin{align*}
\#\{x\in BG(F)|i(x)\in U, H(x)\leq B\}\hskip-2cm&\\&=\sum_{n=0}^Ba_n\\&=\sum_{n=0}^{\lfloor\log_{q}(rB)\rfloor}a_n\\
&\sim_{B\to\infty}C\cdot q^{\lfloor \log_{q}(rB)\rfloor/r}\log(q^{\lfloor\log_{q}(rB) \rfloor/r})^{b(c)-1}\\
&\asymp_{B\to\infty} B\log(B)^{b(c)-1}.
\end{align*}
The statement follows.
%\item  For simplicity, we shorten $\rho:=\rho_1^{G_*(\oF)_1}$. We will use \cite[Example 13.6]{Milne} to establish that $L(s,\rho)$ is rational in~$q^{-s}$. (The representation~$\rho$ is a complex representation, taking values in permutation matrices, hence is a~$\QQ_{\ell}$-representation, where~$\ell$ is a prime.) There exists some open immersion $j:U\hookrightarrow\mathcal C$, such that~$\rho$ factorizes through $\pi_1(U,\overline{\eta_{\mathcal C}})$, where~$\overline{\eta_{\mathcal C}}$ is a geometric generic point of~$\mathcal C$. The induced $\QQ_{\ell}$-representation of $\pi_1(U, \overline{\eta}_{\mathcal C}),$ also denoted by~$\rho,$ defines a locally constant constructible sheaf~$E_{\rho}'$ which satisfies $E_{\rho}(U)=(\QQ_{\ell}\cdot G_*(\oF)_1)^{\rho}.$ We set $E_{\rho}:=j_*(E_{\rho}').$ By using \cite[Example 13.6]{Milne}, we obtain that $$L(s,\rho)=\mathcal Z(\mathcal C, E_{\rho})(q^{-s}),$$ where $\mathcal Z(\mathcal C, E_{\rho})$ is the zeta function associated to the sheaf~$E_{\rho}.$ Now, \cite[Theorem 13.3]{Milne}, gives that \begin{multline*}L(s,\rho)=Z(\mathcal C, E_{\rho})(q^{-s})\\=\frac{\det(1-q^{-s}\Frob |H^1_{c}(\overline {\mathcal C}, \overline {E_{\rho}}))}{\det(1-q^{-s}\Frob|H^0_c(\overline C, \overline {E_{\rho}}))\cdot \det(1-q^{-s}\Frob |H^2_c(\overline C, \overline {E_{\rho}}))},
%\end{multline*}
%where~$\Frob$ is the geometric Frobenius, $\overline{\mathcal C}$ and $\overline{E_{\rho}}$ are base changes to~$\overline {\mathbb F_q}$ and cohomologies are with compact support. POLES ARE IN THE DENOMINATOR. We have that $H^0_c(\overline{\mathcal C}, \overline{E_{\rho}})=\overline {E_{\rho}}(\overline{\mathcal C})=(\QQ_{\ell}\cdot G_*(\oF))^{\rho}=G_*(F).$  On the other side
\end{enumerate}
\end{proof}
When~$F$ is a number field, we prove that $G$-torsors are equidistributed in the product space $\prod_{\vMF}BG(F_v)$.
\begin{thm}%[``Equidistribution of $G$-torsors''] 
\label{thmofequidistribution} Suppose $F$ is a number field. We define a Borel measure~$\nu$ on $\prod_{\vMF}BG(F_v)$ by $$\nu:=\frac{i_*\coun_{BG(F)}}{\#G(F)}.$$
 For $B>0$, we define a Borel measure~$\nu_B$ on $\prod_{\vMF}BG(F_v)$ by $$\nu_B:=\frac{\nu}{\frac{\#\Sh^2(G)}{\# G^*(F)\cdot (b(c)-1)!}\cdot B\cdot \log(B)^{b(c)-1}}\cdot\mathbf 1_{\{(x_v)_v|H((x_v)_v)\leq B\}}.$$ The measures~$\nu_B$ converge vaguely to the measure~$\omega_H$ when~$B\to\infty$. In other words,
\begin{enumerate}
\item for any continuous $f:\prod_{\vMF}BG(F_v)\to\CC$ one has that $$\lim_{B\to\infty}\frac{\frac{1}{\# G(F)}\sum_{\substack{x\in BG(F)\\H(x)\leq B}}f(x)}{B\cdot\log(B)^{b(c)-1}}=\frac{\#\Sh^2(G)}{\# G^*(F)\cdot (b(c)-1)!}\cdot \omega_H(f),$$
\item for every open set $U\subset\prod_{\vMF}BG(F_v)$ the boundary of which is $\omega_H$-negligible, one has that \begin{multline*}\frac{\#\{x|i(x)\in U, H(x)\leq B\}}{\# G(F)}\\=\frac{\#\Sh^2(G)\cdot\omega_H(U)}{\# G^*(F)\cdot (b(c)-1)!}\cdot {B\log(B)^{b(c)-1}}+o(B\log(B)^{b(c)-1}).
\end{multline*}
\end{enumerate} 
 %
%
%$$ $$For $B>0$, we set $\nu_B$ for the measure $$ $$
%
% Let $U\subset\prod_{\vMF}BG(F_v)$ be an open subset such that for its boundary $\partial U,$ one has $\omega_H(\partial U)=0$. One has that 
%$$\lim_{B\to +\infty}\frac{N_H(U, B)}{N_H(B)} =\frac{\omega_H\big(\overline{i(BG(F))}\cap U\big)}{\omega_H\big(\overline{i(BG(F))}\big)}.$$
\end{thm}
\begin{proof}
Let~$f:\prod_{\vMF}BG(F_v)\to\RR_{\geq 0}$ be a function which is a finite linear combination of elementary functions. By Proposition~\ref{genfz} and the Tauberian result from \cite[Theorem 3.3.2]{aniso}, one has that \begin{align*}\nu(\{x\in i(BG(F))| H(x)\leq f(x)\cdot B\})\sim_{B\to\infty}\frac{\#\Sh^2(G)\cdot\omega_H(f)}{\#G^*(F)\cdot(b(c)-1)!}{B\log(B)^{b(c)-1}}.
\end{align*} 
By Part (c) and Part (d) of proof of Lemma \ref{existmeas}, we have that the functions, which are finite linear combinations of elementary functions~$\otimes_vf_v$, are dense in the space of continuous real valued functions $\prod_{\vMF}BG(F_v)\to~\RR$. The statement is now a special case of \cite[Proposition 2.10]{Igusa}. 
%Let $\epsilon>0$. There exist continuous functions $f, g:\prod_{\vMF}BG(F_v)\to\RR$ such that $$0\leq f\leq \mathbf 1_U\leq g$$ and such that $\omega_H(g-f)<\varepsilon.$ Let us set $f_\varepsilon=\varepsilon +(1-\varespsilon)f,$ $h_\varepsilon=\varepsilon +(1-\varepsilon) \mathbf 1_U$ and $g_\varepsilon=\varepsilon+(1-\varepsilon)g$. By Lemma \ref{existmeas}, one has that $\omega_H(f_{\varepsilon}),$ $\omega_H(h_\varepsilon)$ and $\omega_H(g_\varepsilon)$ exist. Moreover, one has that
%Let $\theta:\prod_{\vMF}BG(F_v)\to\RR_{>0}$ be a continuous function such that $\theta$ is a finite linear combination with positive coefficients $\lambda_i>0$ of elementary functions $\theta=\sum_{i=1}^n\lambda_if_i$. Let us establish that \begin{align*}Z(\theta, s):&=\frac{1}{\# G(F)}\sum_{x\in BG(F)}\theta(i(x))^sH(i(x))^{-s}\\
%&=\frac{1}{\# G(F)}\sum_{x\in BG(F)}\big(\sum_{i=1}^n\lambda _if_i(x)\big)^s\cdot H(x)^{-s}.
%\end{align*}
%converges whenever $Z$ converges (hence, at least in the domain $\Re(s)>\eta$ for certain $\eta>0$) and extends to a meromorphic function in the domain $\Re(s)\geq 1$ with a simple pole at $s=1$ of order $b(c)$ and no other poles in the domain $\{\Re(s)\geq 1\}-\{1\}$. 
\end{proof}
%\subsection{Fairness condition}
%We propose the following fairness condition.
%\begin{mydef}
%We say that a counting function~$c: J_*\to\QQ_{\geq 0}$ is fair, if the set $$\big\{g\in \Hom(\muu_e, J)|c([g])=\min_{1\neq h\in J_*}c(h)\text{ and } \Gamma_F g=g\big\},$$ where $[g]$ is the class of~$g$ in $J_*=\Hom(\muu_e, J)/J$, generates the group~$J_*$.
%\end{mydef}
%We say that a height $BG(F)\to\RR_{>0}$ is fair if the corresponding counting function is fair.We present several examples.
%\begin{lem}
%\begin{enumerate}
%\item The function $c_{\Delta, \mathfrak S_n}:(\mathfrak S_n)_*\to \QQ_{\geq 0}$ is a fair function. 
%\item
%\end{enumerate}
%\end{lem}
%\begin{proof}
%\begin{enumerate}
%\item
%\end{enumerate}
%\end{proof}
%\begin{exam}
%\normalfont
%%Let $n$ be a positive integer. By Lemma \ref{slabalema}, a counting function $c:(\muu_n)_*\to \QQ_{\geq 0}$ is fair if and only if $\{g\in (\muu_n)_*| c(g)=\min_{1\neq h\in (\muu_n)_*}c(h)\}$ generates~$(\muu_n)_*$. It follows, in particular, that a quasi-toric height from \cite{darda2021rational} is fair.
% I DON'T KNOW IF A COMMUTATIVE CASE DIFFERENT THAN PRODUCT OF $\mu_n$ EXISTS. MAYBE REMOVE COMPLETELY????
%, because for a (non-canonical) isomorphism $(\muu_n)\cong \ZZ/n\ZZ$ one has that $c(x)=1 $
%\end{exam}
%\begin{conj}\label{fairconj}
%Let $c:G_*(F)\to \QQ_{\geq 0}$ be a normalized fair counting function. We choose a finite subset $S\subset M_F$ containing $\Sigma\cup M_F^\infty$ and for $v\in S$, functions $c_v:BG(F_v)\to \RR_{>0}$. We set $H=H((c_v)_v)$. Let us set \begin{multline*}\tau_H\\:=\Res(\zeta_F,1)^{b(c)}\bigg(\prod_{\vMFz}\zeta_v(1)^{-b(c)}\int_{BG(F_v)}H_v^{-1}\mu_v\bigg)\cdot \prod_{\vMFi}\int_{BG(F_v)}H_v^{-1}\mu_v,
%\end{multline*}
%where $\Res(\zeta_F,1)$ is the residue of the Dedekind zeta function at $s=1$. One has that $$\#\{x\in BG(F)\}\sim_{B\to\infty}=\frac{\tau_H}{(\# G^*(F))\cdot (b(c)-1)!} B\log(B)^{b(c)-1},$$where~$G^*(F)$ stands for the set of $1$-dimensional representations of~$G(\oF)$ defined over~$F$.
%\end{conj}
%We present results of 
%\begin{thm}
%Conjecture \ref{fairconj} is valid in the following cases:
%\begin{enumerate}
%\item For $G$ commutative and a general fair height.
%\item Let $n$ be an integer and let $G=\mathfrak S_n$ and suppose that for the height is taken to be $H=\Delta^\iota$. The prediction of Conjecture \ref{fairconj} is consistent with the prediction of \cite[Conjecture A]{massformulae}. Hence, it is valid for $G= S_3, $ and $S_5$ and $F=\QQ$ by the main results of \cite{davenportheilbronn}, \cite{densityquartic} and  \cite{densityquintic}, respectively.
%\end{enumerate}
%\end{thm}
%\begin{proof}
%\begin{enumerate}
%\item
%%\item Let$~v\in M_F$ be a place. Note that in this situation, one has that $b(c)=1.$ The map \begin{align*}
%B\mathfrak S_n(F_v)&\to \{\textnormal{degree $n$ \' etale algebras over $F_v$ up to isomorphism}\}\\x&\mapsto (x/S_{n-1}),
%\end{align*} where $S_{n-1}$ acts on $x$ via the inclusion $\iota:S_{n-1}\hookrightarrow \mathfrak S_n$, is a automorphism preserving bijection (REFERENCE). Moreover, by definition, one has that $\Delta^{\iota}(x)=\Delta(x/S_{n-1})$. It follows that for~$v$ finite, one has $$\int_{BG(F_v)}(\Delta^{\iota}_v)^{-1}\mu_v=\sum_{\substack{ E\textnormal{ \' etale algebras of degree $n$}\\\textnormal{up to isomorphism}}}\frac{\Delta(E)^{-1}}{\#\Aut(E)}.$$
%By using \cite[Theorem 1.1]{massformulae} to compute the sum on the right hand side, we obtain that for a finite place~$v$, one has that
%\begin{align*}
%\zeta_v(1)^{-1}\int_{BG(F_v)}(\Delta^\iota_v)^{-1}\mu_v\hskip-1cm&\\&=(1-N(\mathfrak p_v)^{-1})\sum_{k=0}^{n-1}q(k,n-k)N(\mathfrak p_v)^{-1}\\
%&=\sum_{k=0}^{n-1}(q(k,n-k)-q(k-1,n-k+1))N(\mathfrak p_v)^{-1},
%\end{align*}
%where $q(k,r)$ denotes the number of partitions of~$k$ into at most~$r$ parts and~$\mathfrak{p}_v$ is the prime ideal corresponding to~$v$. For a real place~$v$, one has by \cite[Proposition 2.4]{massformulae} that $$\int_{BG(F_v)}(\Delta_v^{-1})\mu_v=\frac{r_2(\mathfrak S_n)}{n!},$$where $r_2(\mathfrak S_n)$ stands for the number of elements of order at most~$2$ in~$\mathfrak S_n$. For~$v$ complex, one has by \cite[Page 18]{massformulae} that $$\int_{BG(F_v)}(\Delta^\iota_v)^{-1}\mu_v=\frac{1}{n!}.$$ 
%Note that $\# \mathfrak S_n^*(F)=2$, the $1$-dimensional representations being the sign functions. The statement now follows directly by the comparing our leading constant with the leading constant of \cite[Conjecture~A]{massformulae}.
%\end{enumerate}
%\end{proof}
%\begin{rem}
%\normalfont
%A non-trivial contribution of the factor $\#\Sh^2(F, Z(G))$ can be seen for some non commutative~$G$. Let us set $F=\QQ(\sqrt{7})$ and $G=(\ZZ/8\ZZ)\times S_3$. so that $Z(G)=(\ZZ/8\ZZ)$. By \cite[Section 2, a)]{Sansuc}, one has that $$\#\Sh^2(F,Z(G))=\#\Sh^2(\QQ(\sqrt{7},\ZZ/8\ZZ)=\# \Sh^1(\QQ(\sqrt{7}),\muu_8)=\ZZ/2\ZZ.$$
%Let us set $c=(c', c_{\Delta})$, where $c'$ 
%By Lemma \ref{} and , we have that $$ $$
%\end{rem}
%\begin{rem}\normalfont
%The image of $BG(F)$ for the diagonal map to $\prod_{\vMF}BG(F_v)$ is not always dense. Indeed, by \cite[Lemma 1.4]{Sansuc}, the closure of the image is of index $$\frac{|\Sh^1_S(F,G)|}{|\Sh^1(G^*)|} ,$$ where $G^*$ is the Cartier dual of $G$ and $S$ is the finite set of places for which the decomposition group in $F_{G^*}/F$ is not cyclic.
%\end{rem}
%
We have noted in Paragraph~\ref{parbreakthin} that in the commutative case there are no breaking thin maps. Thus Theorems~\ref{satfunf} and~\ref{thmofequidistribution}, together with the fact that for every height~$H$ of type~$c$, the type of the height~$H^{\frac{1}{a(c)}}$ is normalized, imply that Conjecture~\ref{malletale} is verified for~$G$ commutative. Moreover, in the number field case, for non-normalized heights~$H$ of type~$c$ we get that 
\begin{align*}\frac{\#\{x\in BG(F)|H(x)\leq B\}}{\# G(F)}=\frac{\#\Sh^2(G)\cdot\omega_H\big(\overline{i(BG(F))}\big)}{\# G^*(F)\cdot (b(c)-1)!}B^{a(c)}\log(B)^{b(c)-1}.
\end{align*}
\begin{rem}
 \normalfont Peyre conjectures in \cite{Peyre}, that for a smooth Fano $F$-variety~$X$, the leading constant in the asymptotic formula for the number of points of bounded height is a product of three factors: a factor connected to the location of the anticanonical line bundle in the ample cone, the cardinality of the cohomology group $\# H^1(\Gamma_F, \Pic(\overline F))$ and the volume of the subset~$X(\overline F)\subset \prod_{\vMF}X(F_v)$ for a measure analogue to~$\omega_H$. In our situation, the volume factor appears, but we are unsure whether the other factors can be understood as in the case of varieties.
\end{rem}
\bibliography{bibliography}
\bibliographystyle{acm}
\end{document}

\begin{thebibliography}{10}

\bibitem{statfirstg}
{\sc Alberts, B.}
\newblock Statistics of the first {G}alois cohomology group: a refinement of
  {M}alle's conjecture.
\newblock {\em Algebra Number Theory 15}, 10 (2021), 2513--2569.

\bibitem{odorney}
{\sc Alberts, B., and O'Dorney, E.}
\newblock Harmonic analysis and statistics of the first {G}alois cohomology
  group.
\newblock {\em Res. Math. Sci. 8}, 3 (2021), Paper No. 50, 16.

\bibitem{aniso}
{\sc Batyrev, V.~V., and Tschinkel, Y.}
\newblock Rational points of bounded height on compactifications of anisotropic
  tori.
\newblock {\em Internat. Math. Res. Notices}, 12 (1995), 591--635.

\bibitem{densityquartic}
{\sc Bhargava, M.}
\newblock The density of discriminants of quartic rings and fields.
\newblock {\em Ann. of Math. (2) 162}, 2 (2005), 1031--1063.

\bibitem{densityquintic}
{\sc Bhargava, M.}
\newblock The density of discriminants of quintic rings and fields.
\newblock {\em Ann. of Math. (2) 172}, 3 (2010), 1559--1591.

\bibitem{https://doi.org/10.48550/arxiv.1512.03035}
{\sc Bhargava, M., Shankar, A., and Wang, X.}
\newblock Geometry-of-numbers methods over global fields i: Prehomogeneous
  vector spaces.
\newblock \url{https://arxiv.org/pdf/1512.03035.pdf}, 2015.

\bibitem{sexticsthree}
{\sc Bhargava, M., and Wood, M.~M.}
\newblock The density of discriminants of {$S_3$}-sextic number fields.
\newblock {\em Proc. Amer. Math. Soc. 136}, 5 (2008), 1581--1587.

\bibitem{TSpectrale}
{\sc Bourbaki, N.}
\newblock {\em \'{E}l\'{e}ments de math\'{e}matique. {F}asc. {XXXII}.
  {T}h\'{e}ories spectrales. {C}hapitre {I}: {A}lg\`ebres norm\'{e}es.
  {C}hapitre {II}: {G}roupes localement compacts commutatifs}.
\newblock Actualit\'{e}s Scientifiques et Industrielles [Current Scientific and
  Industrial Topics], No. 1332. Hermann, Paris, 1967.

\bibitem{TopologieGj}
{\sc Bourbaki, N.}
\newblock {\em \'{E}l\'{e}ments de math\'{e}matique. {T}opologie
  g\'{e}n\'{e}rale. {C}hapitres 1 \`a 4}.
\newblock Hermann, Paris, 1971.

\bibitem{TopologieGd}
{\sc Bourbaki, N.}
\newblock {\em \'{E}l\'{e}ments de math\'{e}matique. {T}opologie
  g\'{e}n\'{e}rale. {C}hapitres 5 \`a 10}.
\newblock Hermann, Paris, 1974.

\bibitem{Integrationj}
{\sc Bourbaki, N.}
\newblock {\em Integration. {I}. {C}hapters 1--6}.
\newblock Elements of Mathematics (Berlin). Springer-Verlag, Berlin, 2004.
\newblock Translated from the 1959, 1965 and 1967 French originals by Sterling
  K. Berberian.

\bibitem{Integrationd}
{\sc Bourbaki, N.}
\newblock {\em Integration. {II}. {C}hapters 7--9}.
\newblock Elements of Mathematics (Berlin). Springer-Verlag, Berlin, 2004.
\newblock Translated from the 1963 and 1969 French originals by Sterling K.
  Berberian.

\bibitem{bourquifzh}
{\sc Bourqui, D.}
\newblock Fonction z\^{e}ta des hauteurs des vari\'{e}t\'{e}s toriques non
  d\'{e}ploy\'{e}es.
\newblock {\em Mem. Amer. Math. Soc. 211}, 994 (2011), viii+151.

\bibitem{brauertheorem}
{\sc Brauer, R.}
\newblock On {A}rtin's {$L$}-series with general group characters.
\newblock {\em Ann. of Math. (2) 48\/} (1947), 502--514.

\bibitem{Igusa}
{\sc Chambert-Loir, A., and Tschinkel, Y.}
\newblock Igusa integrals and volume asymptotics in analytic and adelic
  geometry.
\newblock {\em Confluentes Math. 2}, 3 (2010), 351--429.

\bibitem{brauergrothendieck}
{\sc Colliot-Th\'{e}l\`ene, J.-L., and Skorobogatov, A.~N.}
\newblock {\em The {B}rauer-{G}rothendieck group}, vol.~71 of {\em Ergebnisse
  der Mathematik und ihrer Grenzgebiete. 3. Folge. A Series of Modern Surveys
  in Mathematics [Results in Mathematics and Related Areas. 3rd Series. A
  Series of Modern Surveys in Mathematics]}.
\newblock Springer, Cham, [2021] \copyright 2021.

\bibitem{darda:tel-03682761}
{\sc Darda, R.}
\newblock {\em {Rational points of bounded height on weighted projective
  stacks}}.
\newblock Theses, {Universit{\'e} Paris Cit{\'e}}, Sept. 2021.

\bibitem{dardayasudabm}
{\sc Darda, R., and Yasuda, T.}
\newblock The {B}atyrev-{M}anin conjecture for {D}{M} stacks.
\newblock \url{https://arxiv.org/pdf/2207.03645.pdf}, 2022.

\bibitem{dardayasudasemicomm}
{\sc Darda, R., and Yasuda, T.}
\newblock Inverse {G}alois problem for semicommutative finite group schemes,
  2022.

\bibitem{dave}
{\sc Davenport, H., and Heilbronn, H.}
\newblock On the density of discriminants of cubic fields. {II}.
\newblock {\em Proc. Roy. Soc. London Ser. A 322}, 1551 (1971), 405--420.

\bibitem{Ellenberg}
{\sc Ellenberg, J.~S., Satriano, M., and Zureick-Brown, D.}
\newblock {Heights on stacks and a generalized Batyrev-Manin-Malle conjecture}.
\newblock \url{https://arxiv.org/pdf/2106.11340.pdf}, 2021.

\bibitem{Gmethodant}
{\sc Ellenberg, J.~S., and Venkatesh, A.}
\newblock Counting extensions of function fields with bounded discriminant and
  specified {G}alois group.
\newblock In {\em Geometric methods in algebra and number theory}, vol.~235 of
  {\em Progr. Math.} Birkh\"{a}user Boston, Boston, MA, 2005, pp.~151--168.

\bibitem{frankemanintschinkel}
{\sc Franke, J., Manin, Y.~I., and Tschinkel, Y.}
\newblock Rational points of bounded height on {F}ano varieties.
\newblock {\em Invent. Math. 95}, 2 (1989), 421--435.

\bibitem{hassenormprin}
{\sc Frei, C., Loughran, D., and Newton, R.}
\newblock The {H}asse norm principle for {A}belian extensions.
\newblock {\em Amer. J. Math. 140}, 6 (2018), 1639--1685.

\bibitem{prescnorms}
{\sc Frei, C., Loughran, D., and Newton, R.}
\newblock Number fields with prescribed norms.
\newblock {\em Comment. Math. Helv. 97}, 1 (2022), 133--181.
\newblock With an appendix by Yonatan Harpaz and Olivier Wittenberg.

\bibitem{cohomologicalinvariants}
{\sc Garibaldi, S., Merkurjev, A., and Serre, J.-P.}
\newblock {\em Cohomological invariants in {G}alois cohomology}, vol.~28 of
  {\em University Lecture Series}.
\newblock American Mathematical Society, Providence, RI, 2003.

\bibitem{SGAj}
{\sc Grothendieck, A.}
\newblock {\em Rev\^{e}tements \'{e}tales et groupe fondamental ({SGA} 1)},
  vol.~3 of {\em Documents Math\'{e}matiques (Paris) [Mathematical Documents
  (Paris)]}.
\newblock Soci\'{e}t\'{e} Math\'{e}matique de France, Paris, 2003.
\newblock S\'{e}minaire de g\'{e}om\'{e}trie alg\'{e}brique du Bois Marie
  1960--61. [Algebraic Geometry Seminar of Bois Marie 1960-61], Directed by A.
  Grothendieck, With two papers by M. Raynaud, Updated and annotated reprint of
  the 1971 original [Lecture Notes in Math., 224, Springer, Berlin; MR0354651
  (50 \#7129)].

\bibitem{haraprox}
{\sc Harari, D.}
\newblock Quelques propri\'{e}t\'{e}s d'approximation reli\'{e}es \`a la
  cohomologie galoisienne d'un groupe alg\'{e}brique fini.
\newblock {\em Bull. Soc. Math. France 135}, 4 (2007), 549--564.

\bibitem{Hararicohom}
{\sc Harari, D.}
\newblock {\em Galois cohomology and class field theory}.
\newblock Universitext. Springer, Cham, [2020] \copyright 2020.
\newblock Translated from the 2017 French original by Andrei Yafaev.

\bibitem{Hortsch}
{\sc Hortsch, R.}
\newblock Counting elliptic curves of bounded {F}altings height.
\newblock {\em Acta Arith. 173}, 3 (2016), 239--253.

\bibitem{Kluners}
{\sc Kl\"{u}ners, J.}
\newblock A counterexample to {M}alle's conjecture on the asymptotics of
  discriminants.
\newblock {\em C. R. Math. Acad. Sci. Paris 340}, 6 (2005), 411--414.

\bibitem{lagemanartinschreier}
{\sc Lagemann, T.}
\newblock Distribution of {A}rtin-{S}chreier extensions.
\newblock {\em J. Number Theory 132}, 9 (2012), 1867--1887.

\bibitem{geometricconsitencymanin}
{\sc Lehmann, B., Sengupta, A.~K., and Tanimoto, S.}
\newblock Geometric consistency of {M}anin's conjecture.
\newblock \url{https://arxiv.org/pdf/1805.10580.pdf}, 2018.

\bibitem{Malle}
{\sc Malle, G.}
\newblock On the distribution of {G}alois groups. {II}.
\newblock {\em Experiment. Math. 13}, 2 (2004), 129--135.

\bibitem{https://doi.org/10.48550/arxiv.2004.04651}
{\sc Masri, R., Thorne, F., Tsai, W.-L., and Wang, J.}
\newblock Malle's conjecture for ${G} \times {A}$, with ${G} = {S}_3, {S}_4,
  {S}_5$.
\newblock \url{https://arxiv.org/abs/2004.04651v2}, 2020.

\bibitem{Milne}
{\sc Milne, J.~S.}
\newblock {\em \'{E}tale cohomology}.
\newblock Princeton Mathematical Series, No. 33. Princeton University Press,
  Princeton, N.J., 1980.

\bibitem{ArithmeticDuality}
{\sc Milne, J.~S.}
\newblock {\em Arithmetic duality theorems}, second~ed.
\newblock BookSurge, LLC, Charleston, SC, 2006.

\bibitem{CohomologyNF}
{\sc Neukirch, J., Schmidt, A., and Wingberg, K.}
\newblock {\em Cohomology of number fields}, second~ed., vol.~323 of {\em
  Grundlehren der mathematischen Wissenschaften [Fundamental Principles of
  Mathematical Sciences]}.
\newblock Springer-Verlag, Berlin, 2008.

\bibitem{Peyre}
{\sc Peyre, E.}
\newblock Hauteurs et mesures de {T}amagawa sur les vari\'{e}t\'{e}s de {F}ano.
\newblock {\em Duke Math. J. 79}, 1 (1995), 101--218.

\bibitem{peyrethin}
{\sc Peyre, E.}
\newblock Points de hauteur born\'{e}e, topologie ad\'{e}lique et mesures de
  {T}amagawa.
\newblock vol.~15. 2003, pp.~319--349.
\newblock Les XXII\`emes Journ\'{e}es Arithmetiques (Lille, 2001).

\bibitem{Ramakrishnan}
{\sc Ramakrishnan, D., and Valenza, R.~J.}
\newblock {\em Fourier analysis on number fields}, vol.~186 of {\em Graduate
  Texts in Mathematics}.
\newblock Springer-Verlag, New York, 1999.

\bibitem{SchmidtNF}
{\sc Schmidt, W.~M.}
\newblock Number fields of given degree and bounded discriminant.
\newblock No.~228. 1995, pp.~4, 189--195.
\newblock Columbia University Number Theory Seminar (New York, 1992).

\bibitem{corpslocaux}
{\sc Serre, J.-P.}
\newblock {\em Corps locaux}.
\newblock Publications de l'Universit\'{e} de Nancago, No. VIII. Hermann,
  Paris, 1968.
\newblock Deuxi\`eme \'{e}dition.

\bibitem{coursdarithmetique}
{\sc Serre, J.-P.}
\newblock {\em A course in arithmetic}.
\newblock Graduate Texts in Mathematics, No. 7. Springer-Verlag, New
  York-Heidelberg, 1973.
\newblock Translated from the French.

\bibitem{Cohomologiegalois}
{\sc Serre, J.-P.}
\newblock {\em Cohomologie galoisienne}, fifth~ed., vol.~5 of {\em Lecture
  Notes in Mathematics}.
\newblock Springer-Verlag, Berlin, 1994.

\bibitem{turkeli}
{\sc T\"{u}rkelli, S.}
\newblock Connected components of {H}urwitz schemes and {M}alle's conjecture.
\newblock {\em J. Number Theory 155\/} (2015), 163--201.

\bibitem{wangmalle}
{\sc Wang, J.}
\newblock Malle's conjecture for {$S_n \times A$} for {$n=3,4,5$}.
\newblock {\em Compos. Math. 157}, 1 (2021), 83--121.

\bibitem{basicnt}
{\sc Weil, A.}
\newblock {\em Basic number theory}.
\newblock Classics in Mathematics. Springer-Verlag, Berlin, 1995.
\newblock Reprint of the second (1973) edition.

\bibitem{widder}
{\sc Widder, D.~V.}
\newblock {\em The {L}aplace {T}ransform}.
\newblock Princeton Mathematical Series, vol. 6. Princeton University Press,
  Princeton, N. J., 1941.

\bibitem{onprobabilitieslocal}
{\sc Wood, M.~M.}
\newblock On the probabilities of local behaviors in abelian field extensions.
\newblock {\em Compos. Math. 146}, 1 (2010), 102--128.

\bibitem{Wright}
{\sc Wright, D.~J.}
\newblock Distribution of discriminants of abelian extensions.
\newblock {\em Proc. London Math. Soc. (3) 58}, 1 (1989), 17--50.

\bibitem{https://doi.org/10.48550/arxiv.1408.3912}
{\sc Yasuda, T.}
\newblock Densities of rational points and number fields.
\newblock \url{https://arxiv.org/pdf/1408.3912.pdf}, 2014.

\bibitem{Yasuda}
{\sc Yasuda, T.}
\newblock {Manin's conjecture vs. Malle's conjecture}.
\newblock \url{https://arxiv.org/pdf/1505.04555.pdf}, 2015.

\end{thebibliography}
